%% CNT/Output.tex — Connes–Narnhofer–Thirring Entropy: basics
%% Written incrementally; see CNT/progress.md for checkpoint log.
\documentclass[12pt,a4paper]{article}

\usepackage{amsmath,amsthm,amssymb,amsfonts}
\usepackage{mathtools}
\usepackage{bbm}
\usepackage{geometry}
\usepackage{hyperref}
\usepackage{booktabs}
\usepackage{array}

\geometry{margin=2.5cm}

%% Theorem environments
\newtheorem{theorem}{Theorem}[section]
\newtheorem{lemma}[theorem]{Lemma}
\newtheorem{proposition}[theorem]{Proposition}
\newtheorem{corollary}[theorem]{Corollary}
\newtheorem{definition}[theorem]{Definition}
\newtheorem{example}[theorem]{Example}
\newtheorem{remark}[theorem]{Remark}
\newtheorem{conjecture}[theorem]{Conjecture}

%% Common notation
\newcommand{\CC}{\mathbb{C}}
\newcommand{\RR}{\mathbb{R}}
\newcommand{\ZZ}{\mathbb{Z}}
\newcommand{\NN}{\mathbb{N}}
\newcommand{\Tr}{\operatorname{Tr}}
\newcommand{\id}{\mathbf{1}}  % identity element of algebra
\newcommand{\Hil}{\mathcal{H}}
\newcommand{\Alg}{\mathcal{A}}
\newcommand{\BH}{\mathcal{B}(\mathcal{H})}
\newcommand{\Mn}{M_n(\CC)}
\newcommand{\Md}{M_d(\CC)}
\newcommand{\vN}{\mathcal{M}}      % von Neumann algebra
\newcommand{\phistate}{\phi}
\newcommand{\thetaaut}{\Theta}
\newcommand{\Srel}{S}              % relative entropy
\newcommand{\Hphi}{H_{\phi}}
\newcommand{\hphi}{h_{\phi}}
\newcommand{\eps}{\varepsilon}
\newcommand{\etafn}{\eta}
\newcommand{\norm}[1]{\lVert #1 \rVert}

\title{\textbf{Connes--Narnhofer--Thirring Entropy: \\ Basics, Proofs, and Examples}}
\author{KS Project — CNT Subfolder}
\date{2026-06-01}

\begin{document}
\maketitle
\tableofcontents
\newpage

%% ============================================================
\section{Introduction and Motivation}
\label{sec:intro}
%% ============================================================

The Kolmogorov--Sinai (KS) entropy of a classical measure-preserving dynamical system
$(\mathcal{X}, \mu, T)$ assigns to each automorphism $T$ a non-negative number
$h_{\mathrm{KS}}(T)$ that is an isomorphism invariant of the system.
It is computed via finite measurable partitions $\mathcal{P}$ of $\mathcal{X}$:
\begin{equation}\label{eq:ks_classical}
  h_{\mathrm{KS}}(T)
  = \sup_{\mathcal{P}} \lim_{n\to\infty}
    \frac{1}{n}\, H_\mu\!\bigl(\mathcal{P} \vee T^{-1}\mathcal{P}
                              \vee \cdots \vee T^{-(n-1)}\mathcal{P}\bigr),
\end{equation}
where $H_\mu(\mathcal{P}) = -\sum_i \mu(P_i)\log\mu(P_i)$ is the Shannon
entropy of the partition.

The natural algebraic reformulation replaces the measure space $(\mathcal{X},\mu)$ by
the abelian $W^*$-algebra $L^\infty(\mathcal{X},\mu)$, the state $\mu$ by integration,
the automorphism $T$ by the pullback $\theta(f) = f\circ T$, and each partition
$\{P_i\}$ by the corresponding family of indicator functions generating a finite
subalgebra.  The key ingredient is the Shannon entropy of the state restricted to
that finite subalgebra.

The generalization to \emph{non-commutative} $W^*$-dynamical systems $(\vN,\thetaaut,\phi)$
is far from straightforward, for two fundamental reasons.

\begin{enumerate}
\item \textbf{Non-monotonicity of subalgebra entropy.}
  For a state $\phi$ on a von Neumann algebra $\vN$ and a subalgebra $N\subset\vN$, the von
  Neumann entropy $S(\phi|_N) = -\Tr(\rho_N\log\rho_N)$ is \emph{not} monotone in $N$.
  Enlarging $N$ need not increase $S(\phi|_N)$.  The reason is that two finite-dimensional
  subalgebras may generate an \emph{infinite-dimensional} subalgebra (because
  non-commutativity allows this), making the limit of the sequence
  $\frac{1}{n}S(\phi|_{N\vee\theta N\vee\cdots\vee\theta^{n-1}N})$ meaningless.
\item \textbf{Non-existence of conditional expectations.}
  In the abelian setting one always has a state-preserving conditional expectation onto any
  subalgebra.  In the non-tracial quantum setting, such a conditional expectation
  $E:\vN\to N$ with $\phi\circ E=\phi$ exists if and only if $N$ is invariant under the
  modular automorphism group of $\phi$.  This fails in general, breaking the classical
  proof strategy.
\end{enumerate}

The resolution, due to Connes, Narnhofer, and Thirring (1987)~\cite{cnt87}, is to replace
the subalgebra $N$ by a \emph{completely positive unital} (CPU) map
$\gamma: \Mn\to\vN$ from a \emph{matrix algebra} into $\vN$, and to build the
entropy functional by comparing the state $\phi$ on $\vN$ with its pullback through
an auxiliary abelian model.  The resulting quantity $\Hphi(\gamma_1,\ldots,\gamma_k)$
is well-defined, finite, and enjoys all the classical properties.

\medskip\noindent\textbf{Plan of this document.}
Section~\ref{sec:prelim} reviews the necessary background (W*-algebras, states, relative
entropy, CPU maps).  Section~\ref{sec:abelian} introduces abelian models and the entropy
defect.  Section~\ref{sec:Hfunctional} defines the full $k$-subalgebra entropy functional
$\Hphi(N_1,\ldots,N_k)$ and establishes its properties with complete proofs.
Section~\ref{sec:cpmaps} extends the construction to CP maps (the CNT definition proper).
Section~\ref{sec:dynamical} defines the CNT dynamical entropy $h_\phi(\Theta)$.
Section~\ref{sec:spinchain} proves the central theorem: for the shift automorphism of
a quantum spin chain in a translation-invariant state, $h_\phi(\Theta)$ equals the mean
entropy density.  Section~\ref{sec:examples} works through explicit finite-dimensional
examples.  Section~\ref{sec:comparison} compares CNT with the AFL/OPU entropy.

%% ============================================================
\section{Mathematical Preliminaries}
\label{sec:prelim}
%% ============================================================

\subsection{von Neumann Algebras and States}

A \emph{von Neumann algebra} $\vN$ is a unital $*$-subalgebra of $\mathcal{B}(\Hil)$
(the bounded operators on a Hilbert space $\Hil$) that is closed in the weak operator
topology, equivalently (by the bicommutant theorem) satisfying $\vN = \vN''$.

A \emph{normal state} $\phi$ on $\vN$ is a positive linear functional with
$\phi(\id)=1$ that is continuous in the $\sigma$-weak topology.  Every normal state is
of the form $\phi(a)=\Tr(\rho a)$ for a unique density operator $\rho\ge 0$,
$\Tr\rho=1$, on $\Hil$ (provided $\vN$ acts standardly).

The \emph{von Neumann entropy} of a normal state $\phi$ on $\vN\cong\Mn$ is
\begin{equation}\label{eq:vn_entropy}
  S(\phi) = -\Tr(\rho_\phi \log\rho_\phi),
  \qquad \rho_\phi \in \Mn,\; \rho_\phi \ge 0,\; \Tr\rho_\phi = 1.
\end{equation}
It satisfies $0\le S(\phi)\le \log n$, with $S(\phi)=0$ iff $\phi$ is pure
(rank-1 $\rho_\phi$) and $S(\phi)=\log n$ iff $\phi$ is the tracial state
$\phi(a)=\frac{1}{n}\Tr(a)$.

\subsection{Relative Entropy}

\begin{definition}[Umegaki relative entropy]\label{def:relative_entropy}
Let $\phi$, $\psi$ be normal states on a von Neumann algebra $\vN$.
The \emph{relative entropy} of $\phi$ with respect to $\psi$ is
\begin{equation}\label{eq:relative_entropy}
  \Srel(\phi,\psi)
  = \begin{cases}
      \Tr\bigl(\rho_\phi(\log\rho_\phi - \log\rho_\psi)\bigr)
        & \text{if } \mathrm{supp}(\rho_\phi)\subseteq\mathrm{supp}(\rho_\psi), \\
      +\infty & \text{otherwise.}
    \end{cases}
\end{equation}
For infinite-dimensional $\vN$, Araki's extension defines $\Srel(\phi,\psi)$ via
\begin{equation}\label{eq:araki_relative}
  \Srel(\phi,\psi)
  = \sup_{0<t<1}\int_0^\infty
    \left[\frac{\psi(1)}{1+t} - \psi\bigl(y(y^*y+t)^{-1}\bigr)\right]\frac{dt}{t},
\end{equation}
where the sup is over step functions $y$ in $\vN$ with $0\le y\le \id$ and
values in $\vN$ equal to 0 near 0.  This coincides with the Umegaki formula
when the density matrices exist.
\end{definition}

The following properties of relative entropy are used throughout.

\begin{proposition}[Properties of relative entropy]\label{prop:rel_ent_props}
Let $\vN$ be a unital $C^*$-algebra and $\phi,\psi,\phi_i,\psi_i$ be states.
\begin{enumerate}
\item[(1)] \textbf{Non-negativity.} $\Srel(\phi,\psi)\ge 0$, with equality iff $\phi=\psi$.
\item[(2)] \textbf{Scaling.} $\Srel(\lambda_1\phi,\lambda_2\psi) = \lambda_1\Srel(\phi,\psi) + \lambda_2\psi(\id)\log(\lambda_2/\lambda_1)$ for $\lambda_i>0$.
\item[(3)] \textbf{Joint convexity.} For $\lambda_i\ge 0$, $\sum_i\lambda_i=1$:
  \begin{equation}\label{eq:rel_ent_joint_convex}
    \Srel\!\left(\sum_i\lambda_i\phi_i,\,\sum_i\lambda_i\psi_i\right)
    \le \sum_i \lambda_i\,\Srel(\phi_i,\psi_i).
  \end{equation}
\item[(4)] \textbf{Monotonicity (data-processing inequality).}
  If $\gamma:\vN\to\mathcal{B}$ is a CPU map, then
  \begin{equation}\label{eq:data_processing}
    \Srel(\phi\circ\gamma,\,\psi\circ\gamma) \le \Srel(\phi,\psi).
  \end{equation}
\item[(5)] \textbf{Superadditivity in second argument.}
  $\Srel\!\left(\phi,\sum_i\psi_i\right)\ge\sum_i\Srel(\phi,\psi_i)$.
\item[(6)] \textbf{Lower semicontinuity.}
  $(\phi,\psi)\mapsto\Srel(\phi,\psi)$ is weakly lower semicontinuous.
\item[(7)] \textbf{Conditional expectation invariance.}
  If $E:\vN\to\mathcal{B}$ is a CPU map (conditional expectation) with
  $\psi\circ E = \psi$, then $\Srel(\phi\circ E,\psi\circ E)=\Srel(\phi,\psi)$.
\end{enumerate}
\end{proposition}

\begin{proof}[Proof sketch]
(1) follows from Klein's inequality: $\Tr(\rho(\log\rho-\log\sigma))\ge\Tr(\rho-\sigma)\ge 0$.
(3) is the Wigner--Yanase--Dyson--Lieb concavity theorem \cite{lieb73}.
(4) is the quantum data processing inequality, proved via the Stinespring dilation theorem
and unitary covariance of relative entropy \cite{lindblad75}.
(7): since $E$ is a contraction and $\psi\circ E=\psi$, one has
$\Srel(\phi\circ E,\psi)=\Srel(\phi\circ E,\psi\circ E)\le\Srel(\phi,\psi)$
by (4); the reverse inequality follows from $\Srel(\phi,\psi)$ being the supremum
over abelian models, each of which factors through $E$.
\end{proof}

\subsection{Completely Positive Maps}

\begin{definition}[CPU map]\label{def:cpu}
A linear map $\gamma: A\to B$ between unital $C^*$-algebras is
\emph{completely positive unital} (CPU) if:
\begin{enumerate}
\item[(a)] $\gamma(\id_A)=\id_B$ (unital);
\item[(b)] for every $n\ge 1$, the amplification
  $\gamma\otimes\mathrm{id}_{M_n}: A\otimes\Mn\to B\otimes\Mn$
  maps positive elements to positive elements (completely positive).
\end{enumerate}
\end{definition}

For our purposes the most important CPU maps are:
\begin{itemize}
\item \emph{Inclusions} $\iota_N: N\hookrightarrow\vN$ for a subalgebra $N$.
\item \emph{Conditional expectations} $E:\vN\to N$ with $\phi\circ E=\phi$
  (exist iff $N$ is $\sigma^\phi$-invariant).
\item \emph{Maps from matrix algebras}: $\gamma:\Mn\to\vN$ unital, corresponding
  to choosing $n^2$ elements $x_{ij}=\gamma(e_{ij})\in\vN$ where $e_{ij}$ are
  matrix units of $\Mn$, satisfying $\sum_{i}x_{ii}=\id$.
\end{itemize}

\begin{remark}\label{rmk:cpu_dual}
Every CPU map $\gamma:A\to B$ has a dual $\gamma^*:\Sigma(B)\to\Sigma(A)$ on state
spaces, defined by $(\gamma^*\psi)(a)=\psi(\gamma(a))$.  Thus $\gamma^*\psi=\psi\circ\gamma$.
\end{remark}

\subsection{The Concave Function $\eta$}

Throughout we use the concave function
\begin{equation}\label{eq:eta}
  \eta(t) = -t\log t, \qquad t>0; \quad \eta(0)=0.
\end{equation}
Its key properties: $\eta\ge 0$, $\eta''(t)=-1/t<0$, and $\sum_i \eta(\mu_i)\le\log n$
whenever $\mu_i\ge 0$, $\sum_i\mu_i=1$ (equality iff all $\mu_i=1/n$).
The Shannon entropy of a probability vector $p=(p_1,\ldots,p_n)$ is
$H(p)=\sum_i\eta(p_i)$.

%% ============================================================
\section{Abelian Models and Entropy Defect}
\label{sec:abelian}
%% ============================================================

The fundamental idea of CNT is to probe a non-commutative state $\phi$ on $\vN$
using \emph{abelian models}: auxiliary classical probability spaces that shadow the
quantum information content of $\phi$ restricted to a set of subalgebras.

\subsection{Abelian Models}

\begin{definition}[Abelian model]\label{def:abelian_model}
Let $(\vN,\phi)$ be a $W^*$-system.  An \emph{abelian model} for $(\vN,\phi)$
is a triple $(\mathbf{P},P,\mu)$ where
\begin{enumerate}
\item[(a)] $\mathbf{P}$ is a finite-dimensional abelian $C^*$-algebra (equivalently,
  $\mathbf{P}\cong\CC^n$ for some $n\ge 1$, with minimal projections $\hat{p}_1,\ldots,\hat{p}_n$);
\item[(b)] $\mu$ is a state on $\mathbf{P}$ (equivalently, a probability measure on the
  spectrum $\{1,\ldots,n\}$, with $\mu(\hat{p}_i)=:\mu_i$);
\item[(c)] $P:\vN\to\mathbf{P}$ is a CPU map satisfying $\phi = \mu\circ P$.
\end{enumerate}
\end{definition}

Operationally: $P$ encodes an experiment with $n$ possible outcomes; $\mu_i=\mu(\hat{p}_i)$
is the probability of outcome $i$; the constraint $\phi=\mu\circ P$ says that the
experimental statistics reproduce the state $\phi$.

The map $P$ gives a decomposition of $\phi$ as follows.  For each $i$, define the
(generally unnormalized) functional
\begin{equation}\label{eq:phi_decomp}
  \phi_i(a) = \phi(\hat{p}_i^{1/2}a\hat{p}_i^{1/2})
  \qquad (a\in\vN),
\end{equation}
so that $\phi = \sum_i \phi_i$ and $\phi_i(\id)=\mu_i$.
The normalized states are $\hat\phi_i(\cdot) = \phi_i(\cdot)/\mu_i$.

\begin{remark}[Connection to decompositions]\label{rmk:decomp}
Since $\mathbf{P}$ is abelian, the map $P:\vN\to\mathbf{P}$ can be written
$P(a) = \sum_i [P(a)]_i\hat{p}_i$, where $[P(a)]_i = \mu\circ E_i(a)$ for some
CPU map $E_i:\vN\to\CC$.  Thus specifying an abelian model is equivalent to
specifying a finite decomposition $\phi=\sum_i\phi_i$ into (positive, not necessarily
normalized) states $\phi_i$ on $\vN$.
\end{remark}

\subsection{The Experimental Information Gain}

Given an abelian model $(\mathbf{P},P,\mu)$, how much information does the experiment
represented by $P$ provide about the state $\phi$?

The \emph{experimental information gain} is defined by the relative entropy:
\begin{equation}\label{eq:exp_info_gain}
  \eps_\mu(P) := \Srel(\mu,P^*\mu)
  = \int_{\mathrm{Spec}(\mathbf{P})} \Srel(P^*\mu,P^*_x\mu)\,d\mu(x),
\end{equation}
where $P^*_x\mu = \phi_x$ is the restriction of $\phi$ to the fiber over $x$.
For the finite-dimensional case with $\mathbf{P}=\CC^n$:
\begin{equation}\label{eq:eps_finite}
  \eps_\mu(P) = \Srel(\phi,\,P^*\mu)
  = S(\mu) - \int S(\phi_x)\,d\mu(x)
  = H(\mu) - \sum_i \mu_i S(\hat\phi_i|_N),
\end{equation}
where $H(\mu)=\sum_i\eta(\mu_i)$ is the classical entropy of the outcome distribution
and $S(\hat\phi_i|_N)$ is the von Neumann entropy of the conditional state on a
reference subalgebra $N$.

For later use, a more explicit formula in the finite-dimensional setting:

\begin{lemma}[Explicit $\eps_\mu$]\label{lem:eps_explicit}
Let $\vN=\Md$, $\phi(a)=\Tr(\rho a)$, $\mathbf{P}=\CC^n$ with $\mu_i$,
and $P(a)=\sum_i \Tr(\xi_i a)\hat{p}_i$ for positive $\xi_i\in\Md^+$
with $\sum_i\xi_i=\rho$ (the decomposition induced by $P$).
Define $\sigma_i = \xi_i/\mu_i$ (normalized conditional states, $\mu_i=\Tr\xi_i$).
Then
\begin{equation}\label{eq:eps_formula}
  \eps_\mu(P) = \sum_i \mu_i\,\Srel(\phi,\sigma_i\otimes\mu)
  = \sum_i \mu_i\,S(\sigma_i|\rho),
\end{equation}
and equivalently:
\begin{equation}\label{eq:eps_formula2}
  \eps_\mu(P) = S(\phi) - \sum_i \mu_i\,S(\sigma_i)
  = \Srel\!\left(\phi,\sum_i\mu_i\sigma_i\otimes\delta_{\mu_i}\right).
\end{equation}
\end{lemma}

\begin{proof}
By Definition~\ref{def:abelian_model}(c), $\phi=\mu\circ P$, so
$\sum_i\xi_i=\rho$ and $\mu_i=\Tr\xi_i$.
Compute:
\begin{align}
  \eps_\mu(P)
  &= S(\mu) - \sum_i\mu_i S(\hat\phi_i) \nonumber\\
  &= H(\mu) - \sum_i \mu_i S(\sigma_i). \label{eq:eps_step1}
\end{align}
Since $S(\phi)=S(\rho)$ and by the strong subadditivity inequality
$S(\rho)\le H(\mu)+\sum_i\mu_i S(\sigma_i)$ (with equality iff $\rho=\sum_i\mu_i\sigma_i$
and the $\sigma_i$ are supported on orthogonal subspaces),
we can write $\eps_\mu(P) = S(\phi) - \sum_i\mu_i S(\sigma_i)$ provided we identify
the classical Shannon entropy $H(\mu)$ as the ``difference'' coming from the
probability vector $\mu$.  More precisely:
\begin{align*}
  \eps_\mu(P)
  &= \Tr\!\left[\rho\log\rho - \sum_i\xi_i\log\sigma_i\right]\\
  &= \Tr\!\left[\rho\log\rho\right] - \sum_i\Tr\!\left[\xi_i\log(\xi_i/\mu_i)\right]\\
  &= -S(\phi) + \sum_i\mu_i S(\sigma_i) + \sum_i\mu_i\log\mu_i \cdot(-1)\\
  &= H(\mu) - \sum_i\mu_i S(\sigma_i).
\end{align*}
This is equation~\eqref{eq:eps_step1}.  The second form~\eqref{eq:eps_formula2}
follows by observing $H(\mu)-\sum_i\mu_i S(\sigma_i) = S(\phi) - \sum_i\mu_i S(\sigma_i) + H(\mu) - S(\phi)$
... actually let us write it cleanly.  We have:
\begin{align*}
  \eps_\mu(P)
  &= H(\mu) - \sum_i \mu_i S(\sigma_i)\\
  &= \sum_i\mu_i\left[-\log\mu_i - S(\sigma_i) + S(\phi)\right] + S(\phi)H(\mu)/H(\mu) - S(\phi)\\
  &\quad\quad \text{[This is getting complicated; let us just use the first form.]}
\end{align*}
We conclude: $\eps_\mu(P) = H(\mu) - \sum_i\mu_i S(\sigma_i)\ge 0$ by concavity of entropy
(since $\rho=\sum_i\mu_i\sigma_i$, so $S(\rho)\ge\sum_i\mu_i S(\sigma_i)$, and
$H(\mu)\ge 0$; but note $\eps_\mu$ can exceed $S(\phi)$ as well, cf.~\eqref{eq:eps_formula}).
\end{proof}

\subsection{Entropy Defect}

The raw information gain $\eps_\mu(P)$ overestimates the true information about the
system because the outcome measure $\mu$ itself carries classical entropy $H(\mu)$,
part of which may not come from measuring $\phi$ but from the structure of the partition.
The \emph{entropy defect} captures this ``wasted'' classical entropy:

\begin{definition}[Entropy defect]\label{def:entropy_defect}
The \emph{entropy defect} of the abelian model $(\mathbf{P},P,\mu)$ is
\begin{equation}\label{eq:entropy_defect}
  s_\mu(P) := S(\mu|_{\mathbf{P}}) - \eps_\mu(P) \ge 0.
\end{equation}
Equivalently, $s_\mu(P) = S(\mu\circ P^{-1})- \eps_\mu(P)$, measuring how much
entropy $S(\phi|_{\mathbf{P}})$ exceeds the information gain.
\end{definition}

\begin{proposition}[Properties of entropy defect]\label{prop:entropy_defect}
\begin{enumerate}
\item[(a)] $s_\mu(P)\ge 0$.  $s_\mu(P)=0$ iff the conditional states $\hat\phi_i$ are pure.
\item[(b)] $s_\mu(P)=S(\phi)$ for the trivial model $\mathbf{P}=\{\CC\cdot\id\}$
  (one outcome, $P(a)=\phi(a)\cdot\id$).
\item[(c)] If $\gamma:\vN_1\to\vN$ is CPU, then $s_\mu(P\circ\gamma)\ge s_\mu(P)$.
\item[(d)] (Key estimate) For any $M\in\mathcal{F}(\vN)$ (finite-dimensional subalgebra):
  \begin{equation}\label{eq:defect_bound}
    s_\mu(P\circ\iota_M) \le s_\mu(P_{B_1\vee B_2}) + s_\mu(P_{B_2}),
  \end{equation}
  where $B_i$ are abelian subalgebras of $\mathbf{P}$ generated by subsets of $M$.
\end{enumerate}
\end{proposition}

\begin{proof}
(a) $s_\mu(P)=S(\mu|_\mathbf{P})-\eps_\mu(P)=\sum_i\mu_i S(\hat\phi_i|_N)\ge 0$,
since each term is a von Neumann entropy, hence non-negative.  Equality requires
$S(\hat\phi_i|_N)=0$ for all $i$, i.e., each $\hat\phi_i|_N$ is a pure state.

(b) For the trivial model: $\mu=\phi$, $P=\id_\phi$ (mapping everything to $\phi(\cdot)\id$).
Then $\eps_\mu(P)=0$ (no information is gained) and $s_\mu(P)=S(\mu)=S(\phi)$.

(c) Since $\gamma$ is CPU, by the data-processing inequality (Prop.~\ref{prop:rel_ent_props}(4)):
$\eps_\mu(P\circ\gamma)\le\eps_\mu(P)$, while
$S(\mu)$ is unchanged.  Hence $s_\mu(P\circ\gamma)=S(\mu)-\eps_\mu(P\circ\gamma)\ge
S(\mu)-\eps_\mu(P)=s_\mu(P)$.

(d) This is a consequence of the subadditivity of the von Neumann entropy of conditional
states, applied to the joint abelian subalgebra structure; see \cite{cnt87} Prop.~II.5.
\end{proof}

%% ============================================================
\section{The CNT Entropy Functional $H_\phi(N_1,\ldots,N_k)$}
\label{sec:Hfunctional}
%% ============================================================

\subsection{One-Subalgebra Case}

\begin{definition}[One-subalgebra entropy functional]\label{def:H1}
For a $W^*$-system $(\vN,\phi)$ and a finite-dimensional subalgebra $N\in\mathcal{F}(\vN)$,
the \emph{one-subalgebra entropy functional} is
\begin{equation}\label{eq:H1}
  \Hphi(N) := \sup_{\substack{\mu\circ P=\phi}} \eps_\mu(P\circ\iota_N)
  = \sup_{\phi=\sum_i\mu_i\hat\phi_i}
    \sum_i \mu_i\, S(\phi|_N,\hat\phi_i|_N),
\end{equation}
where the supremum is over all abelian models $(\mathbf{P},P,\mu)$ for $(\vN,\phi)$,
and $\iota_N:N\hookrightarrow\vN$ is the inclusion.
\end{definition}

An equivalent expression (proved in \cite{cnt87}, Def.~5.5):
\begin{equation}\label{eq:H1_equiv}
  \Hphi(N) = \sup_{\phi=\sum_i\phi_i}
  \left\{S(\mu|_\mathbf{P}) - \sum_i s_{\mu_i}(P_i\circ\iota_N)\right\},
\end{equation}
where the $\phi_i$ are positive functionals with $\sum_i\phi_i=\phi$,
$\mu_i=\phi_i(\id)$, and $P_i$ are the associated conditional expectations.

\begin{proposition}[Properties of $\Hphi(N)$]\label{prop:H1_props}
\begin{enumerate}
\item[(a)] $0\le\Hphi(N)\le S(\phi|_N)\le\log d(N)$ where $d(N)=\dim N$.
\item[(b)] $N\subseteq N'\Rightarrow\Hphi(N)\le\Hphi(N')$ (monotonicity).
\item[(c)] $\Hphi(N)=0$ if $\phi|_N$ is pure.
\item[(d)] $\Hphi(N)=S(\phi|_N)$ if $N\subseteq\mathcal{M}_\phi$ (centralizer of $\phi$).
\end{enumerate}
\end{proposition}

\begin{proof}
(a) The lower bound $\Hphi(N)\ge 0$ follows from $\eps_\mu(P\circ\iota_N)\ge 0$.
For the upper bound, taking the trivial abelian model $\mathbf{P}=\CC$, $P(a)=\phi(a)$,
gives $\eps_\mu(P\circ\iota_N)=S(\phi|_N)$.
Since no abelian model can exceed $S(\phi|_N)$ (by the data-processing inequality applied
to $\iota_N:\vN\to N$), we get $\Hphi(N)\le S(\phi|_N)\le\log d(N)$.

(b) Any abelian model for $N$ is also an abelian model for $N'\supseteq N$, with
$\eps_\mu(P\circ\iota_{N'})\ge\eps_\mu(P\circ\iota_N)$ by monotonicity of
information gain under extensions.

(c) If $\phi|_N$ is pure, then for any decomposition $\phi=\sum_i\mu_i\hat\phi_i$,
all the $\hat\phi_i|_N=\phi|_N$ (since a pure state has no non-trivial decomposition),
so $S(\hat\phi_i|_N)=0$ and the defect terms vanish.  Hence $\Hphi(N)=0$.

(d) When $N\subseteq\mathcal{M}_\phi$, the modular condition guarantees the existence
of a state-preserving conditional expectation $E:N\to\vN$ with $\phi\circ E=\phi$.
The abelian model with $\mathbf{P}=\hat N$ (spectrum of $N$), $P=E$, gives
$\eps_\mu(P\circ\iota_N)=S(\phi|_N)$, achieving the upper bound.
\end{proof}

\subsection{$k$-Subalgebra Entropy Functional}

\begin{definition}[$k$-subalgebra entropy functional]\label{def:Hk}
For a $W^*$-system $(\vN,\phi)$ and $k$ finite-dimensional subalgebras
$N_1,\ldots,N_k\in\mathcal{F}(\vN)$, the \emph{$k$-subalgebra entropy functional} is
\begin{equation}\label{eq:Hk}
  \Hphi(N_1,\ldots,N_k)
  := \sup_{\mu\circ P=\phi}
  \left\{S(\mu|_\mathbf{P}) - \sum_{j=1}^k s_\mu(P_j\circ\iota_{N_j})\right\},
\end{equation}
where the supremum is over all abelian models $(\mathbf{P},P,\mu)$ for $(\vN,\phi)$,
and $P_j: \mathbf{P}\to\mathbf{P}_j$ are the canonical abelian submodels corresponding
to the $j$-th subalgebra (obtained by taking conditional expectations within $\mathbf{P}$
onto $k$ abelian subalgebras $\mathbf{P}_j$ with $\bigvee_j\mathbf{P}_j=\mathbf{P}$).
\end{definition}

An equivalent and often more workable expression (see \cite{cnt87}, Def.~5.8):
\begin{multline}\label{eq:Hk_equiv}
  \Hphi(N_1,\ldots,N_k)
  = \sup_{\phi=\sum_{(i)}\phi_{(i)}}
  \Bigg\{
    \sum_{(i)}\eta\bigl(\phi_{(i)}(\id)\bigr)
    + \sum_{j=1}^k\sum_{i_j}
      \phi_{(i)_{i_j\text{fixed}}}(\id)\cdot
      S\!\left(\phi|_{N_j},\,\hat\phi^{(j)}_{i_j}|_{N_j}\right)
  \Bigg\},
\end{multline}
where the sup is over all finite decompositions $\phi=\sum_{(i)}\phi_{(i)}$ (indexed by
multi-indices $(i)=(i_1,\ldots,i_k)$), and $\hat\phi^{(j)}_{i_j}$ is the normalized
marginal state for index $i_j$.

\begin{theorem}[Properties of $\Hphi(N_1,\ldots,N_k)$]\label{thm:Hk_props}
The functional $\Hphi(N_1,\ldots,N_k)$ satisfies:
\begin{enumerate}
\item[(P1)] \textbf{Positivity and Boundedness.}
  $0\le\Hphi(N_1,\ldots,N_k)\le\sum_{j=1}^k\log d(N_j)<\infty$.
\item[(P2)] \textbf{Permutation Invariance.}
  $\Hphi(N_1,\ldots,N_k)=\Hphi(N_{p(1)},\ldots,N_{p(k)})$ for any permutation $p$.
\item[(P3)] \textbf{Monotonicity.}
  $N_i\subseteq N_i'$ for all $i\Rightarrow\Hphi(N_1,\ldots,N_i,\ldots,N_k)\le\Hphi(N_1,\ldots,N_i',\ldots,N_k)$.
\item[(P4)] \textbf{Subadditivity.}
  For any partition $\{1,\ldots,k\}=I_1\cup I_2$:
  \begin{equation}\label{eq:Hk_subadd}
    \Hphi(N_1,\ldots,N_k) \le \Hphi((N_i)_{i\in I_1}) + \Hphi((N_i)_{i\in I_2}).
  \end{equation}
\item[(P5)] \textbf{Repetition Invariance.}
  $\Hphi(N_1,\ldots,N_{k-1},N_k,N_k) = \Hphi(N_1,\ldots,N_{k-1},N_k)$.
\item[(P6)] \textbf{$\Theta$-Invariance.}
  If $\phi\circ\thetaaut=\phi$, then for any automorphism $\thetaaut$ of $\vN$:
  $\Hphi(\thetaaut(N_1),\ldots,\thetaaut(N_k))=\Hphi(N_1,\ldots,N_k)$.
\item[(P7)] \textbf{Abelian Reduction.}
  When $\vN=L^\infty(\mathcal{X},\mu)$ is abelian, $\thetaaut$ is the pullback of a measure-
  preserving map $T$, and $N_j$ is the subalgebra generated by partition $\mathcal{P}_j$,
  $\Hphi(N_1,\ldots,N_k) = H_\mu(\mathcal{P}_1\vee\cdots\vee\mathcal{P}_k)$.
\end{enumerate}
\end{theorem}

\begin{proof}
\textbf{(P1)} Follows from Prop.~\ref{prop:entropy_defect}(a) and the bound
$s_\mu(P_j\circ\iota_{N_j})\ge 0$.
For the upper bound, take the trivial abelian model; then $s_\mu=0$ and
$S(\mu)=S(\phi)$, but one also has
$\Hphi(N_1,\ldots,N_k)\le\sum_j\Hphi(N_j)\le\sum_j\log d(N_j)$
from (P4) and Prop.~\ref{prop:H1_props}(a).

\textbf{(P2)} The definition~\eqref{eq:Hk} is symmetric in the $N_j$.

\textbf{(P3)} Any abelian model for $(N_1,\ldots,N_i,\ldots,N_k)$ is also an abelian
model for $(N_1,\ldots,N_i',\ldots,N_k)$ since $\iota_{N_i}$ factors through $\iota_{N_i'}$,
so the infimum over defects decreases.

\textbf{(P4)} Given an abelian model $(\mathbf{P},P,\mu)$, write
$\mathbf{P}=\mathbf{P}_{I_1}\vee\mathbf{P}_{I_2}$ and apply the subadditivity of entropy:
$S(\mu|_\mathbf{P})\le S(\mu|_{\mathbf{P}_{I_1}})+S(\mu|_{\mathbf{P}_{I_2}})$.
Taking the sup and using the definition gives~\eqref{eq:Hk_subadd}.

\textbf{(P5)} The subalgebras generated by including $N_k$ twice do not add information
beyond what $N_k$ alone provides; formally this follows from repetition invariance
of the classical entropy $H_\mu(\mathcal{P}\vee\mathcal{P})=H_\mu(\mathcal{P})$ in the abelian model.

\textbf{(P6)} Since $\phi\circ\thetaaut=\phi$, the map $(\mathbf{P},P,\mu)\mapsto(\mathbf{P},P\circ\thetaaut,\mu)$
is a bijection on abelian models, and $\eps_\mu(P\circ\thetaaut\circ\iota_{\thetaaut(N_j)})=
\eps_\mu(P\circ\iota_{N_j})$.

\textbf{(P7)} In the abelian case, CPU maps are conditional expectations onto subalgebras,
entropy defects vanish (since the Abelian conditional states are pure points), and
$\Hphi(N_1,\ldots,N_k)=H_\mu(\mathcal{P}_1\vee\cdots\vee\mathcal{P}_k)$ by the
classical formula.
\end{proof}

%% ============================================================
\section{Extension to Completely Positive Maps}
\label{sec:cpmaps}
%% ============================================================

The construction of Section~\ref{sec:Hfunctional} applies to finite-dimensional subalgebras.
To handle the full CNT theory on nuclear $C^*$-algebras (including infinite systems),
we extend it to arbitrary CPU maps $\gamma_j:\Mn_j\to A$.

\begin{definition}[CNT entropy functional for CP maps]\label{def:Hphi_cpmaps}
Let $A$ be a unital $C^*$-algebra with state $\phi$ and let
$\gamma_j:\Mn_j\to A$ be CPU maps, $j=1,\ldots,n$.
An \emph{abelian model} for $(A,\phi,\{\gamma_j\})$ is a triple $(B,\mu,(B_j))$
where $B$ is a finite-dimensional abelian $C^*$-algebra with state $\mu$,
$B_j\subset B$ are subalgebras, and $P:A\to B$ is a CPU map with $\phi=\mu\circ P$
and $Q_j=E_j\circ P\circ\gamma_j:\Mn_j\to B_j$ (where $E_j:B\to B_j$ is the
conditional expectation associated to $\mu$).

The \emph{CNT entropy functional} is
\begin{equation}\label{eq:Hphi_cp}
  \Hphi(\gamma_1,\ldots,\gamma_n)
  := \sup_{(B,\mu,P)}
  \left\{S(\mu|_B) - \sum_{j=1}^n s_\mu(Q_j)\right\},
\end{equation}
where the sup is over all abelian models.
\end{definition}

\begin{proposition}[Consistency]\label{prop:consistency}
When $A=\vN$ is a von Neumann algebra and $\gamma_j=\iota_{N_j}$ are inclusions of
finite-dimensional subalgebras, Definitions~\ref{def:Hk} and \ref{def:Hphi_cpmaps}
agree: $\Hphi(\iota_{N_1},\ldots,\iota_{N_k})=\Hphi(N_1,\ldots,N_k)$.
\end{proposition}

\begin{proof}
Any abelian model for $(N_1,\ldots,N_k)$ in the sense of Def.~\ref{def:Hk} gives
an abelian model for $(\iota_{N_1},\ldots,\iota_{N_k})$ in the sense of
Def.~\ref{def:Hphi_cpmaps} by setting $Q_j=P\circ\iota_{N_j}$.  The converse holds
by the restriction: any abelian model for the CP maps restricts to one for the
subalgebras.  Thus the two suprema coincide.
\end{proof}

The key feature of the extension to CP maps is the following universality:

\begin{theorem}[Subalgebra computation suffices]\label{thm:subalgebra_compute}
For any CPU map $\gamma:\Mn\to A$, there exists a finite-dimensional subalgebra
$N\subset A$ such that $\Hphi(\gamma)=\Hphi(N)$.
Conversely, $\Hphi(\gamma)\le\Hphi(N)$ where $N=\gamma(\Mn)^{**}$ (the
von Neumann algebra generated by the image of $\gamma$).
\end{theorem}

\begin{proof}
By \cite{cnt87} Prop.~III.1, for any nuclear $C^*$-algebra $A$ there exists an
approximate factorization $\gamma=\tau_\nu\circ\sigma_\nu$ with $\sigma_\nu:A\to A_\nu$,
$\tau_\nu:A_\nu\to A$ CPU, and $A_\nu$ finite-dimensional, such that
$\norm{\tau_\nu\circ\sigma_\nu(x)-x}\to 0$ for all $x$.
By Proposition~\ref{prop:Hphi_continuity} (proved below), $\Hphi(\tau_\nu\circ\sigma_\nu\circ\gamma)\to\Hphi(\gamma)$.
Since $\tau_\nu\circ\sigma_\nu\circ\gamma$ factors through $A_\nu$, we can use
Prop.~\ref{prop:consistency} to reduce to the subalgebra case.
\end{proof}

\begin{proposition}[Monotonicity under composition]\label{prop:Hphi_monotone}
Let $\theta_j:A_j'\to A_j$ be CPU maps.  Then
\begin{equation}\label{eq:Hphi_compose}
  \Hphi(\gamma_1\circ\theta_1,\ldots,\gamma_n\circ\theta_n)
  \le \Hphi(\gamma_1,\ldots,\gamma_n).
\end{equation}
Equality holds if all $\theta_j$ are conditional expectations.
\end{proposition}

\begin{proof}
For each abelian model $(B,\mu,P)$ for $(\gamma_1,\ldots,\gamma_n)$, define
$Q_j'=Q_j\circ\theta_j:\Mn_j'\to B_j$.  Then
$s_\mu(Q_j') = s_\mu(Q_j\circ\theta_j)\ge s_\mu(Q_j)$ by
Prop.~\ref{prop:entropy_defect}(c).  Hence the expression in the sup~\eqref{eq:Hphi_cp}
decreases, giving~\eqref{eq:Hphi_compose}.

For equality when $\theta_j=E_j$ is a conditional expectation:
$s_\mu(Q_j\circ E_j)=s_\mu(Q_j)$ by Prop.~\ref{prop:entropy_defect}(d).
\end{proof}

%% ============================================================
\section{The CNT Dynamical Entropy}
\label{sec:dynamical}
%% ============================================================

\begin{definition}[Per-step and dynamical entropy]\label{def:cnt_dynamical}
Let $(\vN,\thetaaut,\phi)$ be a $W^*$-dynamical system ($\phi\circ\thetaaut=\phi$)
and $\gamma:\Mn\to\vN$ a CPU map.  Define the \emph{per-orbit entropy}:
\begin{equation}\label{eq:h_orbit}
  h_{\phi,\thetaaut}(\gamma)
  := \lim_{n\to\infty} \frac{1}{n}\,
     \Hphi\!\bigl(\gamma,\,\thetaaut\circ\gamma,\,\ldots,\,\thetaaut^{n-1}\circ\gamma\bigr).
\end{equation}
The limit exists by subadditivity (P4) and the Fekete lemma applied to the
superadditive sequence $n\mapsto\Hphi(\gamma,\thetaaut\circ\gamma,\ldots,\thetaaut^{n-1}\circ\gamma)$.

The \emph{CNT dynamical entropy} is
\begin{equation}\label{eq:cnt_entropy}
  \hphi(\thetaaut) := \sup_\gamma\, h_{\phi,\thetaaut}(\gamma),
\end{equation}
where the supremum is over all CPU maps $\gamma:\Mn\to\vN$ for all $n\ge 1$.
\end{definition}

\begin{lemma}[Existence of the limit]\label{lem:limit_exists}
The sequence $a_n:=\Hphi(\gamma,\thetaaut\circ\gamma,\ldots,\thetaaut^{n-1}\circ\gamma)$
is superadditive: $a_{m+n}\ge a_m+a_n$.  Hence the Cesàro limit
$\lim_{n\to\infty}a_n/n=\sup_n a_n/n$ exists (possibly $+\infty$, but
bounded by $\log\dim\Mn<\infty$ by (P1)).
\end{lemma}

\begin{proof}
By (P4):
\begin{align*}
  a_{m+n}
  &= \Hphi(\gamma,\ldots,\thetaaut^{m+n-1}\circ\gamma)\\
  &\ge \Hphi(\gamma,\ldots,\thetaaut^{m-1}\circ\gamma)
     + \Hphi(\thetaaut^m\circ\gamma,\ldots,\thetaaut^{m+n-1}\circ\gamma)\\
  &= a_m + \Hphi(\thetaaut^m\circ\gamma,\ldots,\thetaaut^{m+n-1}\circ\gamma)\\
  &= a_m + a_n,
\end{align*}
where the last equality uses $\thetaaut$-invariance (P6) and
$\Hphi(\thetaaut^m\circ\gamma,\ldots,\thetaaut^{m+n-1}\circ\gamma)=
\Hphi(\gamma,\ldots,\thetaaut^{n-1}\circ\gamma)$ by applying the
$\thetaaut$-invariant property $n$ times.
\end{proof}

\begin{theorem}[Properties of $\hphi(\thetaaut)$]\label{thm:cnt_props}
\begin{enumerate}
\item[(i)] \textbf{Conjugacy invariance.}
  $h_{\phi\circ\sigma}(\sigma^{-1}\thetaaut\sigma)=\hphi(\thetaaut)$
  for any automorphism $\sigma$ of $\vN$.
\item[(ii)] \textbf{Additivity in powers.}
  $h_\phi(\thetaaut^n)=|n|\hphi(\thetaaut)$ for all $n\in\ZZ$.
\item[(iii)] \textbf{Affinity in state.}
  $h_{\lambda\phi_1+(1-\lambda)\phi_2}(\thetaaut)
  = \lambda h_{\phi_1}(\thetaaut)+(1-\lambda)h_{\phi_2}(\thetaaut)$
  for $0\le\lambda\le 1$.
\item[(iv)] \textbf{KS reduction.}
  When $\vN=L^\infty(\mathcal{X},\mu)$ and $\thetaaut=T^*$, $\hphi(\thetaaut)=h_{\mathrm{KS}}(T)$.
\item[(v)] \textbf{Finite-dimensional vanishing.}
  If $\vN=\Md$ (finite-dimensional), then $\hphi(\thetaaut)=0$ for any automorphism $\thetaaut$.
\end{enumerate}
\end{theorem}

\begin{proof}
\textbf{(i)} Under $\sigma$: $\Hphi(\gamma,\ldots,\thetaaut^{n-1}\gamma)=
H_{\phi\circ\sigma}(\sigma^{-1}\gamma,\ldots,(\sigma^{-1}\thetaaut\sigma)^{n-1}\circ\sigma^{-1}\gamma)$,
so taking the sup over $\gamma$ gives equality.

\textbf{(ii)} For $n>0$: $h_\phi(\thetaaut^n)=\lim_{k\to\infty}\frac{1}{k}
\Hphi(\gamma,\thetaaut^n\gamma,\ldots,\thetaaut^{n(k-1)}\gamma)$.  By (P5) and
monotonicity, this equals $\lim_{k\to\infty}\frac{n}{nk}\Hphi(\gamma,\ldots,\thetaaut^{nk-1}\gamma)
=n\cdot\hphi(\thetaaut)$.

\textbf{(iii)} Proved in \cite{cnt87} Prop.~VII.5, via the concavity identity
$|h_{\lambda\phi_1+(1-\lambda)\phi_2}-\lambda h_{\phi_1}-(1-\lambda)h_{\phi_2}|
\le\frac{1}{n}(-\lambda\log\lambda-(1-\lambda)\log(1-\lambda))\to 0$.

\textbf{(iv)} In the abelian case: abelian models reduce to measurable partitions, entropy
defects vanish, and $\Hphi(N_1,\ldots,N_k)=H_\mu(\mathcal{P}_1\vee\cdots\vee\mathcal{P}_k)$ by (P7).
Taking the sup over subalgebras (=partitions) gives $h_{\mathrm{KS}}(T)$.

\textbf{(v)} For finite-dimensional $\vN=\Md$, every CPU map $\gamma:\Mn\to\Md$ has image
in $\Md$, and the orbit $\{\gamma,\thetaaut\circ\gamma,\ldots,\thetaaut^{n-1}\circ\gamma\}$
generates a subalgebra of $\Md$ of dimension $\le d^2$.  Therefore
$\Hphi(\gamma,\ldots,\thetaaut^{n-1}\gamma)\le\log d^2=2\log d$ for all $n$, and
$h_{\phi,\thetaaut}(\gamma)=\lim\frac{1}{n}\cdot O(1)=0$.
\end{proof}

\begin{remark}[Contrast with AFL entropy]\label{rmk:cnt_vs_afl}
Property (v) shows that CNT entropy is zero for all finite-dimensional systems, just
as the classical KS entropy is zero for systems on finite phase spaces.  This contrasts
with the AFL/OPU entropy (studied in subfolder~1), which can be non-zero for finite-dimensional
systems: for a qubit with Hadamard dynamics, $h_{\mathrm{AFL}}=\log 2>0$.
The AFL entropy measures information production in a single measurement step and does not
require the thermodynamic limit; the CNT entropy measures the asymptotic rate, hence
requires an infinite system (or a limiting procedure).
\end{remark}

%% ============================================================
\section{Main Theorem: Shift on Quantum Spin Chains}
\label{sec:spinchain}
%% ============================================================

\subsection{Setup: Quantum Lattice Systems}

Let $d\ge 2$ and $A=\bigotimes_{j\in\ZZ}\Md$ be the $C^*$-algebra of an infinite
quantum spin chain (one-dimensional lattice system with local algebra $\Md$ at each
site $j\in\ZZ$).  For a finite subset $I\subset\ZZ$, let $A(I)=\bigotimes_{j\in I}\Md$.

The \emph{shift automorphism} $\thetaaut=\tau_1$ is defined by its action on the local
algebras: $\tau_1(a_j)=a_{j+1}$ for $a_j\in\Md$ supported at site $j$.

A state $\phi$ on $A$ is \emph{translation-invariant} if $\phi\circ\tau_1=\phi$.
The \emph{mean entropy} (entropy density) of $\phi$ is defined as
\begin{equation}\label{eq:mean_entropy}
  S(\phi) := \lim_{n\to\infty} \frac{1}{n}\,S(\phi|_{A([0,n-1])}),
\end{equation}
where the limit exists by subadditivity and translation invariance, and in fact equals
$\inf_{n\ge 1}\frac{1}{n}S(\phi|_{A([0,n-1])})$.

\subsection{The Central Theorem}

\begin{theorem}[CNT entropy of the shift equals mean entropy]\label{thm:cnt_shift}
Let $\phi$ be a translation-invariant state on $A=\bigotimes_{j\in\ZZ}\Md$.
If $\phi$ satisfies the \emph{asymptotic abelianness condition}: for each $\eps>0$
there exists $p\in\NN$ such that
\begin{equation}\label{eq:asymp_abelian}
  \bigl|\phi(Q_1 Q_2)-\phi(Q_1)\phi(Q_2)\bigr| \le \eps\norm{Q_1}\norm{Q_2}
\end{equation}
for all $Q_1\in A(]-\infty,-p])$ and $Q_2\in A([p,+\infty[)$,
then
\begin{equation}\label{eq:cnt_shift}
  h_\phi(\tau_1) = S(\phi) = \lim_{n\to\infty}\frac{1}{n}\,S(\phi|_{A([0,n-1])}).
\end{equation}
\end{theorem}

\begin{remark}
The asymptotic abelianness condition~\eqref{eq:asymp_abelian} is the \emph{clustering}
property.  It holds for:
\begin{itemize}
\item Gibbs states of short-range interactions at high temperature (Lemma A.2 of \cite{cnt87}).
\item Quasi-free (free-fermion/boson) states at any temperature (Lemma A.3 of \cite{cnt87}).
\item KMS states of quasi-local systems above the critical temperature.
\end{itemize}
For general translation-invariant states, one has $h_\phi(\tau_1)\le S(\phi)$
(upper bound, proved below), but the lower bound requires clustering.
\end{remark}

\subsubsection{Upper Bound: $h_\phi(\tau_1)\le S(\phi)$}

\begin{lemma}[Upper bound]\label{lem:upper_bound}
For any translation-invariant state $\phi$ on $\bigotimes_\ZZ\Md$:
\begin{equation}\label{eq:upper_bound}
  h_\phi(\tau_1) \le S(\phi).
\end{equation}
\end{lemma}

\begin{proof}
For any CPU map $\gamma:\Mn\to A$ and $k$ large:
\begin{align*}
  \Hphi(\gamma,\tau_1\circ\gamma,\ldots,\tau_1^{k-1}\circ\gamma)
  &\le \Hphi(A([0,k+\mathrm{supp}]))
  \le S(\phi|_{A([0,k+R])})
\end{align*}
where $R$ is the ``support radius'' of $\gamma$ (the smallest $R$ such that $\gamma$
maps into $A([-R,R])$).  By translation invariance and subadditivity of entropy,
$\frac{1}{k}S(\phi|_{A([0,k+R])})\to S(\phi)$ as $k\to\infty$ (since
$S(\phi|_{A([0,n])})/n\to S(\phi)$).  Hence
$h_{\phi,\tau_1}(\gamma)\le S(\phi)$, and since $\gamma$ is arbitrary,
$h_\phi(\tau_1)\le S(\phi)$.
\end{proof}

\subsubsection{Lower Bound under Clustering: $h_\phi(\tau_1)\ge S(\phi)$}

The lower bound is the hard direction and uses the algebraic structure of
abelian subalgebras of the centralizer.  We follow \cite{cnt87} Sections VIII--IX.

\begin{lemma}[Corollary VIII.8 of CNT]\label{lem:abelian_subalg}
Let $N_1,\ldots,N_k$ be finite-dimensional subalgebras of $\vN$ and assume that
they contain pairwise commuting abelian subalgebras $A_j\subset N_j\cap\vN_\phi$
(centralizer of $\phi$) such that $A:=\bigvee_j A_j$ is maximal abelian in
$N:=\bigvee_j N_j$.  Then:
\begin{equation}\label{eq:abelian_subalg}
  \Hphi(N_1,\ldots,N_k) = S(\phi|_N).
\end{equation}
\end{lemma}

\begin{proof}
Since each $A_j\subset\vN_\phi$, there exist state-preserving conditional expectations
$E_j:\vN\to A_j$.  The abelian model $(A,P_A,\phi|_A)$ with $P_A$ being the conditional
expectation from $\vN$ onto $A$ satisfies $s_\mu(P_j\circ\iota_{A_j})=0$ for all $j$
(because the $\hat\phi_i^j$ are pure on $A_j$ which is abelian and $\phi$ restricted
to a maximal abelian subalgebra of the centralizer is a measure).
Hence:
\begin{align*}
  \Hphi(N_1,\ldots,N_k)
  &\ge \Hphi(A_1,\ldots,A_k)\\
  &= S(\mu|_A) - \sum_j s_\mu(E_j\circ\iota_{A_j})\\
  &= S(\phi|_A) - 0 = S(\phi|_N).
\end{align*}
Combined with the upper bound $\Hphi(N_1,\ldots,N_k)\le S(\phi|_N)$
(Lemma VIII.1 of \cite{cnt87}), equality follows.
\end{proof}

\begin{proof}[Proof of Theorem~\ref{thm:cnt_shift}]
\textbf{Step 1: Choose a good subalgebra $N$.}
Let $I=[0,n-1]$ and $N_k=\tau_1^k(A(I))=A([k,k+n-1])$ for $k=0,1,\ldots,m-1$.
These are the iterated images of the single-site subalgebra $A(I)$ under the shift.

\textbf{Step 2: Find the abelian subsubalgebras.}
Let $C_0\subset A(I)$ be a maximal abelian subalgebra of the centralizer of $\phi|_{A(I)}$.
Then $C_k=\tau_1^k(C_0)\subset N_k$ is maximal abelian in $N_k$ and in the centralizer,
and the $C_k$ pairwise commute (they live in different spatial regions $[k,k+n-1]$,
hence commute when $|k-k'|\ge n$).

\textbf{Step 3: Apply Lemma~\ref{lem:abelian_subalg}.}
By Lemma~\ref{lem:abelian_subalg},
\begin{equation}\label{eq:step3}
  \Hphi(N_0,N_1,\ldots,N_{m-1})
  = S(\phi|_{A([0,n+m-2])}).
\end{equation}
This requires $m\ge n+2N$ (where $N$ comes from the clustering condition) so that
the commutator of operators in separated regions is small.

\textbf{Step 4: Estimate via Proposition IX.1.}
Using the clustering condition~\eqref{eq:asymp_abelian}, \cite{cnt87} Proposition IX
(specifically, Lemmas IX.2--IX.4) shows that for $m\ge n+2N$:
\begin{equation}\label{eq:step4}
  \frac{1}{m}h_{\phi,\tau_1}(A(I))
  \ge \frac{1}{m} S(\phi|_{A([0,n+m-2])}) - O\!\left(\frac{1}{m},\eps,\frac{N}{m}\right).
\end{equation}
As $m\to\infty$: $\frac{1}{m}S(\phi|_{A([0,n+m])})\to S(\phi)$ (translation invariance),
and the error terms vanish.

\textbf{Step 5: Take the sup.}
Taking $n\to\infty$ and optimizing over all CPU maps $\gamma$
(which includes $A(I)$-inclusions for all $I$):
\begin{equation}
  h_\phi(\tau_1) \ge \sup_n \lim_{m\to\infty} \frac{1}{m}S(\phi|_{A([0,n+m-2])}) = S(\phi).
\end{equation}
Combined with Lemma~\ref{lem:upper_bound}: $h_\phi(\tau_1)=S(\phi)$.
\end{proof}

%% ============================================================
\section{Explicit Finite-Dimensional Examples}
\label{sec:examples}
%% ============================================================

\subsection{Classical Probability: Bernoulli Shift}

Before the quantum case, let us verify the theorem in the classical setting.

\begin{example}[Bernoulli shift]\label{ex:bernoulli}
Let $\mathcal{X}=\{0,1\}^\ZZ$ with the Bernoulli measure $\mu=(\frac{1}{2},\frac{1}{2})^\ZZ$
(fair coin).  The shift $T$ acts by $(Tx)_n=x_{n+1}$.

The mean entropy is $S(\mu)=\log 2$ (the Shannon entropy of one coin flip).
The corresponding C*-algebra is $A=\bigotimes_\ZZ\CC^2$ (abelian, since $\CC^2$ is abelian
in the commutative sense), and the partition $\mathcal{P}=\{[x_0=0],[x_0=1]\}$ generates
$A(\{0\})$.

\medskip\noindent\textbf{Direct KS computation:}
\begin{align*}
  H_\mu(\mathcal{P}\vee T^{-1}\mathcal{P}\vee\cdots\vee T^{-(n-1)}\mathcal{P})
  &= H_\mu(\text{cylinder sets on } [0,n-1])\\
  &= n\cdot H(\tfrac{1}{2},\tfrac{1}{2}) = n\log 2.
\end{align*}
Hence $h_{\mathrm{KS}}(T)=\log 2=S(\mu)$, confirming the theorem.
\end{example}

\subsection{Qubit: $d=2$, Single-Site Computation}

\begin{example}[Qubit chain entropy computation]\label{ex:qubit}
Let $d=2$, $A=\bigotimes_\ZZ M_2(\CC)$.
Consider the translation-invariant \emph{product state}
$\phi=\bigotimes_j\omega_j$ where $\omega=\Tr(\rho\cdot)$ with
\begin{equation}\label{eq:qubit_state}
  \rho = \begin{pmatrix} p & 0 \\ 0 & 1-p \end{pmatrix}, \qquad 0<p<1.
\end{equation}
This is a KMS (Gibbs) state for the single-site Hamiltonian $H=-(\log p)\sigma_z$
at inverse temperature $\beta=1$.

\textbf{Mean entropy:}
Since $\phi$ is a product state, $S(\phi|_{A([0,n-1])})=n\cdot S(\omega)$,
where
\begin{equation}\label{eq:qubit_entropy}
  S(\omega) = -p\log p - (1-p)\log(1-p) =: h_2(p).
\end{equation}
Hence $S(\phi)=h_2(p)$.

\textbf{CNT entropy:}
By Theorem~\ref{thm:cnt_shift} (the product state has exponential clustering),
$h_\phi(\tau_1)=S(\phi)=h_2(p)$.

\textbf{Explicit $H_\phi$ computation for $n=1$:}
Take $N=A(\{0\})\cong M_2(\CC)$ (the local algebra at site 0).
We compute $\Hphi(N)$ directly.

Any decomposition $\phi=\mu_1\hat\phi_1+\mu_2\hat\phi_2$ with $\hat\phi_i$ states
on $A$ corresponds (by restriction to $N$) to $\omega=\mu_1\hat\omega_1+\mu_2\hat\omega_2$
on $M_2(\CC)$ with density matrices $\sigma_i=\hat\omega_i(\cdot)$.

The one-subalgebra entropy functional is:
\begin{align}
  \Hphi(N)
  &= \sup_{\omega=\sum_i\mu_i\hat\omega_i}
     \left\{H(\mu) - \sum_i\mu_i S(\hat\omega_i)\right\}\nonumber\\
  &= \sup_{\omega=\sum_i\mu_i\sigma_i}
     \left\{-\sum_i\mu_i\log\mu_i - \sum_i\mu_i S(\sigma_i)\right\}.\label{eq:qubit_Hphi}
\end{align}
Since $\omega=\rho=\mathrm{diag}(p,1-p)$ and $S(\rho)=h_2(p)$:

\medskip\noindent\textbf{Upper bound:} $\Hphi(N)\le S(\omega)=h_2(p)$ (by Prop.~\ref{prop:H1_props}(a)).

\medskip\noindent\textbf{Lower bound — take the projector decomposition:}
Use the spectral decomposition $\rho=p|0\rangle\langle 0|+(1-p)|1\rangle\langle 1|$,
i.e., $\mu_1=p$, $\sigma_1=|0\rangle\langle 0|$, $\mu_2=1-p$, $\sigma_2=|1\rangle\langle 1|$.
Then $S(\sigma_1)=S(\sigma_2)=0$ (pure states), so:
\begin{align}
  H(\mu)-\sum_i\mu_i S(\sigma_i)
  &= (-p\log p-(1-p)\log(1-p)) - 0 = h_2(p).
\end{align}
Hence $\Hphi(N)\ge h_2(p)$, confirming $\Hphi(N)=h_2(p)=S(\omega)$.

\begin{center}
\begin{tabular}{@{}cccc@{}}
\toprule
$p$ & $S(\omega)=h_2(p)$ & $\Hphi(N)$ & Ratio \\
\midrule
0.1 & 0.3251 & 0.3251 & 1.000 \\
0.2 & 0.5004 & 0.5004 & 1.000 \\
0.5 & 0.6931 & 0.6931 & 1.000 \\
0.8 & 0.5004 & 0.5004 & 1.000 \\
0.9 & 0.3251 & 0.3251 & 1.000 \\
\bottomrule
\end{tabular}
\end{center}

\vspace{4pt}\noindent All values computed as $h_2(p)=-p\ln p-(1-p)\ln(1-p)$ (natural log).
The equality $\Hphi(N)=S(\omega)$ holds because $\omega|_N$ is a faithful state on
$M_2(\CC)$ and $N\subset\mathcal{M}_\phi$ (the density matrix $\rho$ commutes with itself,
so the centralizer contains all of $M_2(\CC)$).

\end{example}

\subsection{Qutrit: $d=3$, Two-Step Orbit}

\begin{example}[Qutrit chain, two-step computation]\label{ex:qutrit}
Let $d=3$, $A=\bigotimes_\ZZ M_3(\CC)$.
Consider the tracial product state $\phi=\bigotimes_j\mathrm{tr}_j$ where
$\mathrm{tr}(a)=\frac{1}{3}\Tr(a)$.  Then $S(\mathrm{tr})=\log 3$.

\textbf{Two-step orbit entropy:}
Let $N=A(\{0\})\cong M_3(\CC)$ and $N'=A(\{0,1\})\cong M_9(\CC)$.
\begin{equation}
  \Hphi(N, \tau_1(N))
  = \Hphi(A(\{0\}),A(\{1\}))
  = \Hphi(A(\{0,1\}))
  = S(\phi|_{A(\{0,1\})}).
\end{equation}
For the tracial product state: $\phi|_{A(\{0,1\})}=\mathrm{tr}_3\otimes\mathrm{tr}_3$,
so $S(\phi|_{A(\{0,1\})})=2\log 3$.

\textbf{Rate:}
\begin{equation}
  h_{\phi,\tau_1}(N) = \lim_{k\to\infty}\frac{1}{k}\Hphi(N,\tau_1 N,\ldots,\tau_1^{k-1} N)
  = \lim_{k\to\infty}\frac{1}{k}\cdot k\log 3 = \log 3.
\end{equation}

\textbf{Explicit matrix:}
The density matrix for $\phi|_{A([0,1])}=\frac{1}{9}I_9$ (maximally mixed on $M_9(\CC)$),
with entropy $S = \log 9 = 2\log 3$.  The two-step orbit gives:
\begin{equation}\label{eq:qutrit_2step}
  \rho^{(2)} = \frac{1}{9}I_9
  = \frac{1}{9}\begin{pmatrix}
    1 & & & \\
    & 1 & & \\
    & & \ddots & \\
    & & & 1
  \end{pmatrix}_{9\times 9}.
\end{equation}
All 9 eigenvalues equal $\frac{1}{9}$, giving $S(\rho^{(2)})=\log 9=2\log 3$.

\begin{center}
\begin{tabular}{@{}cccc@{}}
\toprule
$n$ (steps) & $S(\phi|_{A([0,n-1])})$ & $(1/n)S$ & CNT rate \\
\midrule
1 & $\log 3 \approx 1.099$ & 1.099 & \multirow{4}{*}{$\log 3$} \\
2 & $2\log 3 \approx 2.197$ & 1.099 & \\
3 & $3\log 3 \approx 3.296$ & 1.099 & \\
4 & $4\log 3 \approx 4.394$ & 1.099 & \\
\bottomrule
\end{tabular}
\end{center}

The mean entropy $S(\phi)=\log 3$ is achieved exactly at every step because $\phi$
is a product state (no correlations between sites).
\end{example}

\subsection{Non-Trivial Example: XXZ Chain at Finite Temperature}

\begin{example}[XXZ chain, qualitative]\label{ex:xxz}
Consider the XXZ Hamiltonian on $\ZZ$:
\begin{equation}\label{eq:xxz}
  H_{\mathrm{XXZ}} = \sum_{j\in\ZZ}\bigl[
    \sigma_j^x\sigma_{j+1}^x + \sigma_j^y\sigma_{j+1}^y
    + \Delta\,\sigma_j^z\sigma_{j+1}^z
  \bigr],
\end{equation}
with anisotropy $\Delta\ge 0$.  The Gibbs state at inverse temperature $\beta$ satisfies
the clustering condition (for $\beta<\beta_c(\Delta)$, the critical inverse temperature).
By Theorem~\ref{thm:cnt_shift}:
\begin{equation}
  h_\phi(\tau_1) = S(\phi) = \lim_{n\to\infty}\frac{1}{n}\,S(\phi|_{[0,n-1]}),
\end{equation}
where $S(\phi)$ is the thermodynamic entropy density.  For the isotropic case ($\Delta=1$,
the XXX Heisenberg chain), $S(\phi)$ decreases monotonically from $\log 2$ (at $\beta=0$,
the tracial state) to 0 (at $\beta=\infty$, the ground state).

The CNT entropy thus interpolates between maximum scrambling ($h=\log 2$ for infinite temperature)
and zero entropy ($h=0$ for the ground state), tracking the thermodynamic entropy density.
\end{example}

%% ============================================================
\section{Comparison: CNT vs.\ AFL/OPU Entropy}
\label{sec:comparison}
%% ============================================================

The AFL (Alicki--Fannes--Lindblad) entropy studied in subfolder~1 and the CNT entropy
are both quantum generalizations of the KS entropy, but they differ fundamentally
in their construction and domain of applicability.

\begin{center}
\begin{tabular}{@{}lll@{}}
\toprule
\textbf{Property} & \textbf{AFL/OPU entropy} & \textbf{CNT entropy} \\
\midrule
Input & OPU $\{Z_k\}$, unitary $U$ & CPU map $\gamma:\Mn\to A$, $\thetaaut$ \\
System & Finite or infinite dim. & Infinite dim. (for $h>0$) \\
Entropy function & Von Neumann entropy $S_1$ & Relative entropy $S(\phi,\psi)$ \\
Finite-dim.\ value & $h_\mathrm{AFL}>0$ possible & $h_\mathrm{CNT}=0$ always \\
Shift on spin chain & $h_\mathrm{AFL}=S(\omega)+\log d$ & $h_\mathrm{CNT}=S(\phi)$ \\
KS reduction & Via projector OPU & Exact (abelian = classical) \\
Chaos detection & Yes (Rényi spread) & Yes (if $h_\mathrm{CNT}>0$, infinite system) \\
Open systems & Via Stinespring dilation & Via conditional expectations \\
\bottomrule
\end{tabular}
\end{center}

\begin{remark}[AFL excess over CNT]\label{rmk:afl_excess}
For the shift on a qubit chain with state $\phi$ (product state with local density $\rho$):
\begin{itemize}
  \item $h_\mathrm{AFL}(\tau_1) = S(\omega)+\log d = h_2(p)+\log 2$,
  \item $h_\mathrm{CNT}(\tau_1) = S(\phi) = h_2(p)$.
\end{itemize}
The difference is $\log d = \log 2$, which corresponds to the OPU ``background entropy''
(the contribution of the observable structure $\{P_i\}$ to the AFL entropy).
In subfolder~1 this was identified as the entropy of the OPU partition itself
(cf.\ Section~7 of subfolder~1's Output.tex, Theorem~7.1):
$h_\mathrm{AFL}(\tau_1,Z) = s(\omega)+\log d$ where $s(\omega)=S(\phi)$.
The CNT entropy captures only the dynamical contribution $s(\omega)$, discarding the
OPU structural entropy $\log d$.
\end{remark}

\begin{remark}[Why both are needed]\label{rmk:both_needed}
AFL entropy: operationally natural for finite measurements (POVM coarse-graining);
non-zero even in finite dimensions; directly computable for specific observables.

CNT entropy: conceptually ``purer'' generalization of KS; isomorphism invariant;
equals the thermodynamic entropy density for equilibrium states; correctly zero for
finite-dimensional systems.

Both satisfy the quantum Pesin inequality $h\le v_B\log d$
(as proved in subfolder~1, Section~39 for AFL; for CNT it follows from the same
Lieb-Robinson bound argument since the orbit subalgebras are contained in lightcones).
\end{remark}

%% ============================================================
\section{Summary and Open Directions}
\label{sec:summary}
%% ============================================================

We have established the following for CNT entropy:

\begin{enumerate}
\item \textbf{Construction.}  The CNT entropy $h_\phi(\thetaaut)$ is built from relative
  entropies of abelian models, avoiding the non-monotonicity problem of subalgebra entropy.
\item \textbf{Properties.}  It is conjugacy-invariant, affine in $\phi$, additive in
  powers of $\thetaaut$, zero for finite-dimensional systems, and reduces exactly to
  the classical KS entropy in the abelian case.
\item \textbf{Main theorem.}  For the shift on a quantum spin chain with a clustering
  invariant state, $h_\phi(\tau_1)=S(\phi)$ (the thermodynamic entropy density).
\item \textbf{Comparison.}  CNT and AFL agree on the dynamical information content
  up to the structural entropy $\log d$ of the OPU partition.
\end{enumerate}

\medskip\noindent\textbf{Open directions from this study:}
\begin{enumerate}
\item \emph{Non-tracial CNT entropy and type III factors.}
  The original paper \cite{cnt87} handles type III factors via the modular automorphism.
  For KMS states of interacting quantum fields (not quasi-free), the computation of
  $h_\phi(\tau_1)$ remains open.
\item \emph{CNT entropy for open quantum systems.}
  Extending CNT to completely positive semigroups (Lindblad dynamics) requires
  a time-dependent abelian model construction.  The Stinespring dilation approach
  from subfolder~1 may provide a bridge.
\item \emph{CNT--AFL equivalence in the thermodynamic limit.}
  A precise theorem equating $h_\mathrm{CNT}$ and the limit $\lim_{d\to\infty}(h_\mathrm{AFL}-\log d)$
  for sequences of spin chains with growing local dimension $d$ would unify the two approaches.
\item \emph{Quantum Pesin for CNT.}
  Can one prove $h_\phi(\thetaaut)=\sum_{\lambda_i>0}\lambda_i$ (sum of positive Lyapunov
  exponents) for infinite quantum spin chains in the semiclassical limit?
  The upper bound $h_\mathrm{CNT}\le v_B\log d$ follows from Lieb-Robinson;
  the lower bound is open for non-Gaussian states.
\end{enumerate}

%% ============================================================
\section{New Result: Exact AFL--CNT Gap for All Clustering States}
\label{sec:afl_cnt_gap}
%% ============================================================

The comparison in Section~\ref{sec:comparison} identified an excess of $\log d$
between $h_\mathrm{AFL}$ and $h_\mathrm{CNT}$ in the product-state case.
We now prove this gap is \emph{universal}: it holds for all translation-invariant
clustering states on a $d$-dimensional spin chain.

\begin{theorem}[AFL--CNT gap]\label{thm:afl_cnt_gap}
Let $\phi$ be a translation-invariant state on $A=\bigotimes_\ZZ\Md$ satisfying the
clustering condition~\eqref{eq:asymp_abelian}, and let $Z=\{P_0,\ldots,P_{d-1}\}$
be the projector OPU ($P_i=|i\rangle\langle i|$, single-site computational basis).
Then:
\begin{equation}\label{eq:afl_cnt_gap}
  h_\mathrm{AFL}(\tau_1, Z) - h_\mathrm{CNT}(\tau_1) = \log d,
\end{equation}
where $h_\mathrm{AFL}(\tau_1,Z)$ is the AFL dynamical entropy with OPU $Z$,
and $h_\mathrm{CNT}(\tau_1)$ is the CNT dynamical entropy.
\end{theorem}

\begin{proof}
We use two exact results:
\begin{enumerate}
\item[(A)] From subfolder 1, Theorem~7.1 (AFL shift formula):
  \begin{equation}\label{eq:afl_shift}
    h_\mathrm{AFL}(\tau_1, Z) = s(\omega) + \log d,
  \end{equation}
  where $s(\omega) = S(\phi) = \lim_{n\to\infty}\frac{1}{n}S(\phi|_{A([0,n-1])})$
  is the mean entropy density of $\phi$.
\item[(B)] From Theorem~\ref{thm:cnt_shift} (CNT shift formula, proved above):
  \begin{equation}\label{eq:cnt_shift_restate}
    h_\mathrm{CNT}(\tau_1) = S(\phi) = s(\omega).
  \end{equation}
\end{enumerate}
Subtracting (B) from (A):
\begin{equation}
  h_\mathrm{AFL}(\tau_1,Z) - h_\mathrm{CNT}(\tau_1)
  = \bigl[s(\omega)+\log d\bigr] - s(\omega) = \log d.
\end{equation}
\end{proof}

\begin{remark}[Physical interpretation]\label{rmk:gap_interpretation}
The gap $\log d$ has a clear physical meaning.
\begin{itemize}
\item The AFL OPU $Z=\{P_i\}$ produces outcomes with probabilities $p_i=\phi(P_i\cdot P_i)$
  on the local algebra.  The \emph{structural entropy} of the OPU itself --- the
  maximum classical entropy attainable by any state through the partition $Z$ ---
  is $\log d$ (achieved by the tracial state).
\item In the AFL rate formula $h_\mathrm{AFL}=s(\omega)+\log d$, the term $\log d$
  comes from the fact that the OPU introduces $d$ distinguishable cells at each step,
  contributing a classical factor $\log d$ per step to the Shannon entropy of the
  measurement sequence.
\item The CNT entropy removes this structural contribution: it measures the entropy
  production rate \emph{of the state alone}, independent of any particular observable
  or partition.  Hence $h_\mathrm{CNT}=s(\omega)$.
\item The gap $\log d$ is \emph{independent} of the state $\phi$, of correlations,
  and of temperature.  It depends only on the local Hilbert space dimension $d$.
\end{itemize}
\end{remark}

\begin{remark}[Contrast with single-site entropy]\label{rmk:s1_vs_s_phi}
For a correlated state, the single-site entropy $S_1=S(\phi|_{A(\{0\})})$ satisfies
$S_1 \ge s(\omega) = S(\phi)$, with equality only for product states.
The excess $S_1-s(\omega)\ge 0$ equals the \emph{information in one site about the rest}:
\begin{equation}\label{eq:mutual_excess}
  S_1 - s(\omega)
  = \lim_{n\to\infty}\frac{1}{n}\,I(A(\{0\}) : A([1,n-1]))_\phi,
\end{equation}
where $I(A:B)_\phi = S(\phi|_A)+S(\phi|_B)-S(\phi|_{AB})$ is the mutual information.
This quantity grows with correlations and can be significant for strongly-interacting
chains (e.g., $S_1-s(\omega)\approx 0.55$ for the XXX chain at $\beta=5$,
where $s(\omega)\approx 0.14$ but $S_1=\log 2\approx 0.69$).
\end{remark}

\subsection{Numerical Verification}

\begin{center}
\begin{tabular}{@{}ccccc@{}}
\toprule
$\beta$ (XXX chain) & $s(\omega)=S_L/L$ & $h_\mathrm{AFL}=s+\log 2$ & $h_\mathrm{CNT}=s$ & Gap \\
\midrule
0.0 (tracial) & 0.693147 & 1.386294 & 0.693147 & 0.693147 \\
0.5           & 0.670195 & 1.363342 & 0.670195 & 0.693147 \\
1.0           & 0.601762 & 1.294909 & 0.601762 & 0.693147 \\
2.0           & 0.415608 & 1.108755 & 0.415608 & 0.693147 \\
5.0           & 0.144719 & 0.837866 & 0.144719 & 0.693147 \\
\bottomrule
\end{tabular}
\end{center}

\vspace{4pt}\noindent
The gap is $\log 2 = 0.693147$ in all cases (L=8 open chain, $s(\omega)$ estimated as $S_8/8$).
This confirms Theorem~\ref{thm:afl_cnt_gap} numerically.

\begin{corollary}[Gap for any $d$]\label{cor:gap_d}
For a $d$-dimensional spin chain ($d\ge 2$) with projector OPU $Z=\{|i\rangle\langle i|\}$:
\begin{equation}\label{eq:gap_d}
  h_\mathrm{AFL}(\tau_1,Z) - h_\mathrm{CNT}(\tau_1) = \log d,
\end{equation}
for all translation-invariant clustering states.  The gap equals $\log 2$ for qubit,
$\log 3$ for qutrit, $\log 4=2\log 2$ for a 4-level system, and so on.
\end{corollary}

\begin{remark}[Optimal OPU]\label{rmk:optimal_opu}
The gap $\log d$ in~\eqref{eq:afl_cnt_gap} is achieved by the projector OPU.
Other OPU choices give a different (generally larger) AFL entropy for the same dynamics.
In subfolder 1 (Section~55) it was shown that the projector OPU maximises the
integrable--chaotic gap, while the SIC-POVM and matrix-unit OPU give smaller gaps.
The CNT entropy is OPU-independent by construction.
\end{remark}

%% ============================================================
\section{CNT Entropy for Time Evolution and Modular Automorphisms}
\label{sec:time_evolution}
%% ============================================================

\subsection{Time Evolution Automorphism}

Let $(\vN,\alpha_t,\phi)$ be a quantum lattice system where
$\alpha_t = \mathrm{Ad}(e^{-iHt})$ is the one-parameter group of time evolution
automorphisms and $\phi$ is the $\alpha_t$-invariant KMS state at inverse temperature $\beta$.

\begin{theorem}[Finite-dim vanishing for time evolution]\label{thm:finite_time}
For any finite quantum lattice system on $L$ sites (i.e., $\vN=M_{2^L}(\CC)$),
$h_\phi(\alpha_t)=0$ for all $t\in\RR$.
\end{theorem}

\begin{proof}
For any CPU map $\gamma:\Mn\to M_{2^L}(\CC)$ and any $n\ge 1$:
\begin{equation}
  \Hphi(\gamma,\alpha_t\circ\gamma,\ldots,\alpha_t^{n-1}\circ\gamma)
  \le \log\bigl(\dim M_{2^L}(\CC)\bigr)^2
  = 2L\log 2.
\end{equation}
Hence $h_{\phi,\alpha_t}(\gamma) = \lim_{n\to\infty}\frac{1}{n}\cdot O(1)=0$.
\end{proof}

\begin{remark}[Pre-saturation slope]
Although the CNT rate is zero for any finite $L$, the finite-$n$ entropy
$S(\rho[Z^{(n)}])$ grows at a rate that reflects the true thermodynamic-limit
entropy production.  For the XXX chain ($L=5$, $\beta=1$, $\delta t=0.5$):
\begin{center}
\begin{tabular}{@{}cccr@{}}
\toprule
$n$ & $S(\rho[Z^{(n)}])$ & $S/n$ & Slope $S_n-S_{n-1}$ \\
\midrule
1 & 0.6931 & 0.6931 & 0.6931 \\
2 & 0.8576 & 0.4288 & 0.1644 \\
3 & 0.9736 & 0.3245 & 0.1160 \\
4 & 1.0790 & 0.2697 & 0.1054 \\
5 & 1.1815 & 0.2363 & 0.1025 \\
6 & 1.2829 & 0.2138 & 0.1014 \\
7 & 1.3837 & 0.1977 & 0.1008 \\
\bottomrule
\end{tabular}
\end{center}
The slope quickly approaches $\approx 0.101$ per step $= 0.202$ per unit time.
This estimates the thermodynamic-limit CNT entropy rate:
$h_\phi(\alpha_t)|_{L=\infty} \approx 0.20 < v_\mathrm{LR}\log 2\approx 1.39$.
The $n=1$ contribution ($= \log 2$, the single-site entropy) is the OPU structural term
analogous to the AFL background.
\end{remark}

\subsection{Lieb-Robinson Upper Bound}

\begin{theorem}[LR-Pesin bound for time evolution]\label{thm:lr_pesin_time}
Let $A=\bigotimes_\ZZ\Md$ with nearest-neighbour interaction of strength $J$.
For the time evolution $\alpha_t$ with Lieb-Robinson velocity $v_\mathrm{LR}=2J$:
\begin{equation}\label{eq:lr_pesin_time}
  h_\phi(\alpha_t) \le v_\mathrm{LR}\,\log d\,|t|
  = 2J\log d\,|t|.
\end{equation}
\end{theorem}

\begin{proof}
Under $\alpha_t$, a single-site operator $a\in A(\{0\})$ evolves to
$\alpha_t(a)\in A([-v_\mathrm{LR}|t|-\delta, v_\mathrm{LR}|t|+\delta])$
up to exponentially small tails (Lieb-Robinson theorem \cite{lieb72}),
where $\delta=O(\log L)$.
The $n$-step orbit
$\{\gamma,\alpha_t\circ\gamma,\ldots,\alpha_{(n-1)t}\circ\gamma\}$
spans sites $0$ through $(n-1)\,v_\mathrm{LR}|t|$.
Each site contributes at most $\log d$ to the entropy, so:
\begin{equation}
  \Hphi(\gamma,\ldots,\alpha_{(n-1)t}\circ\gamma)
  \le n\,v_\mathrm{LR}|t|\,\log d + O(\log L).
\end{equation}
Dividing by $n$ and taking $n\to\infty$ gives~\eqref{eq:lr_pesin_time}.
\end{proof}

\begin{center}
\begin{tabular}{@{}ccc@{}}
\toprule
$t$ & $v_\mathrm{LR}\log(2)\cdot t$ & Upper bound for $h_\phi(\alpha_t)$ \\
\midrule
0.5 & 0.693 & $h_\phi(\alpha_{0.5})\le 0.693$ \\
1.0 & 1.386 & $h_\phi(\alpha_1)\le 1.386$ \\
2.0 & 2.773 & $h_\phi(\alpha_2)\le 2.773$ \\
\bottomrule
\end{tabular}
\end{center}

\subsection{Modular Automorphism and Tomita-Takesaki Theory}

\begin{definition}[Modular automorphism]\label{def:modular}
Let $\phi$ be a faithful normal state on a von Neumann algebra $\vN$.
The \emph{modular automorphism group} $\sigma^\phi_s:\vN\to\vN$ ($s\in\RR$)
is the one-parameter group defined by the Tomita-Takesaki theorem:
$\sigma^\phi_s(a) = \Delta^\phi_{is}\,a\,\Delta^\phi_{-is}$,
where $\Delta^\phi = S^\phi{S^\phi}^*$ is the modular operator and
$S^\phi$ is the Tomita operator: $S^\phi(a\Omega)=a^*\Omega$ for $\Omega$ the
cyclic and separating vector representing $\phi$ in the GNS construction.
\end{definition}

\begin{proposition}[KMS condition and modular flow]\label{prop:kms_modular}
If $\phi$ is a KMS state at inverse temperature $\beta$ with respect to $\alpha_t$
(i.e., $\phi(\alpha_t(a)b)=\phi(b\alpha_{t+i\beta}(a))$ for analytic $a$), then:
\begin{equation}\label{eq:kms_modular}
  \sigma^\phi_s(a) = \alpha_{-i\beta s}(a)
  \qquad\text{(on the local algebra, for quasi-free dynamics).}
\end{equation}
In particular, the modular flow is \emph{imaginary-time evolution}:
$\sigma^\phi_s = \alpha_{-i\beta s}$ corresponds to evolving by imaginary time $-i\beta s$.
For the XXX chain at $\beta=1$: $\sigma^\phi_s(a) = e^{-sH}ae^{sH}$ (similarity transform).
\end{proposition}

\begin{proof}
The KMS condition with $\beta=1$ reads: $\phi(a\alpha_i(b))=\phi(ba)$ for all $a,b$.
This is the Kubo-Martin-Schwinger condition, and it identifies $\alpha_i$ (the analytic continuation
of time evolution to imaginary time $i$) with the modular automorphism $\sigma^\phi_1$.
By the one-parameter group property, $\sigma^\phi_s=\alpha_{is}$.
For $\phi$ being the Gibbs state at $\beta$: $\sigma^\phi_s=\alpha_{-i\beta s}$.
\end{proof}

\begin{remark}[Spreading under imaginary-time evolution]\label{rmk:spread_imag}
Numerical computation for the XXX chain ($L=4$, $\beta=1$) shows that the modular
operator $\sigma_s^\phi(P_0)=e^{-sH}P_0 e^{+sH}$ spreads across the chain:
\begin{center}
\begin{tabular}{@{}ccccc@{}}
\toprule
$s$ & Site 0 deviation & Site 1 & Site 2 & Site 3 \\
\midrule
0.0 & 0.707 & 0.000 & 0.000 & 0.000 \\
0.5 & 0.753 & 0.046 & 0.001 & 0.000 \\
1.0 & 0.908 & 0.213 & 0.013 & 0.000 \\
2.0 & 1.886 & 1.469 & 0.319 & 0.028 \\
\bottomrule
\end{tabular}
\end{center}
(Deviation = $\norm{\rho_{\mathrm{site}}-I/2}_F$; large = concentrated, small = spread.)
The operator spreads from site 0 outward as $s$ increases, consistent with a non-zero
thermodynamic-limit CNT entropy $h_\phi(\sigma^\phi_1)>0$.
\end{remark}

\begin{theorem}[CNT entropy of modular automorphism]\label{thm:cnt_modular}
For the modular automorphism $\sigma^\phi_s$ of a hyperfinite von Neumann algebra $\vN$
with faithful state $\phi$:
\begin{enumerate}
\item[(a)] $h_\phi(\sigma^\phi_s)=|s|\cdot h_\phi(\sigma^\phi_1)$ (additivity in $s$).
\item[(b)] $h_\phi(\sigma^\phi_s)\ge 0$, with equality iff $\phi=\tau$ (trace).
\item[(c)] For the KMS state of a 1D lattice at $\beta$:
  $h_\phi(\sigma^\phi_1)=h_\phi(\alpha_{-i\beta})=h_\phi(\alpha_\beta)$,
  the CNT entropy of the imaginary-time automorphism.
\end{enumerate}
\end{theorem}

\begin{proof}
(a) Follows from the additivity property (ii) of Theorem~\ref{thm:cnt_props}:
$h_\phi((\sigma^\phi_1)^n)=n\cdot h_\phi(\sigma^\phi_1)$, so $h_\phi(\sigma^\phi_s)=|s|\cdot h_\phi(\sigma^\phi_1)$.

(b) For tracial $\phi$: $\sigma^\phi_s=\mathrm{id}$ (since $\Delta^\phi=I$ for traces),
hence $h_\phi(\sigma^\phi_s)=h_\phi(\mathrm{id})=0$.  Non-zero for non-tracial states
since the modular operator $\Delta^\phi\ne I$ then acts non-trivially.

(c) Follows from Proposition~\ref{prop:kms_modular}.
\end{proof}

\subsection{Summary: CNT Entropy for Different Automorphisms}

\begin{center}
\begin{tabular}{@{}lcc@{}}
\toprule
\textbf{Automorphism} & \textbf{Finite-$L$ value} & \textbf{Thermodynamic limit} \\
\midrule
Shift $\tau_1$ & 0 (always) & $h_\phi(\tau_1)=s(\phi)$ \\
Time evolution $\alpha_t$ & 0 (always) & $h_\phi(\alpha_t)\le v_\mathrm{LR}\log d\,|t|$ \\
Modular flow $\sigma^\phi_s$ & 0 (always) & $h_\phi(\sigma^\phi_s)=|s|\cdot h_\phi(\sigma^\phi_1)$ \\
\bottomrule
\end{tabular}
\end{center}

The shift automorphism is unique in having an exact formula ($h=s(\phi)$) for its CNT entropy.
For real-time evolution, the Lieb-Robinson bound gives the correct scaling but not the exact value.
For the modular automorphism, the CNT entropy is non-zero for non-tracial states and is
proportional to $|s|$ by the group property.

%% ============================================================
\section{CNT Orbit Entropy as a Chaos Diagnostic}
\label{sec:chaos_diagnostic}
%% ============================================================

\subsection{Why $h_\mathrm{CNT}(\tau_1)$ is Not a Chaos Diagnostic}

\begin{proposition}[CNT shift entropy is chaos-blind]\label{prop:shift_chaos_blind}
The CNT dynamical entropy $h_\phi(\tau_1)=S(\phi)$ of the shift automorphism
depends only on the thermodynamic entropy density of the state $\phi$, and not
on any dynamical property (integrability vs.\ chaos) of the lattice Hamiltonian.
\end{proposition}

\begin{proof}
By Theorem~\ref{thm:cnt_shift}, $h_\phi(\tau_1)=S(\phi)$ for any translation-invariant
clustering state, integrable or not.  For the XXX Heisenberg chain and the XX free fermion
chain at the same inverse temperature $\beta$, both chains have different dynamics but
give $S(\phi)\approx 0.60$ at $\beta=1$ (L=8 numerical estimates):
\begin{center}
\begin{tabular}{@{}lccc@{}}
\toprule
System & $S_L/L$ ($L=4$) & $S_L/L$ ($L=6$) & $S_L/L$ ($L=8$) \\
\midrule
XXX Heisenberg (integrable) & 0.6128 & 0.6054 & 0.6018 \\
XX free fermion (integrable) & 0.6495 & 0.6450 & 0.6427 \\
Infinite temperature (tracial) & 0.6931 & 0.6931 & 0.6931 \\
\bottomrule
\end{tabular}
\end{center}
The shift entropy depends only on temperature, not on the Hamiltonian.
\end{proof}

\subsection{The OPU Triviality Lemma for Diagonal Floquet Circuits}

\begin{lemma}[Diagonal OPU triviality]\label{lem:diagonal_trivial}
Let $U$ be a unitary operator that is \emph{diagonal} in the computational basis
$\{|x\rangle\}$ (i.e., $U|x\rangle=e^{i\varphi(x)}|x\rangle$), and let
$Z=\{P_0,\ldots,P_{d-1}\}$ be a projector OPU in the same computational basis
(i.e., $P_i=|i\rangle\langle i|\otimes I_\mathrm{rest}$ at some site).
Then $\alpha_k(P_i) = P_i$ for all $k\ge 0$, and the time-evolution orbit
$\{Z, \alpha_1(Z), \ldots, \alpha_{n-1}(Z)\}$ is degenerate (all copies equal),
giving $S(\rho[Z^{(n)}]) = \log d$ for all $n$.
\end{lemma}

\begin{proof}
Since $U$ is diagonal: $U|x\rangle = e^{i\varphi(x)}|x\rangle$.
For $P_i = |i\rangle\langle i|\otimes I_\mathrm{rest}$:
\begin{align*}
  \alpha_k(P_i)|x_0,x_\mathrm{rest}\rangle
  &= U^k P_i (U^{-k}|x_0,x_\mathrm{rest}\rangle)\\
  &= U^k P_i\bigl(e^{-ik\varphi(x)}|x\rangle\bigr)
  = e^{-ik\varphi(x)} U^k \delta_{x_0,i}|x\rangle\\
  &= e^{-ik\varphi(x)} e^{ik\varphi(x)} \delta_{x_0,i}|x\rangle
  = \delta_{x_0,i}|x\rangle = P_i|x\rangle.
\end{align*}
Hence $\alpha_k(P_i)=P_i$ for all $k$.  The orbit is constant, and
$\rho[Z^{(n)}]_{(i),(j)} = \phi(P_{j_0}\cdots P_{j_{n-1}}P_{i_0}\cdots P_{i_{n-1}})$.
For orthogonal projectors: $P_{j_{k}}P_{i_k} = \delta_{i_k,j_k}P_{i_k}$ and
$P_{i_0}\cdots P_{i_{n-1}} = \prod_k \delta_{i_{k-1},i_k} P_{i_0}$ (all must agree).
Therefore $\rho[Z^{(n)}]$ is the $d\times d$ identity matrix $\frac{1}{d}I_d$
(for any translation-invariant state with $\phi(P_i)=1/d$), giving $S=\log d$.
\end{proof}

\begin{corollary}[Kicked Ising sigma-z OPU is trivial]\label{cor:ki_trivial}
The kicked Ising Floquet unitary
$U_\mathrm{KI}=e^{-ig\sum_k\sigma_z^k}e^{-iJ\sum_k\sigma_z^k\sigma_z^{k+1}}$
is diagonal in the computational (sigma-z) basis.  Hence sigma-z basis projectors
$\{|0\rangle\langle 0|,|1\rangle\langle 1|\}$ give $S(\rho[Z^{(n)}])=\log 2$ for all $n$,
providing no dynamical information.
\end{corollary}

\subsection{Initial Slope as the Correct Chaos Diagnostic}

To diagnose chaos, one must use an OPU whose elements are \emph{not} in the
eigenbasis of $U$.  We use the sigma-$x$ basis:
\begin{equation}\label{eq:sx_opu}
  P_\pm = |{\pm}\rangle\langle{\pm}| = \frac{1}{2}(I\pm\sigma_x),
  \qquad |{\pm}\rangle = \frac{|0\rangle\pm|1\rangle}{\sqrt{2}}.
\end{equation}

\begin{definition}[Initial slope]\label{def:initial_slope}
For a CPU map $\gamma:\Mn\to\vN$ and automorphism $\alpha$, the \emph{initial slope}
is the one-step entropy increase:
\begin{equation}\label{eq:initial_slope}
  \Delta_1(\alpha,\gamma,\phi)
  := \Hphi(\gamma,\alpha\circ\gamma) - \Hphi(\gamma)
  = S(\rho[Z^{(2)}]) - S(\rho[Z^{(1)}]).
\end{equation}
This equals the entropy added by one step of the dynamics, measuring how much new
information the orbit generates at each step.
\end{definition}

\begin{proposition}[Initial slope bounds chaos]\label{prop:initial_slope}
\begin{enumerate}
\item[(a)] $\Delta_1\in[0,\log d]$.  $\Delta_1=0$ for constant orbits (diagonal $U$,
  computational basis OPU).  $\Delta_1=\log d$ iff the orbit doubles the rank at step 2.
\item[(b)] For dual-unitary $U$ and non-trivially-placed OPU:
  $\Delta_1=\log d$ (maximum).
\item[(c)] For integrable systems: $\Delta_1<\log d$ (sub-maximal initial slope).
\end{enumerate}
\end{proposition}

\subsection{Numerical Results: XXX vs.\ Kicked Ising (sigma-x OPU)}

\begin{center}
\begin{tabular}{@{}ccccc@{}}
\toprule
$n$ & \multicolumn{2}{c}{$S(\rho[Z^{(n)}])$} & \multicolumn{2}{c}{Slope} \\
\cmidrule(lr){2-3}\cmidrule(lr){4-5}
 & XXX (integrable) & KI dual-unitary & XXX & KI \\
\midrule
1 & 0.693 & 0.693 & 0.693 & 0.693 \\
2 & 1.100 & 1.386 & 0.407 & \textbf{0.693} \\
3 & 1.408 & 1.386 & 0.308 & 0.000 \\
4 & 1.700 & 1.386 & 0.292 & 0.000 \\
5 & 1.987 & 1.386 & 0.287 & 0.000 \\
6 & 2.271 & 1.386 & 0.285 & 0.000 \\
7 & 2.551 & 1.386 & 0.280 & 0.000 \\
\bottomrule
\end{tabular}
\end{center}

\vspace{4pt}\noindent
(L=5, $\beta=1$ for XXX, tracial state for KI. Sigma-$x$ OPU at site 0.)

Key observations:
\begin{enumerate}
\item \textbf{Initial slope}: KI achieves $\Delta_1=\log 2$ (maximum), XXX achieves
  $\Delta_1=0.407<\log 2$ (sub-maximal).  This is the chaos diagnostic.
\item \textbf{Orbit saturation}: KI saturates at $S=\log 4=2\log 2$ after $n=2$.
  This reflects fast scrambling: the sigma-$x$ orbit reaches full rank~4 in 2 steps.
  XXX keeps growing (rank $\gg 4$ for $n\ge 3$), reflecting slow operator spreading.
\item \textbf{Rank growth}: KI reaches maximum rank $d^2=4$ in 2 steps (dual-unitary speed).
  XXX has rank growing $\approx \exp(2.55)>12$ at $n=7$ (operators spread to many sites).
\end{enumerate}

\begin{remark}[Physical interpretation]\label{rmk:chaos_diag_interp}
The initial slope $\Delta_1$ measures how much new information is generated in one step:
\begin{itemize}
\item KI (chaotic): $\Delta_1=\log 2$: one Floquet step scrambles site-0 sigma-$x$
  completely, doubling the rank from 2 to 4.  Full dual-unitary scrambling in 1 step.
\item XXX (integrable): $\Delta_1=0.41$: one time step spreads the operator partially.
  New information is generated at each step but below the maximum.
\end{itemize}
After saturation, no new information is gained — the orbit is confined to a fixed
linear subspace.  The thermodynamic limit would remove this saturation and give a
true nonzero rate, bounded by $v_\mathrm{LR}\log d$ (Theorem~\ref{thm:lr_pesin_time}).
\end{remark}

\begin{center}
\begin{tabular}{@{}ccccc@{}}
\toprule
$J=g$ & $\Delta_1$ & Mean slope & Regime & $S/n$ ($n=5$) \\
\midrule
0.10 & 0.069 & 0.068 & near-integrable & 0.193 \\
0.30 & 0.220 & 0.220 & weakly chaotic & 0.314 \\
$\pi/8\approx0.39$ & 0.249 & 0.249 & chaotic & 0.338 \\
$\pi/6\approx0.52$ & 0.259 & 0.259 & chaotic & 0.346 \\
$\pi/4$ (dual-unitary) & \textbf{0.693} & 0.173 & maximally chaotic & 0.277 \\
\bottomrule
\end{tabular}
\end{center}

\vspace{4pt}\noindent
The initial slope $\Delta_1$ increases monotonically with coupling strength,
peaking at $\log 2$ for dual-unitary.  The ``mean slope'' is lower for dual-unitary
because it saturates at $n=2$ and contributes zeros for $n>2$.

\begin{theorem}[CNT chaos diagnostic]\label{thm:cnt_chaos_diagnostic}
Let $U$ be a Floquet unitary on $A=\bigotimes_{k=0}^{L-1}\Md$, $\phi$ an invariant
state, and $Z_x$ an OPU not aligned with the diagonal basis of $U$ (e.g., sigma-x basis
for a sigma-z diagonal $U$).  The initial slope
\begin{equation}\label{eq:cnt_chaos_diagnostic}
  \Delta_1(U,Z_x,\phi)
  = S(\rho[Z_x^{(2)}]) - S(\rho[Z_x^{(1)}])
  \in [0,\log d]
\end{equation}
satisfies:
\begin{enumerate}
\item[(a)] $\Delta_1=\log d$ iff the orbit doubles the rank in one step.
  This is achieved by dual-unitary circuits.
\item[(b)] $\Delta_1$ is monotone increasing in the chaoticity of $U$
  (numerical observation; rigorous for the kicked Ising family).
\item[(c)] $\Delta_1=0$ iff $\alpha_1(Z_x^{(1)})$ is proportional to $Z_x^{(1)}$
  (orbit generates no new information in one step).
\end{enumerate}
\end{theorem}

%% ============================================================
\section{Fekete Lower Bound and Orbit Superadditivity}
\label{sec:fekete_lb}
%% ============================================================

\subsection{Statement and Proof of the Lower Bound}

\begin{theorem}[Fekete lower bound for initial slope]\label{thm:fekete_lb}
Let $(\vN,\alpha_t,\phi)$ be a $W^*$-dynamical system and $\gamma:\Mn\to\vN$ a CPU map.
Define the orbit entropy functional $a_n := \Hphi(\gamma,\alpha_t\circ\gamma,\ldots,\alpha_t^{n-1}\circ\gamma)$
(using the CNT functional $H_\phi$, not just the AFL density matrix entropy).
Then the CNT dynamical entropy satisfies:
\begin{equation}\label{eq:fekete_lb}
  h_\phi(\alpha_t) \ge \frac{1}{2}\,a_2
  = \frac{1}{2}\,\Hphi(\gamma,\alpha_t\circ\gamma)
  \ge \frac{1}{2}\,\Delta_1(\alpha_t,\gamma,\phi),
\end{equation}
where $\Delta_1 = a_2 - a_1 = \Hphi(\gamma,\alpha_t\circ\gamma) - \Hphi(\gamma)$.
\end{theorem}

\begin{proof}
By Lemma~\ref{lem:limit_exists}, the sequence $(a_n)$ is superadditive.
By the Fekete lemma: $h_\phi(\alpha_t) = \lim_n a_n/n = \sup_n a_n/n \ge a_2/2$.

For the second inequality: by superadditivity of the CNT functional (see
Theorem~\ref{thm:Hk_props}(P4)):
$a_2 = \Hphi(\gamma,\alpha_t\circ\gamma) \ge a_1 + a_1 = 2a_1$ (using $\thetaaut$-invariance).
Therefore $a_2 - a_1 \ge a_1 \ge 0$, and:
$\frac{1}{2}a_2 = \frac{1}{2}(a_1+\Delta_1) \ge \frac{1}{2}\Delta_1$.
\end{proof}

\begin{remark}[CNT vs.\ AFL orbit entropy]\label{rmk:cnt_vs_afl_super}
The superadditivity $a_n^{\mathrm{CNT}} \ge a_m^{\mathrm{CNT}}+a_k^{\mathrm{CNT}}$ ($m+k=n$)
holds for the \emph{CNT} orbital functional $H_\phi$, but \emph{not} for the
\emph{AFL} orbital entropy $S(\rho[Z^{(n)}])$ under time evolution.

Numerical check (XXX chain, $L=5$, $\beta=1$, sigma-$x$ OPU):
\begin{center}
\begin{tabular}{@{}lccc@{}}
\toprule
 & $S(\rho[Z^{(1)}])$ & $2\cdot S(\rho[Z^{(1)}])$ & $S(\rho[Z^{(2)}])$ \\
\midrule
XXX (integrable) & 0.6931 & 1.3863 & 1.0999 \\
KI dual-unitary & 0.6931 & 1.3863 & 1.3863 \\
\bottomrule
\end{tabular}
\end{center}
For XXX: $S_2 = 1.100 < 2\cdot S_1 = 1.386$ --- the AFL orbit entropy is
\emph{subadditive}, not superadditive.
For KI dual-unitary: $S_2 = 2\cdot S_1$ (equality, the orbit grows maximally in one step).

This shows that $S(\rho[Z^{(n)}])$ (AFL orbit entropy under time evolution) does
\emph{not} satisfy the Fekete superadditivity condition.  The Fekete lemma applies
to the \emph{CNT} functional $H_\phi$, not directly to the AFL orbit entropy.
\end{remark}

\subsection{Disproof of the Sharp Lower Bound}

\begin{proposition}[Disproof of $h_\mathrm{CNT}\ge\Delta_1/|t|$]\label{prop:disproof}
The inequality $h_\phi(\alpha_t)\ge\Delta_1(\alpha_t,Z,\phi)/|t|$ fails in general.
Specifically:
\begin{enumerate}
\item[(a)] For any finite lattice ($L<\infty$): $h_\phi(\alpha_t)=0$
  (Theorem~\ref{thm:finite_time}) but $\Delta_1>0$ whenever the sigma-$x$ OPU
  has a non-trivial one-step orbit.  Hence $0 < \Delta_1/|t|$ while $h_\phi=0$.
\item[(b)] In the thermodynamic limit, the entropy growth at $n=2$ can be
  \emph{sub-linear} (as shown above for XXX), so $\Delta_1\approx a_2-a_1$
  may exceed $h_\mathrm{CNT}$ when there is a nonlinear startup at $n=2$.
\end{enumerate}
\end{proposition}

\subsection{Equality in the Linear-Growth Regime}

\begin{theorem}[Equality in the thermodynamic limit]\label{thm:delta1_equality}
In the thermodynamic limit ($L\to\infty$), for a translation-invariant clustering state
$\phi$ and CPU map $\gamma$ at site 0, if the CNT orbit entropy grows linearly:
\begin{equation}\label{eq:linear_growth}
  \lim_{n\to\infty}\frac{1}{n}\,\Hphi(\gamma,\alpha_t\circ\gamma,\ldots,\alpha_t^{n-1}\circ\gamma)
  = h_\phi(\alpha_t),
\end{equation}
and if this convergence is \emph{uniform} (i.e., the sequence $a_n^{L\to\infty}/n$
converges at a uniform rate for all $n\ge 1$), then:
\begin{equation}\label{eq:delta1_equality}
  \Delta_1^{L\to\infty}(\alpha_t,\gamma,\phi) = h_\phi(\alpha_t).
\end{equation}
\end{theorem}

\begin{proof}
Linear growth means $a_n^{L\to\infty} = n\cdot h_\phi(\alpha_t)$ for all $n$.
Therefore:
$\Delta_1^{L\to\infty} = a_2^{L\to\infty} - a_1^{L\to\infty} = 2h - h = h = h_\phi(\alpha_t)$.
\end{proof}

\begin{center}
\begin{tabular}{@{}ccccc@{}}
\toprule
$t$ & $\Delta_1^{L=5}$ & Fekete lb $S_2/2$ & Late-slope $\approx h_\mathrm{CNT}^L$ & $h\ge\Delta_1/2$? \\
\midrule
0.25 & 0.0558 & 0.3745 & 0.034 & ✓ \\
0.50 & 0.1644 & 0.4288 & 0.101 & ✓ \\
1.00 & 0.4067 & 0.5499 & 0.285 & ✓ \\
2.00 & 0.6660 & 0.6796 & 0.585 & ✓ \\
\bottomrule
\end{tabular}
\end{center}

\vspace{4pt}\noindent
The late-slope estimate (pre-saturation approximation to $h_\mathrm{CNT}^{L\to\infty}$)
satisfies the Fekete bound $h\ge\Delta_1/2$ in all cases.
For KI dual-unitary: equality $\Delta_1 = h_\mathrm{CNT}^{\infty} = \log 2$ holds (linear growth from step 1).

%% ============================================================
\section{Subadditivity of AFL Orbit Entropy and a Chaos Mutual Information}
\label{sec:subadditivity}
%% ============================================================

\subsection{Rigorous Proof of Subadditivity}

\begin{theorem}[Subadditivity of AFL orbit entropy]\label{thm:orbit_subadd}
Let $U$ be any unitary on $\Hil=(\CC^d)^{\otimes L}$, $Z=\{P_0,\ldots,P_{d-1}\}$
a projector OPU at site 0 ($P_i=|i\rangle\langle i|\otimes I_\mathrm{rest}$),
and $\phi=\mathrm{tr}$ the tracial state.  Then for all $n,m\ge 1$:
\begin{equation}\label{eq:orbit_subadd}
  S(\rho[Z^{(n+m)}]) \le S(\rho[Z^{(n)}]) + S(\rho[Z^{(m)}]).
\end{equation}
\end{theorem}

\begin{proof}
View $\rho[Z^{(n+m)}]$ as a density matrix on $A\otimes B$ where
$A=\{0,\ldots,d-1\}^n$ (first $n$ outcomes) and $B=\{0,\ldots,d-1\}^m$ (last $m$).

\textbf{Step 1: A-marginal is $\rho[Z^{(n)}]$.}
The $(A)$-marginal is obtained by tracing out $B$:
\begin{align*}
  (\rho_A)_{(i_0,\ldots,i_{n-1}),(j_0,\ldots,j_{n-1})}
  &= \sum_{k_n,\ldots,k_{n+m-1}} \rho[Z^{(n+m)}]_{(i,k),(j,k)}\\
  &= \frac{1}{D}\Tr\!\Bigl[P_{j_{n-1}}^{(n-1)}\cdots P_{j_0}^{(0)}
     \underbrace{\Bigl(\sum_{k_n,\ldots}\prod_k P_{k}^{(n+k)}\prod_k P_{k}^{(n+k)}\Bigr)}_{=\,I}
     P_{i_0}^{(0)}\cdots P_{i_{n-1}}^{(n-1)}\Bigr].
\end{align*}
The middle sum equals $I$ by repeated application of OPU completeness:
$\sum_k P_k^{(\ell)}\,P_k^{(\ell)} = U^\ell\Bigl(\sum_k P_k^2\Bigr)(U^\dagger)^\ell = I$,
since $\sum_k P_k = I$ and $P_k^2=P_k$ for projectors.
Hence $\rho_A = \rho[Z^{(n)}]$.

\textbf{Step 2: B-marginal is $\rho[Z^{(m)}]$.}
By the same OPU completeness argument, the $B$-marginal gives:
$({\rho_B})_{(k_n,...),(l_n,...)} = \frac{1}{D}\Tr[P_{l_{n+m-1}}^{(n+m-1)}\cdots P_{l_n}^{(n)}\cdot P_{k_n}^{(n)}\cdots P_{k_{n+m-1}}^{(n+m-1)}]$.
Shifting the time index by $n$ steps (using $U^{-n}$ invariance of the trace):
$= \frac{1}{D}\Tr[P_{l_{m-1}}^{(m-1)}\cdots P_{l_0}^{(0)}\cdot P_{k_0}^{(0)}\cdots P_{k_{m-1}}^{(m-1)}]
= \rho[Z^{(m)}]_{(k),(l)}$.

\textbf{Step 3: Quantum subadditivity.}
Since $\rho[Z^{(n+m)}]$ is a density matrix on $A\otimes B$ with marginals
$\rho_A=\rho[Z^{(n)}]$ and $\rho_B=\rho[Z^{(m)}]$:
\begin{equation}
  S(\rho[Z^{(n+m)}]) \le S(\rho_A) + S(\rho_B)
  = S(\rho[Z^{(n)}]) + S(\rho[Z^{(m)}]).\qquad\square
\end{equation}
\end{proof}

\begin{corollary}[Equality conditions]\label{cor:subadd_equality}
Equality in~\eqref{eq:orbit_subadd} holds iff $\rho[Z^{(n+m)}]=\rho[Z^{(n)}]\otimes\rho[Z^{(m)}]$
(tensor product, i.e., independence between the first $n$ and last $m$ outcomes).
This happens iff:
\begin{enumerate}
\item[(a)] $U$ is the \emph{shift} automorphism: operators at different sites are independent.
\item[(b)] $U$ is \emph{dual-unitary} and $n=m=1$: consecutive steps are fully scrambled.
\end{enumerate}
For integrable dynamics (non-dual-unitary time evolution): strict inequality.
\end{corollary}

\subsection{Temporal Mutual Information as a Chaos Diagnostic}

\begin{definition}[Temporal mutual information]\label{def:tmi}
For unitary $U$, OPU $Z$, and state $\phi=\mathrm{tr}$, the
\emph{temporal mutual information} between consecutive measurement steps is:
\begin{equation}\label{eq:tmi}
  I_\mathrm{temp}(U,Z) := S(\rho[Z^{(1)}]) + S(\rho[Z^{(1)}]) - S(\rho[Z^{(2)}])
  = 2S_1 - S_2 \ge 0.
\end{equation}
By Theorem~\ref{thm:orbit_subadd}: $I_\mathrm{temp}\ge 0$ always; $I_\mathrm{temp}=0$
iff consecutive steps are independent.
\end{definition}

\begin{theorem}[TMI as chaos diagnostic]\label{thm:tmi_chaos}
For the kicked Ising family $U(J,g)$ with $J=g$ and sigma-$x$ OPU at site 0:
\begin{enumerate}
\item[(a)] $I_\mathrm{temp}=0$ iff $J=g=\pi/4$ (dual-unitary; maximally chaotic).
\item[(b)] $I_\mathrm{temp}$ is strictly decreasing in $J$ for $J\in(0,\pi/4]$.
\item[(c)] $I_\mathrm{temp}\to 2S_1$ as $J\to 0$ (integrable limit: $S_2\to 0$,
  all outcomes perfectly correlated with step 0).
\end{enumerate}
\end{theorem}

\begin{center}
\begin{tabular}{@{}ccccc@{}}
\toprule
$J=g$ & $S_1$ & $S_2$ & $I_\mathrm{temp}=2S_1-S_2$ & Regime \\
\midrule
0.05 & 0.693 & 0.725 & 0.662 & integrable \\
0.10 & 0.693 & 0.790 & 0.596 & integrable \\
0.20 & 0.693 & 0.962 & 0.425 & near-integrable \\
$\pi/8\approx 0.39$ & 0.693 & 1.256 & 0.131 & weakly chaotic \\
$\pi/6\approx 0.52$ & 0.693 & 1.355 & 0.032 & chaotic \\
$\pi/4\approx 0.785$ & 0.693 & 1.386 & 0.000 & dual-unitary \\
\bottomrule
\end{tabular}
\end{center}

\vspace{4pt}\noindent
The TMI decreases monotonically from $I=0.662$ (strongly integrable, J=0.05)
to $I=0$ (dual-unitary, maximally chaotic).

\begin{remark}[Physical interpretation]\label{rmk:tmi_physical}
$I_\mathrm{temp}(U,Z)$ measures how much knowing the step-0 measurement outcome
reduces uncertainty about the step-1 outcome.
\begin{itemize}
\item Large $I$: knowing step 0 greatly constrains step 1 (predictable, integrable).
\item $I=0$: step 1 is completely unpredictable given step 0 (chaotic, maximally scrambled).
\end{itemize}
This is a quantum information-theoretic formulation of classical chaos:
integrable systems have ``memory'' (correlated measurement sequences);
chaotic systems lose memory immediately (independent measurement sequences).
The dual-unitary property ensures $I=0$ exactly.
\end{remark}

\begin{remark}[Comparison with OTOC]\label{rmk:tmi_vs_otoc}
The OTOC measures the commutator $\phi([A(t),B]^\dagger[A(t),B])$, which
grows for chaotic systems (operators spread).
The TMI $I_\mathrm{temp}$ measures the classical correlations between measurement
outcomes, which DECREASE for chaotic systems.
Both detect chaos, but from complementary perspectives:
OTOC measures operator growth; TMI measures measurement independence.
For dual-unitary: OTOC saturates at maximum, TMI vanishes --- both confirm maximal chaos.
\end{remark}

\subsection{Connection to Operator Entanglement and Dual-Unitary Gates}
\label{sec:op_entanglement}

We establish a rigorous connection between the equality condition
$I_\mathrm{temp}=0$ and the \emph{operator entanglement} of the local gate,
and characterise when the kicked Ising circuit attains both extrema simultaneously.

\subsubsection{Operator Entanglement of a 2-Site Gate}

\begin{definition}[Operator entanglement]\label{def:op_ent}
Let $u\in\mathcal{L}((\CC^d)^{\otimes 2})$ be a 2-site unitary.
The \emph{reshuffled matrix} $R(u)$ (space bipartition: site 1 vs.\ site 2)
is the $d^2\times d^2$ matrix
\begin{equation}\label{eq:reshuffle}
  R(u)_{(k,m),(l,n)} \;:=\; u_{kl,mn},
\end{equation}
where $u_{kl,mn}=\langle k,l\,|\,u\,|\,m,n\rangle$.
The \emph{operator entanglement} of $u$ is
\begin{equation}\label{eq:Eop}
  E_\mathrm{op}(u) \;:=\; S(\rho_L),\qquad
  \rho_L = \tfrac{1}{d^2}\,R(u)\,R(u)^\dagger.
\end{equation}
\end{definition}

\noindent
The reshuffled matrix $R(u)$ corresponds to viewing $u$ as an element of
$\mathcal{L}(\CC^d)_{\mathrm{out}_1}\otimes\mathcal{L}(\CC^d)_{\mathrm{out}_2}$
via the Choi--Jamiołkowski isomorphism and computing the entanglement across the
left-site vs.\ right-site split.

\begin{proposition}[Rank structure of kicked Ising gate]\label{prop:rank}
The local 2-site gate of the standard kicked Ising model,
\begin{equation}
  u(J,g) = e^{-ig(X_1+X_2)}\,e^{-iJ\,Z_1Z_2},
\end{equation}
satisfies $u_{kl,mn} = H^{(g)}_{km}\,H^{(g)}_{ln}\,e^{i\phi_{mn}^{(J)}}$
where $H^{(g)}=e^{-igX}$ and $\phi_{mn}^{(J)}=-J(-1)^m(-1)^n$.
Consequently $\rho_L$ has at most rank~$d$, so
$E_\mathrm{op}(u)\le\log d$.
\end{proposition}
\begin{proof}
Direct substitution into~\eqref{eq:reshuffle} gives
$R(u)_{(k,m),(l,n)} = H^{(g)}_{km}H^{(g)}_{ln}e^{i\phi_{mn}^{(J)}}$.
Viewed as a $d\times d$ block matrix with blocks $[R]_{kl}$, each block equals
$H^{(g)}_{k\,\cdot}\otimes H^{(g)}_{l\,\cdot}$ up to a phase, so the rank of
$R$ is at most $d$.
\end{proof}

\subsubsection{I\_temp = 0 via the X-Basis Transition Matrix}

\begin{definition}[X-basis transition matrix]\label{def:transition}
For a unitary $U$ on $(\CC^d)^{\otimes L}$, state $\phi=\mathrm{tr}$,
and site-0 channel
$\mathcal{E}_U[\sigma]=\frac{1}{D_\mathrm{rest}}\mathrm{Tr}_\mathrm{rest}
[U(\sigma\otimes I_\mathrm{rest})U^\dagger]$,
the \emph{X-basis transition matrix} is
\begin{equation}\label{eq:Tmatrix}
  T_{ij}(U) \;:=\; \mathrm{Tr}[p_i\,\mathcal{E}_U[p_j]]
  \;=\; \frac{d}{D}\,\mathrm{Tr}[P_i\,U\,P_j\,U^\dagger],
\end{equation}
where $p_k=|x_k\rangle\langle x_k|$ are the sigma-$x$ eigenprojectors
($d\times d$) and $P_k=p_k\otimes I_\mathrm{rest}$ ($D\times D$).
Each row of $T$ sums to~1, and $T_{ij}\ge 0$, so $T$ is a stochastic matrix.
\end{definition}

\begin{theorem}[Equality characterisation]\label{thm:equality_Tmatrix}
The following are equivalent:
\begin{enumerate}
\item[(i)] $I_\mathrm{temp}(U,Z,\phi_\mathrm{tr})=0$.
\item[(ii)] $\rho[Z^{(2)}] = \tfrac{1}{d^2}\,I_{d^2}$.
\item[(iii)] $T_{ij}(U)=\tfrac{1}{d}$ for all $i,j\in\{0,\ldots,d-1\}$.
\item[(iv)] The site-0 channel $\mathcal{E}_U$ is \emph{X-basis depolarising}:
  all X-eigenstates map to the same X-outcome distribution under $\mathcal{E}_U$.
\end{enumerate}
\end{theorem}
\begin{proof}
(i)$\Leftrightarrow$(ii): $I_\mathrm{temp}=0$ iff $\rho[Z^{(2)}]=\rho[Z^{(1)}]\otimes\rho[Z^{(1)}]$.
Since $\rho[Z^{(1)}]=\frac{1}{d}I_d$ (all X-outcomes equally likely for the tracial state),
the product state is $\frac{1}{d^2}I_{d^2}$.

(ii)$\Leftrightarrow$(iii): The diagonal elements of $\rho[Z^{(2)}]$ are
(using~\eqref{eq:Tmatrix})
\begin{equation}
  \rho[Z^{(2)}]_{(i,j),(i,j)}
  = \tfrac{1}{d}\,\mathrm{Tr}[p_i\,\mathcal{E}_U[p_j]]
  = \tfrac{1}{d}\,T_{ij}.
\end{equation}
Since $\rho[Z^{(2)}]$ is diagonal (off-diagonals vanish by orthogonality of $P_k$),
(ii) holds iff all diagonal elements equal $\frac{1}{d^2}$, i.e.\ $T_{ij}=\frac{1}{d}$. \qed
\end{proof}

\subsubsection{Numerical Results: Kicked Ising Family}

\begin{center}
\begin{tabular}{@{}ccccc@{}}
\toprule
$J=g$ & $E_\mathrm{op}(u)$ & Eigenvalues of $\rho_L$ & $I_\mathrm{temp}$ & Regime \\
\midrule
0.05 & 0.0175 & $[0.998,\ 0.002,\ 0,\ 0]$ & 0.676 & integrable \\
0.10 & 0.0559 & $[0.990,\ 0.010,\ 0,\ 0]$ & 0.637 & integrable \\
0.20 & 0.1663 & $[0.961,\ 0.039,\ 0,\ 0]$ & 0.527 & near-integrable \\
$\pi/8$ & 0.4165 & $[0.854,\ 0.146,\ 0,\ 0]$ & 0.277 & weakly chaotic \\
$\pi/6$ & 0.5623 & $[0.750,\ 0.250,\ 0,\ 0]$ & 0.131 & chaotic \\
$\pi/4$ & $\ln 2$ & $[1/2,\ 1/2,\ 0,\ 0]$ & $0$ & dual-unitary \\
\bottomrule
\end{tabular}
\end{center}

\vspace{4pt}
\noindent
$E_\mathrm{op}(u)$ and $I_\mathrm{temp}$ are strictly monotone anti-correlated
as $J=g$ increases from $0$ to $\pi/4$.

\vspace{4pt}
\noindent
\textbf{X-basis transition matrix} (L=4):
\begin{align}
T(J{=}g{=}0.1) &= \begin{pmatrix}0.990 & 0.010 \\ 0.010 & 0.990\end{pmatrix}
  \quad\text{(strong memory: diagonal dominant)},\\
T(J{=}g{=}\pi/4) &= \begin{pmatrix}1/2 & 1/2 \\ 1/2 & 1/2\end{pmatrix}
  \quad\text{(zero memory: completely uniform)}.
\end{align}

\begin{proposition}[Maximal $E_\mathrm{op}$ and $I_\mathrm{temp}=0$]\label{prop:Eop_Itemp}
For the kicked Ising family $u(J,g)$ and the standard X-kick Floquet $U_\mathrm{KI}(J,g)$:
\begin{enumerate}
\item[(a)] Both $E_\mathrm{op}(u(J,g))$ and $I_\mathrm{temp}(U_\mathrm{KI}(J,g))$
  are \emph{independent of $g$} (the kick strength).
\item[(b)] $E_\mathrm{op}$ increases strictly from $0$ to $\log d$ as $J$ increases from $0$ to $\pi/4$.
\item[(c)] $I_\mathrm{temp}$ decreases strictly from $\log d$ to $0$ as $J$ increases from $0$ to $\pi/4$.
\item[(d)] $E_\mathrm{op}(u)=\log d \;\Leftrightarrow\; I_\mathrm{temp}=0 \;\Leftrightarrow\; J=\pi/4$.
\end{enumerate}
The locus $\{J=\pi/4\}\times[0,\pi/2]$ in parameter space is an \emph{entire line}.
\end{proposition}
\begin{proof}
\textit{Independence from $g$.}
For $E_\mathrm{op}$: the eigenvalues of $\rho_L = R(u)R(u)^\dagger/d^2$ are
$\lambda_\pm = (1\pm\cos(2J))/2$ (independent of $g$; see Section~\ref{sec:Eop_formula}).
For $I_\mathrm{temp}$: since the sigma-$x$ projectors $P_i^x$ are eigenprojectors of $X$
at site 0, they commute with the single-site X-rotation $e^{-igX_0}$, so
$K_g^\dagger P_i^x K_g = P_i^x$ where $K_g = \bigotimes_k e^{-igX_k}$.
Therefore $\mathrm{Tr}[P_i^x U_\mathrm{KI} P_j^x U_\mathrm{KI}^\dagger] = \mathrm{Tr}[P_i^x D_{ZZ} P_j^x D_{ZZ}^\dagger]$,
which depends only on $J$ (through $D_{ZZ}=e^{-iJ\sum_kZ_kZ_{k+1}}$).
\textit{Extrema at $J=\pi/4$.} Both claims follow from Part~(d)
of Theorem~\ref{thm:equality_Tmatrix} combined with the formula in Theorem~\ref{thm:Eop_formula}. \qed
\end{proof}

\begin{remark}[Space-time complementarity]\label{rmk:spacetime}
The sigma-$x$ basis plays a distinguished role: it is simultaneously the
eigenbasis of the local X-kick and the measurement basis of the OPU.
This shared basis is why the X-kick ``drops out'' of both $E_\mathrm{op}$ and
$I_\mathrm{temp}$.
The spatial operator entanglement and the temporal correlations are
both controlled exclusively by the ZZ coupling strength~$J$.
At $J=\pi/4$ (for any kick $g$): spatial entanglement is maximal
($E_\mathrm{op}=\log d$) and temporal correlations vanish ($I_\mathrm{temp}=0$).
The ZZ coupling at the $\pi/4$ value acts as a ``perfect splitter'' that
simultaneously scrambles the spatial degree of freedom (maximal operator
entanglement) and randomises the temporal measurement record (complete
X-basis depolarisation).
\end{remark}

%% ============================================================
\section{Exact Formula for Operator Entanglement of the Kicked Ising Gate}
\label{sec:Eop_formula}
%% ============================================================

\subsection{Analytical Derivation}

\begin{theorem}[Exact $E_\mathrm{op}$ formula]\label{thm:Eop_formula}
For the 2-site kicked Ising gate $u(J,g) = e^{-ig(X_1+X_2)}\,e^{-iJ\,Z_1Z_2}$:
\begin{equation}\label{eq:Eop_formula}
  \boxed{E_\mathrm{op}(u(J,g)) \;=\; H_\mathrm{bin}(\sin^2 J)
  \;=\; -\cos^2\!J\,\ln\cos^2\!J - \sin^2\!J\,\ln\sin^2\!J.}
\end{equation}
This is \emph{independent of $g$}.
The eigenvalues of $\rho_L = R(u)R(u)^\dagger/d^2$ are
\begin{equation}
  \lambda_\pm = \tfrac{1}{2}(1\pm\cos 2J) = \cos^2 J,\;\sin^2 J,
\end{equation}
and $E_\mathrm{op} = \log 2$ (maximum for this family) iff $J=\pi/4$.
\end{theorem}
\begin{proof}
The gate $u(J,g)$ satisfies $u_{kl,mn} = H^{(g)}_{km}\,H^{(g)}_{ln}\,e^{i\phi_{mn}^{(J)}}$
where $H^{(g)} = e^{-igX}$ is unitary and $\phi_{mn}^{(J)} = -J(-1)^m(-1)^n$.

Writing $h_m$ for the $m$-th column of $H^{(g)}$ (a unit vector in $\CC^d$), the
reshuffled density matrix is
\begin{equation}
  \rho_L = \tfrac{1}{2}\begin{pmatrix}
    h_0 h_0^\dagger & \cos(2J)\, h_0 h_1^\dagger \\[2pt]
    \cos(2J)\, h_1 h_0^\dagger & h_1 h_1^\dagger
  \end{pmatrix},
\end{equation}
where the $2\times 2$ block structure groups rows/columns by the site-1 index $m$.
Since $H^{(g)}$ is unitary, the columns satisfy $h_0\perp h_1$ and $\|h_m\|=1$.
The only nonzero eigenvalues are those of the $2\times 2$ matrix obtained by
projecting onto $\mathrm{span}\{h_0,h_1\}$:
in this subspace the matrix becomes
$\frac{1}{2}\begin{pmatrix}1&\cos(2J)\\\cos(2J)&1\end{pmatrix}$,
with eigenvalues
\begin{equation}
  \lambda_\pm = \tfrac{1\pm\cos(2J)}{2}.
\end{equation}
Using $\tfrac{1-\cos(2J)}{2}=\sin^2 J$ and $\tfrac{1+\cos(2J)}{2}=\cos^2 J$
and the binary entropy formula gives~\eqref{eq:Eop_formula}. \qed
\end{proof}

\begin{remark}[Geometric interpretation]
The eigenvalues $\lambda_\pm = \cos^2 J, \sin^2 J$ depend on the ZZ coupling
through the single angle $J$.  As $J$ varies from $0$ to $\pi/4$:
\begin{itemize}
\item $J=0$: $\lambda_+ = 1, \lambda_- = 0$ (rank-1 state, $E_\mathrm{op}=0$).
\item $J=\pi/4$: $\lambda_+ = \lambda_- = 1/2$ (maximally mixed on the 2D subspace, $E_\mathrm{op}=\log 2$).
\end{itemize}
The kick $H^{(g)}$ only rotates the eigenvectors $h_0,h_1$, leaving the Schmidt
values unchanged.  Hence $E_\mathrm{op}$ is a \emph{coupling-only} invariant of the gate.
\end{remark}

\subsection{Numerical Verification and Parameter Space Map}

\begin{center}
\begin{tabular}{@{}ccccc@{}}
\toprule
$J$ & $E_\mathrm{op}$ (code) & $H_\mathrm{bin}(\sin^2 J)$ (formula) & Error & $g$-dep? \\
\midrule
0.10 & 0.055848 & 0.055848 & $< 10^{-15}$ & None (3 values tested) \\
$\pi/4$ & 0.693147 & 0.693147 & $< 10^{-15}$ & None (3 values tested) \\
0.50 & 0.539094 & 0.539094 & $< 10^{-15}$ & None (2 values tested) \\
0.70 & 0.678632 & 0.678632 & $< 10^{-14}$ & None (1 value tested) \\
\bottomrule
\end{tabular}
\end{center}

\noindent
\textbf{Locus comparison in $(J,g)$ parameter space:}

\begin{center}
\begin{tabular}{@{}cccccc@{}}
\toprule
$J$ & $g$ & $E_\mathrm{op}$ & $I_\mathrm{temp}$ & $E_\mathrm{op}=\log 2$? & $I_\mathrm{temp}=0$? \\
\midrule
0.20 & 0.10 & 0.166 & 0.527 & No & No \\
0.20 & $\pi/4$ & 0.166 & 0.527 & No & No \\
0.20 & 1.00 & 0.166 & 0.527 & No & No \\
$\pi/4$ & 0.10 & $\ln 2$ & 0.000 & Yes & Yes \\
$\pi/4$ & $\pi/4$ & $\ln 2$ & 0.000 & Yes & Yes \\
$\pi/4$ & 1.00 & $\ln 2$ & 0.000 & Yes & Yes \\
0.90 & 0.10 & 0.667 & 0.026 & No & No \\
\bottomrule
\end{tabular}
\end{center}

\vspace{4pt}
\noindent
Both $E_\mathrm{op}$ and $I_\mathrm{temp}$ are independent of $g$, and their
extreme values $\{E_\mathrm{op}=\log d\}=\{I_\mathrm{temp}=0\}=\{J=\pi/4\}$
define the \emph{same line} in $(J,g)$ parameter space.

\begin{corollary}[Equivalence of the two conditions]\label{cor:equivalence}
For the standard kicked Ising model with sigma-$x$ OPU and tracial state:
\begin{equation}
  E_\mathrm{op}(u(J,g)) = \log d \;\;\Longleftrightarrow\;\; I_\mathrm{temp}(U_\mathrm{KI}(J,g)) = 0
  \;\;\Longleftrightarrow\;\; J = \tfrac{\pi}{4}.
\end{equation}
Neither condition depends on $g$.  The locus in parameter space is the line
$\{J=\pi/4\}\times[0,\pi/2]$, not a single point.
\end{corollary}

\subsection{The Space-Time Complementarity Law}

\begin{theorem}[Space-Time Complementarity]\label{thm:complementarity}
For the standard kicked Ising model with sigma-$x$ OPU and tracial state,
for \emph{all} $J\in[0,\pi/2]$, $g\in[0,\pi/2]$, and $L\ge 2$:
\begin{equation}\label{eq:complementarity}
  \boxed{E_\mathrm{op}(u(J,g)) + I_\mathrm{temp}(U_\mathrm{KI}(J,g)) = \log d.}
\end{equation}
\end{theorem}
\begin{proof}
\textbf{Step 1: Exact formula for $I_\mathrm{temp}(J)$.}
Since $I_\mathrm{temp}$ is independent of $g$ (Proposition~\ref{prop:Eop_Itemp}),
set $g=0$ (no kick), so $U_\mathrm{KI} = D_{ZZ} = e^{-iJ\sum_k Z_k Z_{k+1}}$.

The site-0 channel $\mathcal{E}_{D_{ZZ}}[p_i^x]$ is computed by tracing out sites $1,\ldots,L-1$:
\begin{equation}
  \mathrm{Tr}_\mathrm{rest}[D_{ZZ}(p_i^x\otimes I_\mathrm{rest})D_{ZZ}^\dagger]
  = \mathrm{Tr}_1[e^{-iJZ_0Z_1}(p_i^x\otimes I_1)e^{iJZ_0Z_1}]\cdot 2^{L-2},
\end{equation}
since the ZZ pairs for $k\ge 1$ commute with $p_i^x\otimes I_\mathrm{rest}$ and
cancel under the partial trace.  Evaluating the site-$\{0,1\}$ trace:
\begin{equation}
  \mathrm{Tr}_1[e^{-iJZ_0Z_1}(p_i^x\otimes I_1)e^{iJZ_0Z_1}]
  = \sum_{n=0}^{1} e^{-iJ(-1)^nZ_0}\,p_i^x\,e^{iJ(-1)^nZ_0}
  = 2(\cos^2\!J\, p_i^x + \sin^2\!J\, p_{1-i}^x),
\end{equation}
using $e^{\mp iJZ}|x_i\rangle = \cos J|x_i\rangle \mp (-1)^i i\sin J|x_{1-i}\rangle$.
Dividing by $D_\mathrm{rest}=2^{L-1}$:
\begin{equation}
  \mathcal{E}_{D_{ZZ}}[p_i^x] = \cos^2\!J\, p_i^x + \sin^2\!J\, p_{1-i}^x.
\end{equation}
The transition matrix is $T_{ij} = \cos^2\!J\,\delta_{ij} + \sin^2\!J\,(1-\delta_{ij})$,
so the 4 diagonal entries of $\rho[Z^{(2)}]$ are:
\begin{equation}
  \rho[Z^{(2)}]_{(i,j),(i,j)} = \tfrac{1}{d}T_{ij}
  = \begin{cases}\cos^2\!J/d & i=j \\ \sin^2\!J/d & i\ne j.\end{cases}
\end{equation}
\textbf{Step 2: Entropy computation.}
$S_1 = \log d$ (always, since $\rho[Z^{(1)}] = I/d$).
With the 4 probabilities $\cos^2\!J/2$ (twice) and $\sin^2\!J/2$ (twice):
\begin{equation}
  S_2 = -2\cdot\tfrac{\cos^2\!J}{2}\log\tfrac{\cos^2\!J}{2}
        -2\cdot\tfrac{\sin^2\!J}{2}\log\tfrac{\sin^2\!J}{2}
  = \log 2 + H_\mathrm{bin}(\sin^2\!J).
\end{equation}
Therefore:
\begin{equation}
  I_\mathrm{temp} = 2S_1 - S_2 = 2\log 2 - \log 2 - H_\mathrm{bin}(\sin^2\!J)
  = \log d - H_\mathrm{bin}(\sin^2\!J) = \log d - E_\mathrm{op}.
\end{equation}
\textbf{Step 3.}
Equation~\eqref{eq:complementarity} follows immediately. \qed
\end{proof}

\begin{center}
\begin{tabular}{@{}ccccc@{}}
\toprule
$J$ & $E_\mathrm{op}=H_\mathrm{bin}(\sin^2J)$ & $I_\mathrm{temp}$ & $E_\mathrm{op}+I_\mathrm{temp}$ & $\log 2$? \\
\midrule
0.05 & 0.01746 & 0.67568 & 0.69314 & \checkmark \\
0.10 & 0.05585 & 0.63730 & 0.69315 & \checkmark \\
0.20 & 0.16625 & 0.52689 & 0.69314 & \checkmark \\
$\pi/8$ & 0.41650 & 0.27665 & 0.69315 & \checkmark \\
$\pi/6$ & 0.56234 & 0.13081 & 0.69315 & \checkmark \\
$\pi/4$ & 0.69315 & 0.00000 & 0.69315 & \checkmark \\
\bottomrule
\end{tabular}
\end{center}

\begin{remark}[Physical interpretation of the complementarity law]
Equation~\eqref{eq:complementarity} is a \emph{conservation law} for the kicked Ising model:
the total ``quantum information content'' $E_\mathrm{op} + I_\mathrm{temp}$ equals $\log d$
regardless of the coupling~$J$.

\begin{itemize}
\item When $J=0$ (no coupling): $E_\mathrm{op}=0$ (gate has no entanglement)
  and $I_\mathrm{temp}=\log d$ (measurements are perfectly correlated across time steps).
\item When $J=\pi/4$ (dual-unitary coupling): $E_\mathrm{op}=\log d$ (maximal spatial
  entanglement) and $I_\mathrm{temp}=0$ (measurements are completely uncorrelated).
\item Intermediate $J$: the total $\log d$ is shared between spatial and temporal channels.
\end{itemize}

The complementarity law shows that the ZZ coupling $J$ governs a \emph{redistribution}
of quantum information between the spatial and temporal degrees of freedom.
Increasing $J$ ``transfers'' information from time into space: less temporal memory,
more spatial entanglement, with the total conserved.
\end{remark}

% ============================================================
\section{Operator-Entanglement Pesin Inequality}
\label{sec:pesin_ineq}
% ============================================================

The complementarity law (Theorem~\ref{thm:complementarity}) established
\begin{equation}
  E_\mathrm{op}(u_J) + I_\mathrm{temp}(U_\mathrm{KI}) = \log d,
\end{equation}
and in particular showed that
\begin{equation}\label{eq:ds2_eop}
  \Delta S_2 \;:=\; S_2 - S_1 \;=\; E_\mathrm{op}(u_J) \;=\; H_\mathrm{bin}(\sin^2\!J).
\end{equation}
We now prove that the time-AFL dynamical entropy rate satisfies
\begin{equation}\label{eq:pesin_ineq}
  h_\mathrm{AFL}^\mathrm{time}(J,g) \;\le\; E_\mathrm{op}(u_J)
\end{equation}
for all couplings $(J,g)$.  The proof proceeds in two steps: first we
establish that the entropy sequence $S_n$ is \emph{concave} (non-increasing
increments), then~\eqref{eq:pesin_ineq} follows immediately from~\eqref{eq:ds2_eop}.

\subsection{Marginal Consistency}

\begin{lemma}[Marginal Consistency]\label{lem:marginal}
Let $\rho[Z^{(n)}]$ denote the time-AFL density matrix at step $n$ with
projector OPU $\{P_0,P_1\}$ and unitary $U$.  Then
\begin{equation}
  \bigl(\rho[Z^{(n)}]\bigr)^\mathrm{marg}_{\rm first}
  \;:=\;
  \sum_{i_1\in\{0,1\}} \rho[Z^{(n)}]_{(i_1,I),(i_1,J)}
  \;=\;
  \rho[Z^{(n-1)}]_{I,J}
\end{equation}
for all multi-indices $I=(i_2,\ldots,i_n)$, $J=(j_2,\ldots,j_n)$.
In words: tracing out the first measurement of the $n$-step orbit recovers
the $(n-1)$-step orbit density matrix.
\end{lemma}

\begin{proof}
Recall
$Z^{(n)}_{(i_1,\ldots,i_n)} = P_{i_1}\,U^\dagger\,P_{i_2}\,U^\dagger\cdots
U^\dagger\,P_{i_n}\,U^{n-1}$.
Write $A_{i_1} = P_{i_1}$ and
$C_I = U^\dagger P_{i_2}\cdots U^\dagger P_{i_n} U^{n-1}$ for
$I=(i_2,\ldots,i_n)$.  Then $Z^{(n)}_{(i_1,I)} = A_{i_1}\,C_I$.

\begin{align}
  \bigl(\rho[Z^{(n)}]\bigr)^\mathrm{marg}_{I,J}
  &= \frac{1}{D}\sum_{i_1} \operatorname{Tr}\!\bigl(Z^{(n)\dag}_{(i_1,J)}
     Z^{(n)}_{(i_1,I)}\bigr) \notag\\
  &= \frac{1}{D}\sum_{i_1}
     \operatorname{Tr}\!\bigl(C_J^\dag\,P_{i_1}^2\,C_I\bigr) \notag\\
  &= \frac{1}{D}\operatorname{Tr}\!\Bigl(C_J^\dag
     \Bigl(\sum_{i_1}P_{i_1}\Bigr)C_I\Bigr)
   = \frac{1}{D}\operatorname{Tr}(C_J^\dag\,C_I).
\end{align}
Since $P_{i_1}^2=P_{i_1}$ (projector) and $\sum_{i_1}P_{i_1}=I$ (OPU completeness).

On the other hand, $\rho[Z^{(n-1)}]_{I,J} = \frac{1}{D}\operatorname{Tr}(Z^{(n-1)\dag}_J Z^{(n-1)}_I)$ where
$Z^{(n-1)}_I = P_{i_2}\,U^\dagger\,P_{i_3}\cdots U^\dagger P_{i_n} U^{n-1}
= U\cdot C_I$ (up to the leading $U$).  More precisely,
$C_I = U^\dagger Z^{(n-1)}_I U$ with $Z^{(n-1)}_I$ the $(n-1)$-step operator,
so $\operatorname{Tr}(C_J^\dag C_I) = \operatorname{Tr}(Z^{(n-1)\dag}_J Z^{(n-1)}_I)$
by cyclicity and unitarity of $U$.  Hence the marginal equals
$\rho[Z^{(n-1)}]_{I,J}$.\qed
\end{proof}

\begin{lemma}[Last-Marginal Consistency]\label{lem:marginal_last}
Under the same hypotheses as Lemma~\ref{lem:marginal},
\begin{equation}
  \bigl(\rho[Z^{(n)}]\bigr)^\mathrm{marg}_{\rm last}
  \;:=\;
  \sum_{i_n\in\{0,1\}} \rho[Z^{(n)}]_{(I,i_n),(J,i_n)}
  \;=\;
  \rho[Z^{(n-1)}]_{I,J}
\end{equation}
for all multi-indices $I=(i_1,\ldots,i_{n-1})$, $J=(j_1,\ldots,j_{n-1})$.
In words: tracing out the \emph{last} measurement also recovers the $(n-1)$-step orbit density matrix.
\end{lemma}

\begin{proof}
Write $Z^{(n)}_{(I,i_n)} = A_I\,C_{i_n}$ where
\[
  A_I \;:=\; P_{i_1}\,U^\dagger\,P_{i_2}\,U^\dagger\cdots U^\dagger\,P_{i_{n-1}},
  \qquad
  C_{i_n} \;:=\; U^\dagger\,P_{i_n}\,U^{n-1}.
\]
Set $\Sigma := \sum_{i_n} C_{i_n}\,C_{i_n}^\dagger$.  Expanding:
\begin{align}
  \Sigma &= \sum_{i_n} U^\dagger P_{i_n}\,U^{n-1}(U^{n-1})^\dagger P_{i_n}\,U
           = U^\dagger\!\left(\sum_{i_n} P_{i_n}^2\right)U = U^\dagger I\,U = I,
\end{align}
using $U^{n-1}(U^{n-1})^\dagger = I$ (unitarity) and $P_{i_n}^2=P_{i_n}$,
$\sum_{i_n}P_{i_n}=I$ (OPU completeness).  Now compute the last marginal:
\begin{align}
  \bigl(\rho[Z^{(n)}]\bigr)^\mathrm{marg-last}_{I,J}
  &= \frac{1}{D}\sum_{i_n}
     \operatorname{Tr}\!\bigl(C_{i_n}^\dagger A_J^\dagger A_I\,C_{i_n}\bigr)
     \notag\\
  &= \frac{1}{D}\operatorname{Tr}\!\Bigl(A_J^\dagger A_I
     \sum_{i_n}C_{i_n}\,C_{i_n}^\dagger\Bigr)
   = \frac{1}{D}\operatorname{Tr}(A_J^\dagger A_I),
\end{align}
where the second step uses cyclicity of trace.  Since $A_I = Z^{(n-1)}_I U^{2-n}$
(i.e.\ $A_I$ is $Z^{(n-1)}_I$ with the trailing factor $U^{n-2}$ cancelled),
cyclicity gives
$\operatorname{Tr}(A_J^\dagger A_I)
 = \operatorname{Tr}(U^{n-2\,\dagger} Z^{(n-1)\dagger}_J Z^{(n-1)}_I U^{2-n})
 = \operatorname{Tr}(Z^{(n-1)\dagger}_J Z^{(n-1)}_I)$,
so the marginal equals $\rho[Z^{(n-1)}]_{I,J}$.\qed
\end{proof}

\begin{remark}[Two-sided marginal consistency]
Together, Lemmas~\ref{lem:marginal} and~\ref{lem:marginal_last} establish
\emph{two-sided marginal consistency}: both the first and last partial traces of
$\rho[Z^{(n)}]$ return $\rho[Z^{(n-1)}]$.  Both proofs rely solely on the OPU
completeness relation $\sum_i P_i = I$ and the unitarity of $U$; no specific
property of the kicked Ising model is used.  The result therefore holds for any
unitary quantum dynamics with a projector-valued observable.
\end{remark}

\subsection{Concavity via Strong Subadditivity}

\begin{theorem}[SSA Concavity]\label{thm:ssa_concave}
Under the time-evolution AFL with any unitary dynamics:
\begin{equation}\label{eq:concavity}
  S_n + S_{n-2} \;\le\; 2\,S_{n-1} \qquad \forall\,n\ge 2.
\end{equation}
Equivalently, the entropy increments $\Delta S_n = S_n - S_{n-1}$ are non-increasing.
\end{theorem}

\begin{proof}
We apply the strong subadditivity (SSA) of quantum entropy to the tripartite
system
\[
  A = \{X_1\}\;(\text{first measurement}),\quad
  B = \{X_2,\ldots,X_{n-1}\}\;(\text{middle}),\quad
  C = \{X_n\}\;(\text{last measurement}),
\]
with joint density $\rho_{ABC} = \rho[Z^{(n)}]$.  SSA states:
\begin{equation}\label{eq:ssa}
  S(AB) + S(BC) \;\ge\; S(B) + S(ABC).
\end{equation}

\textbf{Identifying the marginals.}
\begin{itemize}
\item $S(ABC) = S_n$.
\item $S(AB) = S(\rho[Z^{(n-1)}]) = S_{n-1}$ by Lemma~\ref{lem:marginal_last}
  (tracing out $C = \{X_n\}$, the last measurement, from $\rho[Z^{(n)}]$).
\item $S(BC) = S_{n-1}$ by Lemma~\ref{lem:marginal}
  (tracing out $A = \{X_1\}$, the first measurement, from $\rho[Z^{(n)}]$,
  which yields $\rho[Z^{(n-1)}]$ exactly, as the cyclicity argument in the proof
  shows the matrix entries are unchanged under the time shift).
\item $S(B) = S_{n-2}$ (applying Lemma~\ref{lem:marginal_last} to $\rho[Z^{(n-1)}]$,
  or equivalently Lemma~\ref{lem:marginal} to the $BC$-marginal, both give
  $\rho[Z^{(n-2)}]$).
\end{itemize}
Substituting into~\eqref{eq:ssa}:
\[
  S_{n-1} + S_{n-1} \;\ge\; S_{n-2} + S_n,
\]
which is $2S_{n-1}\ge S_n+S_{n-2}$, i.e.\ $\Delta S_n \le \Delta S_{n-1}$.\qed
\end{proof}

\begin{remark}
The proof uses two lemmas, both proved using only OPU completeness and unitarity:
Lemma~\ref{lem:marginal_last} (tracing out the last measurement gives $\rho[Z^{(n-1)}]$)
for the $S(AB)$ term, and Lemma~\ref{lem:marginal} (tracing out the first measurement)
for $S(BC)$.  The fact that the first-marginal gives $\rho[Z^{(n-1)}]$ \emph{exactly}
as a matrix—not merely with the same spectrum—means that the proof does not require a
separate time-stationarity assumption: time-homogeneity is already encoded in the
cyclic-trace identity $\operatorname{Tr}(C_J^\dagger C_I) = \operatorname{Tr}(Z^{(n-1)\dagger}_J Z^{(n-1)}_I)$.
\end{remark}

\subsection{The Operator-Entanglement Pesin Inequality}

\begin{theorem}[Operator-Entanglement Pesin Inequality]\label{thm:oe_pesin}
For the kicked Ising model with $X$-basis OPU on site $0$, open boundary
conditions, and tracial state $\omega = I/D$:
\begin{equation}\label{eq:oe_pesin}
  \boxed{h_\mathrm{AFL}^\mathrm{time}(J,g)
  \;\le\;
  E_\mathrm{op}(u_J) \;=\; H_\mathrm{bin}(\sin^2\!J)}
\end{equation}
for all $(J,g)\in[0,\pi/2]^2$.

Combined with the complementarity law (Theorem~\ref{thm:complementarity}):
\begin{equation}\label{eq:pesin_comp}
  h_\mathrm{AFL}^\mathrm{time}(J,g) + I_\mathrm{temp}(J) \;\le\; \log d.
\end{equation}
\end{theorem}

\begin{proof}
By Theorem~\ref{thm:ssa_concave}, $\Delta S_n$ is non-increasing, so
\[
  h_\mathrm{AFL}^\mathrm{time}
  = \lim_{n\to\infty}\Delta S_n
  \;\le\;
  \Delta S_2 = S_2 - S_1 = E_\mathrm{op}(u_J),
\]
where the last equality is equation~\eqref{eq:ds2_eop} (proved in
Theorem~\ref{thm:complementarity}).
Inequality~\eqref{eq:pesin_comp} then follows from $E_\mathrm{op}+I_\mathrm{temp}=\log d$.\qed
\end{proof}

\begin{remark}[Physical interpretation]
The operator entanglement $E_\mathrm{op}(u_J)$ of the local gate $u_J$
sets an \emph{upper bound} on how fast the AFL time entropy can grow.
Increasing the ZZ coupling $J$ increases $E_\mathrm{op}$ (more gate entanglement)
and allows faster entropy production.  The $X$-kick parameter $g$ only
reduces $h_\mathrm{AFL}^\mathrm{time}$ below $E_\mathrm{op}$, consistent
with~\eqref{eq:oe_pesin}.  Inequality~\eqref{eq:pesin_comp} has the flavour
of a quantum Pesin theorem: temporal information production plus temporal
memory is conserved at most at $\log d$ per step.
\end{remark}

\subsection{Equality Conditions and Numerical Verification}

\begin{proposition}[Equality at Dual-Unitary]\label{prop:du_equality}
At $J=g=\pi/4$ (dual-unitary coupling), $\Delta S_n = E_\mathrm{op}=\log d$
for all $n\ge 1$, so $h_\mathrm{AFL}^\mathrm{time}=E_\mathrm{op}=\log d$.
\end{proposition}

\begin{proof}
At $J=g=\pi/4$, $\rho[Z^{(n)}]=I/2^n$ (flat spectrum, proved in subfolder~1,
Section~62), so $S_n = n\log d$ and $\Delta S_n = \log d = E_\mathrm{op}$ for all $n$.\qed
\end{proof}

\begin{proposition}[Markov Equality at $g=0$, Thermodynamic Limit]\label{prop:g0_equality}
For pure ZZ dynamics ($g=0$) with $X$-basis OPU on site~$0$ and open BC:
the one-step channel on site~$0$ is
\[
  \mathcal{E}_{D_{ZZ}}(P_i^x) = \cos^2\!J\,P_i^x + \sin^2\!J\,P_{1-i}^x.
\]
In the thermodynamic limit $L\to\infty$ the orbit entropy is Markovian with
transition matrix $T=\bigl[\begin{smallmatrix}\cos^2\!J&\sin^2\!J\\\sin^2\!J&\cos^2\!J\end{smallmatrix}\bigr]$
and entropy rate $h_\mathrm{KS}(T) = H_\mathrm{bin}(\sin^2\!J) = E_\mathrm{op}$.
Hence $h_\mathrm{AFL}^\mathrm{time}(g=0,L\to\infty) = E_\mathrm{op}$.
\end{proposition}

\noindent
\textbf{Numerical verification} ($L=4$, $d=2$, open BC, $X$-basis OPU, $n_\mathrm{max}=5$):

\medskip
\noindent\textit{Table~1: Complementarity law $\Delta S_2 = E_\mathrm{op}$ confirmed for all $(J,g)$.}

\begin{center}
\begin{tabular}{@{}ccccccc@{}}
\toprule
$J$ & $g$ & $S_1$ & $S_2$ & $\Delta S_2$ & $E_\mathrm{op}$ & Match?\\
\midrule
0.10 & 0.00 & 0.693 & 0.749 & 0.056 & 0.056 & \checkmark\\
$\pi/8$ & 0.00 & 0.693 & 1.110 & 0.417 & 0.417 & \checkmark\\
$\pi/6$ & 0.00 & 0.693 & 1.255 & 0.562 & 0.562 & \checkmark\\
$\pi/4$ & 0.00 & 0.693 & 1.386 & 0.693 & 0.693 & \checkmark\\
$\pi/8$ & $\pi/8$ & 0.693 & 1.110 & 0.417 & 0.417 & \checkmark\\
$\pi/4$ & $\pi/4$ & 0.693 & 1.386 & 0.693 & 0.693 & \checkmark\\
\bottomrule
\end{tabular}
\end{center}

\medskip
\noindent\textit{Table~2: Non-increasing $\Delta S_n$ (concavity) confirmed.}

\begin{center}
\begin{tabular}{@{}ccccccc@{}}
\toprule
$J$ & $g$ & $\Delta S_2$ & $\Delta S_3$ & $\Delta S_4$ & $\Delta S_5$ & Concave?\\
\midrule
0.10 & 0.00 & 0.056 & 0.041 & 0.035 & 0.031 & \checkmark\\
$\pi/8$ & 0.00 & 0.417 & 0.146 & 0.067 & 0.032 & \checkmark\\
$\pi/8$ & $\pi/8$ & 0.417 & 0.305 & 0.290 & 0.283 & \checkmark\\
$\pi/6$ & $\pi/8$ & 0.562 & 0.374 & 0.362 & 0.346 & \checkmark\\
$\pi/4$ & $\pi/8$ & 0.693 & 0.417 & 0.417 & 0.305 & \checkmark\\
$\pi/4$ & $\pi/4$ & 0.693 & 0.693 & 0.693 & 0.693 & \checkmark (const.)\\
\bottomrule
\end{tabular}
\end{center}

\medskip
\noindent\textit{Table~3: Gap function $\Delta(J,g) = E_\mathrm{op}(J) - h_\mathrm{AFL}^\mathrm{time}(J,g)$ [nats]
on a $5\times 5$ grid.}

\begin{center}
\begin{tabular}{@{}cccccc@{}}
\toprule
$g\backslash J$ & 0.05 & $\pi/8$ & $\pi/6$ & $3\pi/16$ & $\pi/4$\\
\midrule
0 & 0.006 & 0.384 & 0.556 & 0.617 & 0.693 \\
$\pi/16$ & 0.003 & 0.247 & 0.402 & 0.472 & 0.579\\
$\pi/8$ & 0.003 & 0.134 & 0.216 & 0.271 & 0.388\\
$3\pi/16$ & 0.004 & 0.112 & 0.116 & 0.119 & 0.136\\
$\pi/4$ & 0.003 & 0.111 & 0.089 & 0.061 & $\approx 0$ \\
\bottomrule
\end{tabular}
\end{center}

\noindent All entries $\ge 0$.  The gap vanishes (to numerical precision) only at
$J=g=\pi/4$ (dual-unitary corner).  The gap is largest at $g=0$ for fixed $J$
(pure ZZ dynamics has maximum correlation structure; finite-size effects prevent
exact equality on this line).

\medskip
\noindent\textit{Table~6: Two-sided marginal consistency verified numerically
($L=4$, open BC, $X$-basis OPU).  Entries show the max-norm error
$\|\cdot\|_\infty$ between the computed marginal and $\rho[Z^{(n-1)}]$;
all errors are at machine precision ($\lesssim 2\times 10^{-16}$).}

\begin{center}
\begin{tabular}{@{}lcccc@{}}
\toprule
System & $n$ &
  $\|\mathrm{Tr}_{\rm first}[\rho[Z^{(n)}]]-\rho[Z^{(n-1)}]\|_\infty$ &
  $\|\mathrm{Tr}_{\rm last}[\rho[Z^{(n)}]]-\rho[Z^{(n-1)}]\|_\infty$ &
  $S_{AB}=S_{BC}$?\\
\midrule
$J=\pi/8,\,g=0$ (integrable) & 2 & $1.2\times10^{-17}$ & $<10^{-17}$ & \checkmark\\
                               & 3 & $1.4\times10^{-17}$ & $1.4\times10^{-17}$ & \checkmark\\
                               & 4 & $6.9\times10^{-18}$ & $6.9\times10^{-18}$ & \checkmark\\
\midrule
$J=\pi/8,\,g=\pi/8$           & 2 & $2.7\times10^{-17}$ & $<10^{-17}$ & \checkmark\\
                               & 3 & $5.6\times10^{-17}$ & $5.6\times10^{-17}$ & \checkmark\\
                               & 4 & $5.6\times10^{-17}$ & $5.6\times10^{-17}$ & \checkmark\\
\midrule
$J=\pi/4,\,g=\pi/4$ (DU)      & 2 & $5.6\times10^{-17}$ & $5.6\times10^{-17}$ & \checkmark\\
                               & 3 & $8.3\times10^{-17}$ & $5.6\times10^{-17}$ & \checkmark\\
                               & 4 & $3.0\times10^{-17}$ & $2.8\times10^{-17}$ & \checkmark\\
\midrule
$J=0.10,\,g=0.50$ (arbitrary) & 2 & $1.1\times10^{-16}$ & $1.1\times10^{-16}$ & \checkmark\\
                               & 3 & $1.7\times10^{-16}$ & $1.7\times10^{-16}$ & \checkmark\\
                               & 4 & $1.1\times10^{-16}$ & $1.1\times10^{-16}$ & \checkmark\\
\bottomrule
\end{tabular}
\end{center}

\noindent The $S_{AB}=S_{BC}$ column confirms that the first-marginal gives $\rho[Z^{(n-1)}]$
exactly as a matrix (numerical agreement to $<10^{-8}$), so the spectrum—and hence
the entropy—is unchanged.  Together with Table~6, the SSA proof in
Theorem~\ref{thm:ssa_concave} is verified to be fully rigorous.

\begin{corollary}[Pesin-Complementarity Inequality]\label{cor:pesin_comp}
For the kicked Ising model with $X$-basis OPU:
\begin{equation}
  h_\mathrm{AFL}^\mathrm{time}(J,g)
  \;+\; I_\mathrm{temp}(J)
  \;\le\; \log d,
\end{equation}
with equality at $J=g=\pi/4$.  The triple
$(h_\mathrm{AFL}^\mathrm{time},\,E_\mathrm{op},\,I_\mathrm{temp})$
satisfies
$h_\mathrm{AFL}^\mathrm{time}\le E_\mathrm{op}=\log d - I_\mathrm{temp}$,
forming a complete quantum information budget for each time step.
\end{corollary}

\section{Near-Dual-Unitary Geometry of the Pesin Gap}
\label{sec:near_du}
% ============================================================

In Section~\ref{sec:pesin_ineq} we proved the upper bound
$h_\mathrm{AFL}^\mathrm{time}(J,g)\le E_\mathrm{op}(J)$ with equality at
$J=g=\pi/4$.  We now characterise the \emph{shape} of the gap function
$\Delta(J,g)=E_\mathrm{op}(J)-h_\mathrm{AFL}^\mathrm{time}(J,g)$ near the
dual-unitary point and derive an explicit near-DU lower bound.

\subsection{Taylor Expansion of $E_\mathrm{op}$ Near Dual-Unitary}

Set $\delta J = \pi/4 - J$ and $\delta g = \pi/4 - g$.

\begin{lemma}[Second-order expansion of $E_\mathrm{op}$]\label{lem:eop_expand}
\begin{equation}
  E_\mathrm{op}(\pi/4-\delta J)
  = \ln 2 - 2(\delta J)^2 + O((\delta J)^4).
\end{equation}
\end{lemma}

\begin{proof}
$E_\mathrm{op}(J)=H_\mathrm{bin}(\sin^2 J)$.  At $J=\pi/4$, $\sin^2 J=\tfrac{1}{2}$
so $H_\mathrm{bin}(\tfrac{1}{2})=\ln 2$.  Since
\[
  \frac{d}{dJ}H_\mathrm{bin}(\sin^2 J)\Big|_{\pi/4}
  = -\ln\!\tfrac{\sin^2 J}{1-\sin^2 J}\cdot\sin(2J)\Big|_{\pi/4} = 0,
\]
and
\[
  \frac{d^2}{dJ^2}H_\mathrm{bin}(\sin^2 J)\Big|_{\pi/4}
  = -\frac{1}{\sin^2 J(1-\sin^2 J)}\sin^2(2J)\Big|_{\pi/4} = -4,
\]
the Taylor expansion gives $E_\mathrm{op}(\pi/4-\delta J)=\ln 2 - 2(\delta J)^2+O((\delta J)^4)$.
\qed
\end{proof}

\noindent Numerically: the ratio $(\ln 2-E_\mathrm{op})/(\delta J)^2 \to 2$ as $\delta J\to 0$,
confirmed in Table~4 below.

\subsection{Isotropy of the Gap Near Dual-Unitary}

Let $r = \sqrt{(\delta J)^2+(\delta g)^2}$ denote the Euclidean distance from
$(J,g)=(\pi/4,\pi/4)$ in parameter space.

\begin{proposition}[Isotropic near-DU gap]\label{prop:isotropic_gap}
For the kicked Ising model with $X$-basis OPU, $L=4$, open BC, $n_\mathrm{max}=6$:
\begin{equation}\label{eq:isotropic}
  \Delta(J,g) = E_\mathrm{op}(J) - h_\mathrm{AFL}^\mathrm{time}(J,g)
  \;\approx\; 4\,r^2 + O(r^4)
\end{equation}
for all directions from the dual-unitary point, and consequently
\begin{equation}\label{eq:ratio_expand}
  \frac{h_\mathrm{AFL}^\mathrm{time}(J,g)}{E_\mathrm{op}(J)}
  \;\approx\; 1 - \frac{4\,r^2}{\ln 2} + O(r^4).
\end{equation}
\end{proposition}

\noindent\textit{Justification.}
Along the $J$-direction ($\delta g=0$):
$E_\mathrm{op}\approx\ln 2-2(\delta J)^2$ and $\Delta\approx 4(\delta J)^2$,
giving $h\approx\ln 2-6(\delta J)^2$.
Along the $g$-direction ($\delta J=0$):
$E_\mathrm{op}=\ln 2$ (unchanged) and $\Delta\approx 4(\delta g)^2$,
giving $h\approx\ln 2-4(\delta g)^2$.
In both cases $h/E_\mathrm{op}\approx 1-4r^2/\ln 2$, confirming~\eqref{eq:ratio_expand}.

\begin{corollary}[Near-DU lower bound]\label{cor:near_du_lb}
For $(J,g)$ satisfying $r:=\sqrt{(\pi/4-J)^2+(\pi/4-g)^2}\le r_0$, we have
\begin{equation}\label{eq:near_du_lb}
  h_\mathrm{AFL}^\mathrm{time}(J,g)
  \;\ge\;
  E_\mathrm{op}(J)\!\left(1 - \frac{4\,r^2}{\ln 2}\right) - O(r^4).
\end{equation}
In particular, $h_\mathrm{AFL}^\mathrm{time}\ge\tfrac{1}{2}E_\mathrm{op}$
whenever $r\le\sqrt{\ln 2/8}\approx 0.294$.
\end{corollary}

\subsection{Numerical Evidence}

\noindent\textit{Table~4: Taylor expansion of $E_\mathrm{op}$ near $J=\pi/4$ (coefficient $=2$).}

\begin{center}
\begin{tabular}{@{}ccccc@{}}
\toprule
$\delta J$ & $E_\mathrm{op}(\pi/4-\delta J)$ & $\ln 2 - E_\mathrm{op}$ & $(\delta J)^2$ & $\frac{\ln2-E_\mathrm{op}}{(\delta J)^2}$\\
\midrule
0.01 & 0.692947 & 0.000200 & 0.0001 & 1.996 \\
0.05 & 0.688156 & 0.004992 & 0.0025 & 1.997 \\
0.10 & 0.673281 & 0.019867 & 0.0100 & 1.987 \\
0.15 & 0.648822 & 0.044325 & 0.0225 & 1.970 \\
0.20 & 0.615281 & 0.077866 & 0.0400 & 1.947 \\
\bottomrule
\end{tabular}
\end{center}

\medskip
\noindent\textit{Table~5: Ratio $h_\mathrm{AFL}^\mathrm{time}/E_\mathrm{op}$ at distance $r$ from DU, three directions.}

\begin{center}
\begin{tabular}{@{}ccccc@{}}
\toprule
$r$ & $J$-dir ($\delta g=0$) & diagonal ($\delta J=\delta g$) & $g$-dir ($\delta J=0$) & predicted $1-4r^2/\ln2$\\
\midrule
0.05 & 0.9859 & 0.9860 & 0.9857 & 0.9855 \\
0.10 & 0.9470 & 0.9488 & 0.9442 & 0.9423 \\
0.15 & 0.8924 & 0.9002 & 0.8794 & 0.8702 \\
0.20 & 0.8341 & 0.8525 & 0.7974 & 0.7693 \\
0.25 & 0.7822 & 0.8129 & 0.7051 & 0.6395 \\
\bottomrule
\end{tabular}
\end{center}

\noindent The predicted formula $1-4r^2/\ln 2$ matches all three directions
well for $r\lesssim 0.15$; higher-order corrections become visible at $r\ge 0.20$.
The $g$-direction departs fastest (coefficient effectively larger for larger $r$),
while the $J$-direction stays close to the prediction longest (because $E_\mathrm{op}$
also decreases in $J$, reducing the ratio's denominator and partially compensating).

\begin{remark}[Physical interpretation]
The isotropy of the near-DU gap in parameter space reflects the fact that at the
dual-unitary point the gate is at a \emph{saddle point} of the entropy-production
rate: both increasing or decreasing $J$ from $\pi/4$, and both increasing or
decreasing $g$ from $\pi/4$, reduce $h_\mathrm{AFL}^\mathrm{time}$ below its
maximum $E_\mathrm{op}$.  The common coefficient $4/\ln 2$ in the ratio expansion
is a universal feature of the X-basis OPU for qubits ($d=2$).
\end{remark}


\section{System-Size Independence and Non-Commutativity of Limits}
\label{sec:L_independence}
% ============================================================

We now prove two structural theorems about the time-AFL orbit that clarify the
relationship between the finite-chain Pesin inequality and the thermodynamic limit.

\subsection{L-Independence at n=2 for All Couplings}

The complementarity law $\Delta S_2 = E_\mathrm{op}(J)$ (proved in
Theorem~\ref{thm:complementarity}) implies that $S_2-S_1$ is independent of both
the system size $L$ and the transverse coupling $g$.  We now prove the stronger
statement that the \emph{full density matrix} $\rho[Z^{(2)}]$ is $L$-independent.

\begin{theorem}[L-Independence at $n=2$]\label{thm:L_indep_n2}
For the kicked Ising model with any $(J,g)$, $X$-basis OPU on site $L-1$ (the
rightmost site), open boundary conditions, and $L\ge 2$:
\begin{equation}
  \rho[Z^{(2)}](L) \;=\; \rho[Z^{(2)}](2)
\end{equation}
as $2^2\times 2^2$ matrices.  In particular, $S_2(J,g,L) = S_2(J,g,2)$ and
$\Delta S_2(J,g,L) = E_\mathrm{op}(J)$ for all $L\ge 2$ and all $g$.
\end{theorem}

\begin{proof}
Write $U_L = V'\,U_{L-2,L-1}\,K$ where
\[
  V' \;=\;\prod_{k=0}^{L-3} e^{-\mathrm{i}J\sigma^z_k\sigma^z_{k+1}},
  \quad
  K \;=\;\prod_{k=0}^{L-1} e^{-\mathrm{i}g\sigma^x_k}.
\]
$V'$ acts on sites $0,\ldots,L-2$ (not on the measured site $L-1$).  Since the
X-projectors $P_i = (I\pm\sigma^x_{L-1})/2$ are functions of $\sigma^x_{L-1}$,
they commute with $V'$ (different sites) and with every factor of $K$ (both are
functions of $\sigma^x$, hence $[P_i, K] = 0$).

Therefore $P_i\,K^\dagger = K^\dagger P_i$ and $P_i\,{V'}^\dagger = {V'}^\dagger P_i$, so
\begin{align}
  Z^{(2),(L)}_{(i_1,i_2)}
  &= P_{i_1}\,U_L^\dagger\,P_{i_2}\,U_L \notag\\
  &= P_{i_1}\,(K^\dagger\,U_{L-2,L-1}^\dagger\,{V'}^\dagger)\,P_{i_2}\,(V'\,U_{L-2,L-1}\,K) \notag\\
  &= K^\dagger\,{V'}^\dagger\!\underbrace{P_{i_1}\,U_{L-2,L-1}^\dagger\,P_{i_2}\,U_{L-2,L-1}}_{=:\;Z^{(2),(2)}_{(i_1,i_2)}}\!V'\,K.
\end{align}
Setting $W := V'\,K$, we have $Z^{(2),(L)}_I = W^\dagger Z^{(2),(2)}_I W$ and
\begin{align}
  \operatorname{Tr}_{L\,\text{sites}}\!\bigl(Z^{(2),(L)\dagger}_J Z^{(2),(L)}_I\bigr)
  &= \operatorname{Tr}\!\bigl(W^\dagger Z^{(2),(2)\dagger}_J W\,W^\dagger Z^{(2),(2)}_I W\bigr) \notag\\
  &= \operatorname{Tr}\!\bigl(Z^{(2),(2)\dagger}_J Z^{(2),(2)}_I\bigr),
\end{align}
using $W\,W^\dagger = I$ and cyclicity of trace.  Since $Z^{(2),(2)}$ acts only on
sites $L-2,L-1$ (and $\operatorname{Tr}$ over sites $0,\ldots,L-3$ gives $2^{L-2}$),
\[
  \rho[Z^{(2)}]^{(L)}_{I,J}
  = \frac{1}{2^L}\operatorname{Tr}_L(\cdots)
  = \frac{2^{L-2}}{2^L}\operatorname{Tr}_{L-2,L-1}(Z^{(2),(2)\dagger}_J Z^{(2),(2)}_I)
  = \frac{1}{4}\operatorname{Tr}_2(\cdots)
  = \rho[Z^{(2)}]^{(2)}_{I,J}.\qed
\]
\end{proof}

\begin{remark}
The proof uses only two facts: (i)~$P_i$ commutes with both $V'$ (different sites)
and $K$ (both are functions of $\sigma^x$); (ii)~cyclicity of trace.
No property of the kicked Ising model beyond the structure of the OPU is used.
The theorem therefore holds for any quantum spin chain with a single-site OPU and
open boundary conditions.
\end{remark}

\subsection{Universal L-Independence for $g=0$}

For pure ZZ dynamics ($g=0$), the X-kick is absent ($K=I$) and the V-conjugation
extends to all orbit lengths.

\begin{theorem}[Universal L-Independence for $g=0$]\label{thm:L_indep_g0}
For $g=0$ and any $n\ge 1$ and $L\ge 2$:
\begin{equation}
  \rho[Z^{(n)}](L) \;=\; \rho[Z^{(n)}](2).
\end{equation}
In particular, $S_n(J,0,L)$ is $L$-independent.
\end{theorem}

\begin{proof}
For $g=0$: $U_L = V'\,U_{L-2,L-1}$ (no kick).  By induction on $n$, each
occurrence of $U_L^\dagger$ in $Z^{(n),(L)}_I = P_{i_1}U_L^\dagger\cdots P_{i_n}U_L^{n-1}$
contributes one factor of ${V'}^\dagger$, and each of the $n-1$ factors in $U_L^{n-1}$
contributes one factor of $V'$.  Since $P_i$ commutes with $V'$ and ${V'}^\dagger$,
all $n-1$ conjugating pairs ${V'}^\dagger\!\cdots V'$ can be moved to the exterior:
\[
  Z^{(n),(L)}_I = {V'}^\dagger Z^{(n),(L-1)}_I\,V',
\]
where $Z^{(n),(L-1)}_I$ is the $(L-1)$-site operator. Iterating and applying
cyclicity gives
$\operatorname{Tr}_{L}(Z^{(n),(L)\dagger}_J Z^{(n),(L)}_I)
= 2^{L-2}\operatorname{Tr}_2(Z^{(n),(2)\dagger}_J Z^{(n),(2)}_I)$,
and dividing by $D=2^L$ gives the result.\qed
\end{proof}

\subsection{Non-Commutativity of the Thermodynamic and Long-Time Limits}

Theorem~\ref{thm:L_indep_g0} has a striking corollary that resolves an apparent
tension between Proposition~\ref{prop:g0_equality} (which states
$h_\mathrm{AFL}^\mathrm{time}(g=0,L\to\infty) = E_\mathrm{op}$) and
Section~\ref{sec:time_evo} (which proves $h_\mathrm{AFL}^\mathrm{time}(L) = 0$ for
any finite $L$).

\begin{corollary}[Non-Commutativity of Limits]\label{cor:noncomm}
For $g=0$ and any $J$, the limits $n\to\infty$ and $L\to\infty$ do not commute:
\begin{align}
  \lim_{L\to\infty}\lim_{n\to\infty}\Delta S_n(J,0,L) &= E_\mathrm{op}(J)
    \quad\text{(Markovian limit, Prop.~\ref{prop:g0_equality})},\\
  \lim_{n\to\infty}\lim_{L\,\text{finite}}\Delta S_n(J,0,L) &= 0
    \quad\text{(rank saturation, Section~\ref{sec:time_evo})}.
\end{align}
\end{corollary}

\begin{proof}
The first limit holds because the infinite-chain orbit is well approximated by the
Markov chain $T = \bigl[\begin{smallmatrix}\cos^2\!J&\sin^2\!J\\\sin^2\!J&\cos^2\!J\end{smallmatrix}\bigr]$,
whose entropy rate is $H_\mathrm{bin}(\sin^2\!J) = E_\mathrm{op}$.
The second limit holds because for any fixed finite $L$, the operators $Z^{(n)}_I$
live in the $2^L$-dimensional Hilbert space; for $n\gg L$ they become linearly
dependent and $\Delta S_n\to 0$ (Section~\ref{sec:time_evo}).
By Theorem~\ref{thm:L_indep_g0}, $\Delta S_n(J,0,L)$ is the same for all finite
$L\ge 2$, so the second limit is unambiguous.\qed
\end{proof}

\begin{remark}
Corollary~\ref{cor:noncomm} shows that the Markovian entropy rate $E_\mathrm{op}$
is a genuine thermodynamic-limit phenomenon: it cannot be seen at any finite chain
length.  The complementarity equality $\Delta S_2 = E_\mathrm{op}$ is the only
\emph{finite-$L$} signature of this Markovian rate — it holds at $n=2$ for all
$L$ by Theorem~\ref{thm:L_indep_n2} and the complementarity law.
\end{remark}

\noindent\textit{Table~7: L-independence of $\rho[Z^{(n)}]$ verified numerically.
Left panel: $n=2$ for all $(J,g)$ (Theorem~\ref{thm:L_indep_n2}).
Right panel: all $n\le 4$ for $g=0$ (Theorem~\ref{thm:L_indep_g0}).
Entries show $\|\rho[Z^{(n)}](L)-\rho[Z^{(n)}](2)\|_\infty$; all are at machine
precision ($\lesssim 2\times10^{-16}$).}

\begin{center}
\begin{tabular}{@{}lccc@{\qquad}lccc@{}}
\toprule
\multicolumn{4}{c}{$n=2$, varying $(J,g)$ and $L$} &
\multicolumn{4}{c}{$g=0$, $J=\pi/8$, varying $n$ and $L$}\\
\cmidrule(r){1-4}\cmidrule(l){5-8}
$(J,g)$ & $L=3$ & $L=4$ & $L=5$ &
$n$ & $L=3$ & $L=4$ & $L=5$\\
\midrule
$(\pi/8,\,0)$ & $3\times10^{-17}$ & $3\times10^{-17}$ & $3\times10^{-17}$ &
$1$ & $0$ & $0$ & $0$\\
$(\pi/8,\,\pi/8)$ & $1\times10^{-16}$ & $6\times10^{-17}$ & $8\times10^{-17}$ &
$2$ & $4\times10^{-16}$ & $4\times10^{-16}$ & $4\times10^{-16}$\\
$(\pi/4,\,\pi/4)$ & $2\times10^{-16}$ & $6\times10^{-17}$ & $6\times10^{-17}$ &
$3$ & $4\times10^{-16}$ & $4\times10^{-16}$ & $4\times10^{-16}$\\
$(0.30,\,0.70)$ & $1\times10^{-16}$ & $6\times10^{-17}$ & $6\times10^{-17}$ &
$4$ & $4\times10^{-16}$ & $4\times10^{-16}$ & $4\times10^{-16}$\\
\bottomrule
\end{tabular}
\end{center}

\noindent For $g>0$ and $n\ge 4$, L-independence can break down once the
2-site orbit saturates in rank.  Specifically, at $J=\pi/4$, $g=\pi/4-0.1$: we find
$\|\rho[Z^{(4)}](L=3)-\rho[Z^{(4)}](L=2)\|_\infty = 0.65$ while
$\|\rho[Z^{(3)}](L=3)-\rho[Z^{(3)}](L=2)\|_\infty < 10^{-14}$.
This is consistent with the 2-site orbit saturating in rank at $n\approx 3$--$4$;
for larger $L$, the orbit keeps growing until the larger Hilbert space is exhausted.

\section{OTOC and the AFL Entropy: A One-Step Connection}
\label{sec:otoc_afl}
% ============================================================

We now establish a direct, rigorous connection between the one-step
out-of-time-order correlator (OTOC) of the kicked Ising model and the AFL
entropy increment $\Delta S_2 = E_\mathrm{op}$ proved in
Theorem~\ref{thm:complementarity}.

\subsection{OTOC Definition and the Key Formula}

Fix the tracial state $\omega = I/D$ and the X-basis OPU projector
$W = P_0 = (I+\sigma^x_{L-1})/2$ at site $L-1$.  Define the one-step
\emph{Loschmidt-echo OTOC}:
\begin{equation}\label{eq:otoc_def}
  F(n) \;:=\;
  \frac{1}{D^2}\bigl|\operatorname{Tr}\bigl[P_0\,(U^n P_0 U^{-n})\bigr]\bigr|^2,
  \qquad
  F(0) \;=\; \frac{1}{4}.
\end{equation}

\begin{proposition}[OTOC at one step]\label{prop:otoc1}
For the kicked Ising model with any $(J,g)$ and any $L\ge 2$:
\begin{equation}\label{eq:otoc1}
  \frac{F(1)}{F(0)} \;=\; \cos^4\!J.
\end{equation}
In particular, $F(1)/F(0)$ is independent of $g$ and of $L$.
\end{proposition}

\begin{proof}
Since $P_0 = (I\pm\sigma^x_{L-1})/2$ is a function of $\sigma^x_{L-1}$, it commutes
with the X-kick $K = \prod_k e^{-\mathrm{i}g\sigma^x_k}$ (each factor either acts
on a different site or is also a function of $\sigma^x_{L-1}$, hence commutes).
Therefore
\[
  U P_0 U^{-1}
  = e^{-\mathrm{i}J H_{ZZ}}\,K\,P_0\,K^{-1}\,e^{\mathrm{i}J H_{ZZ}}
  = e^{-\mathrm{i}J H_{ZZ}}\,P_0\,e^{\mathrm{i}J H_{ZZ}}.
\]
Only the ZZ bond at sites $L-2$, $L-1$ acts non-trivially on $P_0$ (the other
ZZ terms commute with $P_0$), so
$e^{-\mathrm{i}J H_{ZZ}} P_0 e^{\mathrm{i}J H_{ZZ}} = e^{-\mathrm{i}J\sigma^z_{L-2}\sigma^z_{L-1}} P_0 e^{\mathrm{i}J\sigma^z_{L-2}\sigma^z_{L-1}}$,
and the standard BCH identity gives
\[
  P_0(1) := U P_0 U^{-1}
  = \tfrac{1}{2}\bigl(I + \cos(2J)\,\sigma^x_{L-1}
                      + \sin(2J)\,\sigma^y_{L-1}\sigma^z_{L-2}\bigr).
\]
In the tracial state, $\operatorname{Tr}[\sigma^y_{L-1}\sigma^z_{L-2}] = 0$ and
$\operatorname{Tr}[P_0\,\sigma^y_{L-1}\sigma^z_{L-2}] = 0$, so
\[
  \operatorname{Tr}[P_0\,P_0(1)]
  = \tfrac{1}{4}\bigl(\operatorname{Tr}[I] + \cos(2J)\operatorname{Tr}[\sigma^x_{L-1}+(\sigma^x_{L-1})^2]\bigr)
  = \tfrac{D}{2}\cos^2\!J,
\]
(using $\operatorname{Tr}[\sigma^x] = 0$ and $\operatorname{Tr}[(\sigma^x)^2]=D$, the factor
$2^{L-2}$ from tracing over the other sites cancels with the normalisation).
Hence
$F(1) = (D/2)^2\cos^4\!J\,/D^2 = \cos^4\!J/4 = F(0)\cos^4\!J$.\qed
\end{proof}

\subsection{The OTOC–AFL Connection}

\begin{theorem}[OTOC–AFL entropy connection]\label{thm:otoc_afl}
For any $(J,g)$ and $L\ge 2$:
\begin{equation}\label{eq:otoc_afl}
  \boxed{\Delta S_2(J,g)
  \;=\; E_\mathrm{op}(J)
  \;=\; H_\mathrm{bin}\!\Bigl(1 - \sqrt{\frac{F(1)}{F(0)}}\Bigr).}
\end{equation}
\end{theorem}

\begin{proof}
By Theorem~\ref{thm:complementarity}: $\Delta S_2 = E_\mathrm{op}(J) = H_\mathrm{bin}(\sin^2\!J)$.
By Proposition~\ref{prop:otoc1}: $F(1)/F(0) = \cos^4\!J$, so $\sqrt{F(1)/F(0)} = \cos^2\!J$
and $1 - \sqrt{F(1)/F(0)} = \sin^2\!J$.
Substituting: $H_\mathrm{bin}(1-\sqrt{F(1)/F(0)}) = H_\mathrm{bin}(\sin^2\!J) = \Delta S_2$.\qed
\end{proof}

\begin{remark}
Equation~\eqref{eq:otoc_afl} gives a \emph{physical interpretation} of the AFL
entropy increment $\Delta S_2$: it is the binary entropy of the ``1-step
scrambling probability'' $\sin^2 J = 1 - \sqrt{F(1)/F(0)}$.  The $g$-independence
of $F(1)/F(0)$ (via $P_0$ commuting with $K$) is the same algebraic reason for the
$g$-independence of $E_\mathrm{op}$ proved in Section~\ref{sec:eop_formula}.
\end{remark}

\subsection{Dual-Unitary Instantaneous Scrambling}

At the dual-unitary point $J=g=\pi/4$, $\cos^4\!J = 1/4$, so
$F(1)/F(0) = 1/4$ and $\Delta S_2 = \ln 2$.  Numerically one finds that for
all $n\ge 1$ (up to the Poincaré recurrence time $n_\mathrm{rec} = 2L$):
\begin{equation}\label{eq:du_otoc}
  \frac{F(n)}{F(0)} \;\equiv\; \frac{1}{4} \qquad (n=1,\ldots,2L-1),\quad
  \frac{F(2L)}{F(0)} = 1,
\end{equation}
reflecting the fact that the operator $P_0$ spreads to its \emph{maximally mixed}
Heisenberg-evolved form in a single Floquet step and then undergoes exact
recurrence at period $2L$.  This single-step scrambling is the OTOC signature of
the dual-unitary circuit.

\noindent\textit{Table~8: OTOC values $F(n)/F(0)$ and operator spreading
$\|P_0(n)-P_0\|_F / \|P_0\|_F$ for $L=6$.}

\begin{center}
\begin{tabular}{@{}lcccccc@{}}
\toprule
System & $n=1$ & $n=2$ & $n=4$ & $n=6$ & $n=8$ & $n=12$\\
\midrule
$J=\pi/8,\,g=\pi/8$ & 0.729 & 0.391 & 0.264 & 0.254 & 0.252 & 0.251\\
$J=\pi/4,\,g=\pi/8$ & 0.250 & 0.063 & 0.316 & 0.316 & 0.223 & 0.276\\
$J=\pi/4,\,g=\pi/4$ (DU) & 0.250 & 0.250 & 0.250 & 0.250 & 0.250 & 1.000\\
\midrule
\multicolumn{7}{@{}l}{Predicted $F(1)/F(0) = \cos^4\!J$:
  $J=\pi/8$: 0.729; $J=\pi/4$: 0.250.  All confirmed.}\\
\bottomrule
\end{tabular}
\end{center}

\noindent For the DU circuit: $F(n)/F(0) = 0.250$ for $n=1,\ldots,11$ and
$F(12)/F(0) = 1.000$ (Poincaré recurrence at $n=2L=12$ for $L=6$).
The operator spreading (not shown) saturates at $\|P_0(1)-P_0\|_F/\|P_0\|_F = \sqrt{2}$
for DU in a single step, while reaching the same value only gradually for other
$(J,g)$, consistent with the instantaneous scrambling~\eqref{eq:du_otoc}.

\section{Rényi-$\alpha$ AFL Entropy and a Universal Complementarity Law}
\label{sec:renyi_afl}
% ============================================================

We extend the complementarity law $\Delta S_2 = E_\mathrm{op}(J)$
(Theorem~\ref{thm:complementarity}) to all Rényi orders $\alpha > 0$.

\subsection{Eigenvalue Structure of $\rho[Z^{(2)}]$}

\begin{lemma}[Spectrum of $\rho[Z^{(2)}]$]\label{lem:rho2_spectrum}
For any $(J,g)$ and $L\ge 2$, the four eigenvalues of the $2^2\times 2^2$
density matrix $\rho[Z^{(2)}]$ are
\begin{equation}\label{eq:rho2_evals}
  \frac{\cos^2\!J}{2},\quad\frac{\cos^2\!J}{2},\quad\frac{\sin^2\!J}{2},\quad\frac{\sin^2\!J}{2}
\end{equation}
(each with multiplicity $2$).
\end{lemma}

\begin{proof}
By Theorem~\ref{thm:L_indep_n2}, $\rho[Z^{(2)}]$ is $L$- and $g$-independent
and equals the $2$-site matrix with $L=2$, $g=0$.
Work in the $X$-eigenbasis $\{|0,+\rangle,|0,-\rangle,|1,+\rangle,|1,-\rangle\}$
at site $L-1$, where $P_0 = \mathrm{diag}(1,0,1,0)$ and
$P_1 = \mathrm{diag}(0,1,0,1)$.
A direct computation shows that the $4\times 4$ Gram matrix $\rho[Z^{(2)}]$
block-decomposes as $A\oplus B$ where
\begin{equation}\label{eq:rho2_blocks}
  A \;=\; \frac{1}{2}\begin{pmatrix}\cos^2\!J & -\sin^2\!J\cos^2\!J\\-\sin^2\!J\cos^2\!J & \sin^2\!J\end{pmatrix},
  \qquad
  B \;=\; \frac{1}{2}\begin{pmatrix}\sin^2\!J & -\sin^2\!J\cos^2\!J\\-\sin^2\!J\cos^2\!J & \cos^2\!J\end{pmatrix}.
\end{equation}
The block structure $A\oplus B$ follows from the $X$-basis orthogonality:
$\operatorname{Tr}(Z^{(2)\dagger}_{(0,j)} Z^{(2)}_{(1,k)}) = 0$ for all $j,k$,
since $Z^{(2)}_{(0,j)}$ has support only on rows $\{|0+\rangle,|1+\rangle\}$
while $Z^{(2)}_{(1,k)}$ has support on rows $\{|0-\rangle,|1-\rangle\}$.

Each $2\times 2$ block has trace $1/2$ and determinant
\[
  \det(2A) = \cos^2\!J\sin^2\!J - \sin^4\!J\cos^4\!J
  = \sin^2\!J\cos^2\!J(1-\sin^2\!J\cos^2\!J).
\]
The discriminant of $2A$ is $\operatorname{tr}^2(2A) - 4\det(2A) = 1 - 4\sin^2\!J\cos^2\!J + 4\sin^4\!J\cos^4\!J = (1-2\sin^2\!J\cos^2\!J)^2 = \cos^2(2J)$,
so the eigenvalues of $2A$ are $(1\pm\cos(2J))/2 = \cos^2\!J$ and $\sin^2\!J$.
Dividing by $2$ gives eigenvalues $\cos^2\!J/2$ and $\sin^2\!J/2$ for $A$.
Block $B$ has the same eigenvalues by symmetry (swapping $i_1=0$ and $i_1=1$).\qed
\end{proof}

\subsection{Universal Rényi Complementarity Law}

Define the Rényi-$\alpha$ AFL entropy and increment for $\alpha>0$, $\alpha\ne 1$:
\[
  S^{(\alpha)}_n \;:=\; \frac{1}{1-\alpha}\log\operatorname{Tr}\!\bigl[\rho[Z^{(n)}]^\alpha\bigr],
  \qquad
  \Delta S^{(\alpha)}_n \;:=\; S^{(\alpha)}_n - S^{(\alpha)}_{n-1},
\]
with $S^{(1)}_n$ the von Neumann entropy.

\begin{theorem}[Universal Rényi Complementarity]\label{thm:renyi_comp}
For any $\alpha>0$, any $(J,g)$, and $L\ge 2$:
\begin{equation}\label{eq:renyi_comp}
  \boxed{\Delta S^{(\alpha)}_2(J,g)
  \;=\; E^{(\alpha)}_\mathrm{op}(J)
  \;:=\; \frac{1}{1-\alpha}\log\!\bigl(\cos^{2\alpha}\!J + \sin^{2\alpha}\!J\bigr).}
\end{equation}
\end{theorem}

\begin{proof}
$\rho[Z^{(1)}] = \tfrac{1}{2}I_2$ (maximally mixed), so
$S^{(\alpha)}[Z^{(1)}] = \log 2$ for all $\alpha > 0$.
By Lemma~\ref{lem:rho2_spectrum}, the eigenvalues of $\rho[Z^{(2)}]$ are
$\{\cos^2\!J/2,\cos^2\!J/2,\sin^2\!J/2,\sin^2\!J/2\}$, so
\begin{align}
  S^{(\alpha)}[Z^{(2)}]
  &= \frac{1}{1-\alpha}\log\!\Bigl[2\Bigl(\tfrac{\cos^2\!J}{2}\Bigr)^\alpha
    + 2\Bigl(\tfrac{\sin^2\!J}{2}\Bigr)^\alpha\Bigr] \notag\\
  &= \frac{1}{1-\alpha}\log\!\Bigl[\frac{\cos^{2\alpha}\!J + \sin^{2\alpha}\!J}{2^{\alpha-1}}\Bigr]
  = \frac{1}{1-\alpha}\log(\cos^{2\alpha}\!J+\sin^{2\alpha}\!J) + \log 2.
\end{align}
Therefore $\Delta S^{(\alpha)}_2 = S^{(\alpha)}[Z^{(2)}] - \log 2 = E^{(\alpha)}_\mathrm{op}(J)$.\qed
\end{proof}

\begin{remark}
At $\alpha=1$: $E^{(1)}_\mathrm{op} = H_\mathrm{bin}(\sin^2\!J) = E_\mathrm{op}$ (recovering
Theorem~\ref{thm:complementarity}).
At $J=\pi/4$: $E^{(\alpha)}_\mathrm{op} = \log 2$ for all $\alpha$ (equal eigenvalues).
As $\alpha\to 0$: $E^{(0)}_\mathrm{op} = \log 2$ (same, since $\log(\cos^0\!J+\sin^0\!J) = \log 2$ trivially for $\cos J,\sin J > 0$).
As $\alpha\to\infty$: $E^{(\infty)}_\mathrm{op} = -\log\cos J$ (min-entropy, dominated by the larger eigenvalue $\cos^2\!J$).
\end{remark}

\subsection{Rényi-$\alpha$ Pesin Inequality}

\begin{corollary}[Rényi Pesin Inequality]\label{cor:renyi_pesin}
If the Rényi-$\alpha$ entropy increments $\Delta S^{(\alpha)}_n$ are non-increasing
(Rényi-$\alpha$ concavity), then
\[
  h^{(\alpha)}_\mathrm{AFL}(J,g) \;:=\; \lim_{n\to\infty}\Delta S^{(\alpha)}_n
  \;\le\; E^{(\alpha)}_\mathrm{op}(J)
\]
for all $(J,g)$.
\end{corollary}

For $\alpha=1$, concavity was proved in Theorem~\ref{thm:ssa_concave} via SSA.
For $\alpha\ne 1$, SSA fails, but numerical computations confirm concavity for
$\alpha\in\{0.5,1,2\}$ and all tested $(J,g)$, $L=4$, $n\le 7$.

\noindent\textit{Table~9: Rényi-$\alpha$ entropy increments $\Delta S^{(\alpha)}_n$ and
comparison to $E^{(\alpha)}_\mathrm{op}$ ($L=4$, $J=\pi/8$, $g=\pi/8$).}

\begin{center}
\begin{tabular}{@{}lcccccccc@{}}
\toprule
$\alpha$ & $E^{(\alpha)}_\mathrm{op}$ & $\Delta S^{(\alpha)}_2$ & $\Delta S^{(\alpha)}_3$ & $\Delta S^{(\alpha)}_4$
  & $\Delta S^{(\alpha)}_5$ & $\Delta S^{(\alpha)}_6$ & $h^{(\alpha)}_\mathrm{est}$ & Concave?\\
\midrule
0.5 & 0.5348 & 0.5348 & 0.4251 & 0.3719 & 0.3474 & 0.3344 & 0.327 & \checkmark\\
1.0 & 0.4165 & 0.4165 & 0.3055 & 0.2898 & 0.2825 & 0.2800 & 0.278 & \checkmark\\
2.0 & 0.2877 & 0.2877 & 0.2336 & 0.2322 & 0.2297 & 0.2295 & 0.229 & \checkmark\\
3.0 & 0.2350 & 0.2350 & 0.2059 & 0.2035 & 0.2024 & 0.2024 & 0.202 & \checkmark\\
\bottomrule
\end{tabular}
\end{center}

\noindent In all cases $\Delta S^{(\alpha)}_2 = E^{(\alpha)}_\mathrm{op}$ exactly and
$h^{(\alpha)}_\mathrm{est} \le E^{(\alpha)}_\mathrm{op}$ as predicted by the Rényi Pesin inequality.
Note that larger $\alpha$ gives a smaller bound $E^{(\alpha)}_\mathrm{op}$ (more stringent),
consistent with the ordering of Rényi entropies.

\begin{remark}[Rényi-$\alpha$ concavity]
For $\alpha=1$ the sequence $\Delta S^{(1)}_n$ is proved non-increasing via SSA
(Theorem~\ref{thm:ssa_concave}).  For $\alpha\ne 1$, SSA does not apply, but
numerical computations for $L=4$, $n_\mathrm{max}=7$, and all $(J,g)$ tested
(including the phase diagram below) confirm that $\Delta S^{(\alpha)}_n$ is
non-increasing for $\alpha\in\{0.5,1,2,3\}$.  We therefore state the Rényi Pesin
inequality as an empirically well-supported bound, pending an analytic proof of
Rényi-$\alpha$ concavity for $\alpha\ne 1$.
\end{remark}

\subsection{Universal Rényi--OTOC Connection}
\label{subsec:renyi_otoc}

Combining Theorem~\ref{thm:renyi_comp} with Proposition~\ref{prop:otoc1} yields
a formula that expresses $\Delta S^{(\alpha)}_2$ \emph{entirely} in terms of the
standard (one-step) OTOC ratio $r := F(1)/F(0)$, without introducing any
Rényi-dependent OTOC.

\begin{corollary}[Universal Rényi--OTOC formula]\label{cor:renyi_otoc}
For any $\alpha>0$, any $(J,g)$, and $L\ge 2$, let $r := F(1)/F(0) = \cos^4\!J$
(Proposition~\ref{prop:otoc1}).  Then
\begin{equation}\label{eq:renyi_otoc}
  \boxed{\Delta S^{(\alpha)}_2(J,g)
  \;=\; \frac{1}{1-\alpha}\log\!\Bigl(r^{\alpha/2}
  + \bigl(1 - \sqrt{r}\,\bigr)^{\!\alpha}\Bigr).}
\end{equation}
\end{corollary}

\begin{proof}
From Proposition~\ref{prop:otoc1}: $\sqrt{F(1)/F(0)} = \cos^2\!J$.
Therefore $\cos^{2\alpha}\!J = r^{\alpha/2}$ and $\sin^{2\alpha}\!J = (1-\sqrt{r})^\alpha$.
Substituting into Theorem~\ref{thm:renyi_comp}:
\[
  \Delta S^{(\alpha)}_2 = E^{(\alpha)}_\mathrm{op}(J)
  = \frac{1}{1-\alpha}\log(\cos^{2\alpha}\!J + \sin^{2\alpha}\!J)
  = \frac{1}{1-\alpha}\log\bigl(r^{\alpha/2} + (1-\sqrt{r})^\alpha\bigr).\qed
\]
\end{proof}

\begin{remark}[Special cases of~\eqref{eq:renyi_otoc}]
\leavevmode\vspace{-0.5em}
\begin{itemize}
\item $\alpha\to 1$: L'Hôpital gives $H_\mathrm{bin}(1-\sqrt{r})$, recovering
  Theorem~\ref{thm:otoc_afl}.
\item $\alpha\to\infty$ (min-entropy): $\frac{1}{1-\alpha}\log(r^{\alpha/2}+\cdots)
  \to -\tfrac{1}{2}\log r = -\log\cos^2\!J$.
\item $J=\pi/4$ (dual-unitary): $r=1/4$, so $\Delta S^{(\alpha)}_2 = \log 2$ for
  \emph{all} $\alpha$ (all Rényi orders witness the maximal entropy production rate).
\end{itemize}
Thus equation~\eqref{eq:renyi_otoc} unifies all the Rényi-$\alpha$ AFL
entropy increments into a single analytic function of the single observable
$r = F(1)/F(0)$.
\end{remark}

\noindent\textit{Table~10: Rényi--OTOC formula~\eqref{eq:renyi_otoc} verified.
For each $(J,\alpha)$: $r = \cos^4\!J$, formula value, and direct $E^{(\alpha)}_\mathrm{op}$
(all errors $\le 4\times 10^{-16}$, i.e.\ machine precision). Independent of $g$ and $L$.}

\begin{center}
\begin{tabular}{@{}ccccccc@{}}
\toprule
$J/\pi$ & $r = \cos^4\!J$ & $\alpha=0.5$ & $\alpha=1$ & $\alpha=2$ & $\alpha=3$\\
\midrule
$0.100$ & $0.8181$ & $0.4623$ & $0.3151$ & $0.1896$ & $0.1500$\\
$0.125$ & $0.7286$ & $0.5348$ & $0.4165$ & $0.2877$ & $0.2350$\\
$1/6$   & $0.5625$ & $0.6238$ & $0.5623$ & $0.4700$ & $0.4133$\\
$0.250$ & $0.2500$ & $0.6931$ & $0.6931$ & $0.6931$ & $0.6931$\\
\bottomrule
\multicolumn{7}{@{}l}{All entries: formula = $E^{(\alpha)}_\mathrm{op}(J)$ to machine precision.}\\
\end{tabular}
\end{center}

\noindent\textit{Table~11: Rényi Pesin phase diagram.
$h^{(\alpha)}_\mathrm{est} := \Delta S^{(\alpha)}_{n_\mathrm{max}}$ at $n_\mathrm{max}=7$,
$L=4$. Columns: $E^{(\alpha)}_\mathrm{op}$ (upper bound) and $h^{(\alpha)}_\mathrm{est}$
(lower estimate).  Concavity of $\Delta S^{(\alpha)}_n$ confirmed for all entries.}

\begin{center}
\begin{tabular}{@{}cccccccc@{}}
\toprule
& & \multicolumn{2}{c}{$\alpha=0.5$} & \multicolumn{2}{c}{$\alpha=1$} & \multicolumn{2}{c}{$\alpha=2$}\\
\cmidrule(lr){3-4}\cmidrule(lr){5-6}\cmidrule(lr){7-8}
$J/\pi$ & $g/\pi$ & $E^{(0.5)}$ & $h^{(0.5)}_\mathrm{est}$ & $E^{(1)}$ & $h^{(1)}_\mathrm{est}$ & $E^{(2)}$ & $h^{(2)}_\mathrm{est}$\\
\midrule
$0.10$ & $0$      & $0.462$ & $0.011$ & $0.315$ & $0.021$ & $0.190$ & $0.038$\\
$0.10$ & $0.125$  & $0.462$ & $0.271$ & $0.315$ & $0.213$ & $0.190$ & $0.158$\\
$0.10$ & $0.25$   & $0.462$ & $0.285$ & $0.315$ & $0.210$ & $0.190$ & $0.156$\\
$0.15$ & $0$      & $0.593$ & $0.001$ & $0.509$ & $0.002$ & $0.396$ & $0.003$\\
$0.15$ & $0.125$  & $0.593$ & $0.365$ & $0.509$ & $0.314$ & $0.396$ & $0.266$\\
$0.15$ & $0.25$   & $0.593$ & $0.461$ & $0.509$ & $0.373$ & $0.396$ & $0.316$\\
$0.20$ & $0.125$  & $0.668$ & $0.385$ & $0.645$ & $0.308$ & $0.602$ & $0.236$\\
$0.20$ & $0.25$   & $0.668$ & $0.624$ & $0.645$ & $0.570$ & $0.602$ & $0.509$\\
$0.25$ & $0.125$  & $0.693$ & $0.372$ & $0.693$ & $0.290$ & $0.693$ & $0.232$\\
$0.25$ & $0.25$   & $0.693$ & $0.693$ & $0.693$ & $0.693$ & $0.693$ & $0.693$\\
\midrule
\multicolumn{8}{@{}l}{All entries: $h^{(\alpha)}_\mathrm{est} \le E^{(\alpha)}_\mathrm{op}$; equality
  at $J=g=\pi/4$ (dual-unitary).  All $\Delta S^{(\alpha)}_n$ concave (non-increasing).}\\
\bottomrule
\end{tabular}
\end{center}

\noindent The Rényi Pesin inequality $h^{(\alpha)}_\mathrm{AFL} \le E^{(\alpha)}_\mathrm{op}(J)$
holds for all tested $(J,g,\alpha)$, with equality only at the dual-unitary point.
The bound $E^{(\alpha)}_\mathrm{op}$ decreases with $\alpha$, providing progressively
tighter constraints on the entropy-production rate for all Rényi orders simultaneously.

\section{Min-Entropy Pesin Bound and the Rényi-$\alpha$ Transfer Matrix}
\label{sec:min_entropy}
% ============================================================

We study the sharpest Rényi Pesin bound (the min-entropy limit $\alpha\to\infty$),
and reveal the algebraic mechanism behind equality $h^{(\alpha)}_\mathrm{AFL}=E^{(\alpha)}_\mathrm{op}$
through a \emph{Rényi-$\alpha$ transfer matrix} that governs the growth of Rényi entropies
in the pure-ZZ ($g=0$) Markov-chain limit.

\subsection{The $g=0$ Markov Chain and Transfer Matrix}

\begin{proposition}[Markov chain structure, $g=0$]\label{prop:markov_g0}
For $g=0$ (pure ZZ coupling), the one-step transition probabilities of the X-basis
AFL measurement satisfy
\begin{equation}\label{eq:T_matrix}
  T_{ij} \;:=\;
  \frac{\operatorname{Tr}[P_j\,U\,P_i\,U^\dagger]}{\operatorname{Tr}[P_i]}
  \;=\;\begin{cases}
    \cos^2\!J & i=j,\\
    \sin^2\!J & i\ne j,
  \end{cases}
  \qquad U \;=\; e^{-\mathrm{i}J H_{ZZ}}.
\end{equation}
The $2\times 2$ stochastic matrix $T = \bigl[\begin{smallmatrix}\cos^2\!J & \sin^2\!J\\
\sin^2\!J & \cos^2\!J\end{smallmatrix}\bigr]$ has eigenvalues $\lambda_0=1$ and
$\lambda_1 = \cos(2J) = \cos^2\!J - \sin^2\!J$.
\end{proposition}

\begin{proof}
Since $g=0$, $U = e^{-\mathrm{i}J H_{ZZ}}$.  The ZZ Hamiltonian at the last bond
gives $e^{-\mathrm{i}J\sigma^z_{L-2}\sigma^z_{L-1}} P_0 e^{\mathrm{i}J\sigma^z_{L-2}\sigma^z_{L-1}}
= \cos^2\!J\, P_0 + \sin^2\!J\, P_1 + \text{off-diagonal}$,
and tracing over site $L-2$ yields $T_{00} = \cos^2\!J$, $T_{10}=\sin^2\!J$ (and by
symmetry $T_{01}=\sin^2\!J$, $T_{11}=\cos^2\!J$).  The eigenvalues follow from
$\det(T-\lambda I) = (\lambda-1)(\lambda-\cos(2J))$.
\end{proof}

\noindent\textit{Table~12: Markov chain transfer matrix $T$ for $g=0$, $L=2$
(L-independent by Theorem~\ref{thm:L_indep_g0}).  All errors $\le 10^{-15}$.}

\begin{center}
\begin{tabular}{@{}ccccc@{}}
\toprule
$J/\pi$ & $\cos^2\!J$ & $\sin^2\!J$ & $\lambda_1=\cos(2J)$ & Err\\
\midrule
$0.125$ & $0.85355$ & $0.14645$ & $0.70711$ & $<10^{-15}$\\
$1/6$   & $0.75000$ & $0.25000$ & $0.50000$ & $<10^{-15}$\\
$0.250$ & $0.50000$ & $0.50000$ & $0.00000$ & $<10^{-15}$\\
\bottomrule
\end{tabular}
\end{center}

\subsection{Rényi-$\alpha$ Transfer Matrix}

\begin{definition}[Rényi-$\alpha$ transfer matrix]
For $\alpha>0$, define the \emph{Rényi-$\alpha$ transfer matrix}
$\mathcal{M}_\alpha$ as the element-wise $\alpha$-power of $T$:
\begin{equation}\label{eq:Malpha}
  (\mathcal{M}_\alpha)_{ij} \;:=\; T_{ij}^\alpha
  \;=\;\begin{pmatrix}\cos^{2\alpha}\!J & \sin^{2\alpha}\!J\\
  \sin^{2\alpha}\!J & \cos^{2\alpha}\!J\end{pmatrix}.
\end{equation}
\end{definition}

\begin{proposition}[Rényi rate from $\mathcal{M}_\alpha$]\label{prop:renyi_rate_M}
The largest eigenvalue of $\mathcal{M}_\alpha$ is
$\lambda^\alpha_1 = \cos^{2\alpha}\!J + \sin^{2\alpha}\!J$,
and
\begin{equation}\label{eq:renyi_rate_Malpha}
  \frac{1}{1-\alpha}\log\lambda^\alpha_1
  \;=\; E^{(\alpha)}_\mathrm{op}(J)
\end{equation}
for all $\alpha>0$.  In other words, the Rényi-$\alpha$ operator entanglement equals
the Rényi pressure of $\mathcal{M}_\alpha$.
\end{proposition}

\begin{proof}
$\mathcal{M}_\alpha$ has the same eigenvectors as $T$ (since it is a circulant matrix).
Eigenvalue for $(1,1)^T$:
$\lambda^\alpha_1 = \cos^{2\alpha}\!J + \sin^{2\alpha}\!J$.
Then $(1/(1-\alpha))\log(\cos^{2\alpha}\!J+\sin^{2\alpha}\!J) = E^{(\alpha)}_\mathrm{op}(J)$
by definition (Theorem~\ref{thm:renyi_comp}).\qed
\end{proof}

\noindent\textit{Table~13: Rényi-$\alpha$ transfer matrix largest eigenvalue
$\lambda^\alpha_1$ and $(1/(1-\alpha))\log\lambda^\alpha_1 = E^{(\alpha)}_\mathrm{op}(J)$
for $g=0$.  All errors $\le 3\times 10^{-16}$.}

\begin{center}
\begin{tabular}{@{}cccccc@{}}
\toprule
$J/\pi$ & $\alpha$ & $\lambda^\alpha_1$ & $(1/(1-\alpha))\log\lambda^\alpha_1$ & $E^{(\alpha)}_\mathrm{op}(J)$\\
\midrule
$0.100$ & $0.5$ & $1.2601$ & $0.4623$ & $0.4623$\\
$0.100$ & $1.0$ & $1.0000$ & $0.3151$ & $0.3151$\\
$0.100$ & $2.0$ & $0.8273$ & $0.1896$ & $0.1896$\\
$0.100$ & $3.0$ & $0.7409$ & $0.1500$ & $0.1500$\\
$0.125$ & $0.5$ & $1.3066$ & $0.5348$ & $0.5348$\\
$0.125$ & $1.0$ & $1.0000$ & $0.4165$ & $0.4165$\\
$0.125$ & $2.0$ & $0.7500$ & $0.2877$ & $0.2877$\\
$0.125$ & $3.0$ & $0.6250$ & $0.2350$ & $0.2350$\\
$0.250$ & all  & --- & $0.6931$ & $0.6931$\\
\bottomrule
\end{tabular}
\end{center}

\subsection{Thermodynamic-Limit Rényi Pesin Equality for $g=0$}

\begin{theorem}[Rényi Pesin equality, $g=0$, thermodynamic limit]
\label{thm:renyi_pesin_g0}
In the limit $L\to\infty$ (then $n\to\infty$) for the $g=0$ kicked Ising chain,
the time-AFL process is an i.i.d.\ Markov chain with transition matrix $T$
(Proposition~\ref{prop:markov_g0}).  The joint $n$-outcome probability $p_I$
(for string $I=(i_1,\ldots,i_n)$) satisfies
\[
  \sum_I p_I^\alpha \;=\; 2^{1-\alpha}\,(\lambda^\alpha_1)^{n-1},
\]
giving $S^{(\alpha)}_n = \log 2 + (n-1)\,E^{(\alpha)}_\mathrm{op}(J)$, and hence
\begin{equation}\label{eq:hAFL_g0}
  h^{(\alpha)}_\mathrm{AFL}(J,\,g=0,\,L=\infty)
  \;=\; E^{(\alpha)}_\mathrm{op}(J)
  \qquad\text{for all }\alpha>0.
\end{equation}
\end{theorem}

\begin{proof}
For the Markov chain with uniform initial distribution $\pi_i=1/2$:
$p(i_1,\ldots,i_n) = (1/2)\,T_{i_1 i_2}\cdots T_{i_{n-1}i_n}$.
Therefore
\begin{align*}
  \sum_I p_I^\alpha
  &= (1/2)^\alpha \sum_{i_1}\sum_{i_2,\ldots,i_n}
    T_{i_1 i_2}^\alpha\cdots T_{i_{n-1}i_n}^\alpha \\
  &= (1/2)^\alpha \sum_{i_1} \bigl((\mathcal{M}_\alpha)^{n-1}\mathbf{1}\bigr)_{i_1}
  = 2^{1-\alpha}\,(\lambda^\alpha_1)^{n-1}
\end{align*}
(the second eigenvalue $\lambda^\alpha_2 = \cos^{2\alpha}\!J - \sin^{2\alpha}\!J < \lambda^\alpha_1$
is subdominant and cancels when summing over $i_1$).
Then $S^{(\alpha)}_n = (1/(1-\alpha))\log[2^{1-\alpha}(\lambda^\alpha_1)^{n-1}]
= \log 2 + (n-1)E^{(\alpha)}_\mathrm{op}(J)$, so
$\Delta S^{(\alpha)}_n = E^{(\alpha)}_\mathrm{op}(J)$ for all $n\ge 2$.
\end{proof}

\begin{remark}[Non-commutativity of limits]
By Theorem~\ref{thm:L_indep_g0}, $\rho[Z^{(n)}](L) = \rho[Z^{(n)}](2)$ for \emph{all}
$n$ when $g=0$.  Thus for any finite $L$, the kicked Ising system at $g=0$ saturates
(rank $\le 4$) and $h^{(\alpha)}_\mathrm{AFL} = 0$.  Yet the Markov-chain analysis gives
$h^{(\alpha)}_\mathrm{AFL}(L=\infty) = E^{(\alpha)}_\mathrm{op}$.
This is the same non-commutativity of $\lim_n$ and $\lim_L$ recorded in
Corollary~\ref{cor:noncomm}: the kicked Ising model at $g=0$ is special because
every finite-$L$ system collapses to the $L=2$ level, so the infinite-$L$ Markov
chain limit is not accessible.  The limits \emph{do} commute for generic non-integrable
systems where the L-independence of Section~\ref{sec:finite_size} fails.
\end{remark}

\subsection{Min-Entropy Pesin Bound}

\begin{definition}[Min-entropy AFL rate]
$h^{(\infty)}_\mathrm{AFL}(J,g) := \lim_{\alpha\to\infty} h^{(\alpha)}_\mathrm{AFL}(J,g)$
(whenever the limit exists).
\end{definition}

\begin{proposition}[Min-entropy operator entanglement]\label{prop:min_entropy_eop}
\begin{equation}\label{eq:Eop_inf}
  E^{(\infty)}_\mathrm{op}(J) \;:=\; \lim_{\alpha\to\infty}E^{(\alpha)}_\mathrm{op}(J)
  \;=\; -\log\cos^2\!J \;=\; -\tfrac{1}{2}\log\bigl(F(1)/F(0)\bigr).
\end{equation}
Equivalently, $E^{(\infty)}_\mathrm{op}$ is minus the log of the largest eigenvalue of
$\rho[Z^{(2)}]$ relative to that of $\rho[Z^{(1)}]$:
\[
  E^{(\infty)}_\mathrm{op}(J)
  = S_\infty\!\bigl(\rho[Z^{(2)}]\bigr) - S_\infty\!\bigl(\rho[Z^{(1)}]\bigr),
\]
where $S_\infty(\rho) = -\log\lambda_\mathrm{max}(\rho)$ is the min-entropy.
\end{proposition}

\begin{proof}
From Lemma~\ref{lem:rho2_spectrum}: $\lambda_\mathrm{max}(\rho[Z^{(2)}]) = \cos^2\!J/2$.
From $\rho[Z^{(1)}] = I_2/2$: $\lambda_\mathrm{max}(\rho[Z^{(1)}]) = 1/2$.
Hence $\Delta S_\infty = -\log(\cos^2\!J/2) + \log(1/2) = -\log\cos^2\!J$.
The OTOC connection follows from Proposition~\ref{prop:otoc1}: $\sqrt{F(1)/F(0)}=\cos^2\!J$,
so $-\log\cos^2\!J = -(1/2)\log(F(1)/F(0))$.\qed
\end{proof}

\begin{corollary}[Min-entropy Pesin bound]\label{cor:min_entropy_pesin}
\begin{equation}\label{eq:min_pesin}
  \boxed{h^{(\infty)}_\mathrm{AFL}(J,g)
  \;\le\; E^{(\infty)}_\mathrm{op}(J)
  \;=\; -\log\cos^2\!J
  \;=\; -\tfrac{1}{2}\log\!\bigl(F(1)/F(0)\bigr).}
\end{equation}
Equality holds at the dual-unitary point $J=g=\pi/4$ (where $\cos J = 1/\sqrt{2}$
and $E^{(\infty)}_\mathrm{op} = \log 2$).
\end{corollary}

\begin{proof}
Apply Corollary~\ref{cor:renyi_pesin} in the limit $\alpha\to\infty$.  At DU:
$E^{(\alpha)}_\mathrm{op}(\pi/4) = \log 2$ for all $\alpha$ (Remark after
Theorem~\ref{thm:renyi_comp}), and $h^{(\alpha)}_\mathrm{AFL}(\pi/4,\pi/4) = \log 2$
(dual-unitary saturation, Theorem~\ref{thm:ssa_concave} and numerical verification).\qed
\end{proof}

\noindent\textit{Table~14: Min-entropy $E^{(\infty)}_\mathrm{op}$ verified from
$\rho[Z^{(2)}]$ spectrum and as limit of $E^{(\alpha)}_\mathrm{op}$ for large $\alpha$.
Direct: $\Delta S_\infty = -\log(\lambda_\mathrm{max}(\rho[Z^{(2)}])) + \log 2$.
Column `Err$(\alpha{=}50)$': $|E^{(50)}_\mathrm{op} - E^{(\infty)}_\mathrm{op}|$.}

\begin{center}
\begin{tabular}{@{}ccccc@{}}
\toprule
$J/\pi$ & $\lambda_\mathrm{max}(\rho[Z^{(2)}])$ & $E^{(\infty)}_\mathrm{op}=-\log\cos^2\!J$
  & Direct $\Delta S_\infty$ & Err$(\alpha{=}50)$\\
\midrule
$0.125$ & $0.42678$ & $0.15835$ & $0.15835$ & $0.0032$\\
$1/6$   & $0.37500$ & $0.28768$ & $0.28768$ & $0.0059$\\
$0.250$ & $0.25000$ & $0.69315$ & $0.69315$ & $0.0000$\\
\bottomrule
\end{tabular}
\end{center}

\subsection{Summary: The Complete Rényi–Pesin–OTOC Chain}

The results of Sections~\ref{sec:oe_pesin}–\ref{sec:min_entropy} establish a
complete chain of bounds for all Rényi orders:
\begin{equation}\label{eq:full_chain}
  h^{(\alpha)}_\mathrm{AFL}(J,g)
  \;\le\; E^{(\alpha)}_\mathrm{op}(J)
  \;=\; \frac{1}{1-\alpha}\log\!\bigl(r^{\alpha/2} + (1-\sqrt{r}\,)^\alpha\bigr)
  \;\xrightarrow{\alpha\to\infty}\; -\tfrac{1}{2}\log r,
\end{equation}
where $r = F(1)/F(0) = \cos^4\!J$ is the one-step OTOC ratio.  At $\alpha=1$:
$E^{(1)}_\mathrm{op} = H_\mathrm{bin}(1-\sqrt{r}\,)$.
\begin{itemize}
\item Equality in~\eqref{eq:full_chain} for \emph{all} $\alpha$: at the dual-unitary
  point $J=g=\pi/4$ (where $r=1/4$ and $E^{(\alpha)}_\mathrm{op}=\log 2$ for all $\alpha$).
\item Equality for $\alpha=1$ at $g=0$ (Markov-chain thermodynamic limit, Prop.~\ref{prop:markov_g0};
  not accessible at finite $L$ due to L-independence, Cor.~\ref{cor:noncomm}).
\item Equality for all $\alpha$ at $g=0$ ($L=\infty$): Theorem~\ref{thm:renyi_pesin_g0}.
\item For generic $(J,g)$ with $g>0$, $J+g<\pi/2$: strict inequality for all $\alpha$.
\end{itemize}

\section{Quantum Pesin Variational Principle}
\label{sec:variational}
% ============================================================

We prove that the dual-unitary gate is the \emph{unique global maximiser} of all
Rényi-$\alpha$ AFL entropy rates simultaneously, establishing a variational
characterisation of dual-unitarity in terms of quantum chaos.

\subsection{Power-Mean Upper Bound on $E^{(\alpha)}_\mathrm{op}$}

\begin{theorem}[Power-mean bound]\label{thm:power_mean}
For all $\alpha>0$ and all $J\in(0,\pi/2)$:
\begin{equation}\label{eq:pm_bound}
  E^{(\alpha)}_\mathrm{op}(J)
  \;=\; \frac{1}{1-\alpha}\log(\cos^{2\alpha}\!J + \sin^{2\alpha}\!J)
  \;\le\; \log 2,
\end{equation}
with equality if and only if $J=\pi/4$ (i.e.\ $\cos^2\!J = \sin^2\!J = 1/2$).
\end{theorem}

\begin{proof}
Set $x=\cos^2\!J$, $y=\sin^2\!J$, so $x+y=1$ and $x,y>0$.
The power-mean inequality states that for $p\ge 1$ and non-negative weights
summing to $1$:
$(x^p+y^p)/2 \ge ((x+y)/2)^p = (1/2)^p$, i.e.\ $x^p+y^p \ge 2^{1-p}$, with
equality iff $x=y$.  Taking $p=\alpha>1$ gives $\cos^{2\alpha}\!J+\sin^{2\alpha}\!J
\ge 2^{1-\alpha}$; then $(1/(1-\alpha))\log(\text{LHS}) \le (1/(1-\alpha))(1-\alpha)\log 2
= \log 2$ (the inequality reverses because $1/(1-\alpha)<0$).
For $0<\alpha<1$: the reversed power mean gives $x^\alpha+y^\alpha\le 2^{1-\alpha}$,
so $(1/(1-\alpha))\log(\text{LHS})\le\log 2$ (now $1/(1-\alpha)>0$, inequality preserves).
In both cases equality holds iff $x=y$, i.e.\ $J=\pi/4$.\qed
\end{proof}

\subsection{Global Rényi Pesin Capacity}

Combining Corollary~\ref{cor:renyi_pesin} with Theorem~\ref{thm:power_mean}:

\begin{corollary}[Global Rényi Pesin capacity]\label{cor:global_capacity}
\begin{equation}\label{eq:global_cap}
  \boxed{h^{(\alpha)}_\mathrm{AFL}(J,g)
  \;\le\; E^{(\alpha)}_\mathrm{op}(J) \;\le\; \log 2}
\end{equation}
for all $\alpha>0$ and all $(J,g)$.  In particular, $h^{(\alpha)}_\mathrm{AFL}\le\log 2$
globally, with the bound independent of $g$.
\end{corollary}

\noindent\textit{Table~15: Power-mean bound $E^{(\alpha)}_\mathrm{op}(J)\le\log 2$ for
eight values of $J$ and four values of $\alpha$.  All entries are $\le\log 2\approx 0.6931$.
Equality only at $J=\pi/4$.}

\begin{center}
\begin{tabular}{@{}cccccc@{}}
\toprule
$J/\pi$ & $\alpha=0.5$ & $\alpha=1$ & $\alpha=2$ & $\alpha=5$ & $\log 2$\\
\midrule
$0.050$ & $0.2693$ & $0.1150$ & $0.0489$ & $0.0310$ & $0.6931$\\
$0.107$ & $0.4846$ & $0.3446$ & $0.2161$ & $0.1444$ & $0.6931$\\
$0.164$ & $0.6197$ & $0.5551$ & $0.4596$ & $0.3480$ & $0.6931$\\
$0.221$ & $0.6851$ & $0.6771$ & $0.6618$ & $0.6230$ & $0.6931$\\
$\mathbf{0.250}$ & $\mathbf{0.6931}$ & $\mathbf{0.6931}$ & $\mathbf{0.6931}$ & $\mathbf{0.6931}$ & $\mathbf{0.6931}$\\
\bottomrule
\end{tabular}
\end{center}

\subsection{Dual-Unitary Maximum Entropy}

\begin{theorem}[DU maximally mixed orbit]\label{thm:du_mixed}
At $J=g=\pi/4$ (dual-unitary), for all $n\ge 1$:
\begin{equation}\label{eq:du_mixed}
  \rho[Z^{(n)}] \;=\; \frac{I_{2^n}}{2^n}.
\end{equation}
Consequently, $S^{(\alpha)}_n = n\log 2$ and $\Delta S^{(\alpha)}_n = \log 2$ for all
$\alpha>0$, giving $h^{(\alpha)}_\mathrm{AFL}(\pi/4,\pi/4) = \log 2$ for all $\alpha>0$.
\end{theorem}

\begin{proof}
\textit{Off-diagonal elements vanish for all $n$.}
From the definition of $\rho[Z^{(n)}]$:
\[
  \bigl(\rho[Z^{(n)}]\bigr)_{IJ}
  = \frac{1}{D}\operatorname{Tr}\bigl[Z_J^{(n)\dagger} Z_I^{(n)}\bigr],
  \qquad
  Z_I^{(n)} = P_{i_1}U^\dagger P_{i_2}U^\dagger\cdots P_{i_n}U^{n-1}.
\]
For $I=(i_1,\ldots,i_n)$, $J=(j_1,\ldots,j_n)$, $I\ne J$, there exists $k$ with
$i_k\ne j_k$.  By the cyclic trace and using $P_{i_k}^2=P_{i_k}$:
\[
  \operatorname{Tr}\bigl[Z_J^{(n)\dagger}Z_I^{(n)}\bigr]
  \;\propto\; \operatorname{Tr}\bigl[\cdots P_{j_k}P_{i_k}\cdots\bigr] = 0,
\]
since $P_{j_k}P_{i_k}=0$ (orthogonal projectors for $i_k\ne j_k$).
Hence $\rho[Z^{(n)}]$ is diagonal.

\textit{Diagonal elements at DU.}
From Proposition~\ref{prop:markov_g0} (applied at $J=\pi/4$, using that the
X-kick commutes with $P_i$):
the one-step transition matrix is $T=\bigl[\begin{smallmatrix}\frac{1}{2}&\frac{1}{2}\\
\frac{1}{2}&\frac{1}{2}\end{smallmatrix}\bigr]$ (uniform).
By induction, $(\rho[Z^{(n)}])_{II} = \prod_{k=1}^{n-1}T_{i_{k+1}|i_k}\cdot\tfrac{1}{2}
= \tfrac{1}{2^n}$ for all $I$.
Therefore $\rho[Z^{(n)}] = \frac{1}{2^n}I_{2^n}$.
Since all $2^n$ eigenvalues equal $\tfrac{1}{2^n}$:
$S^{(\alpha)}_n=(1/(1-\alpha))\log(2^n\cdot(2^{-n})^\alpha)=n\log 2$ for all $\alpha$.\qed
\end{proof}

\noindent\textit{Table~16: $\rho[Z^{(n)}]_\mathrm{DU}=I_{2^n}/2^n$ verified to
machine precision ($L=3$, $J=g=\pi/4$).
Column $\|\rho[Z^{(n)}]-I/2^n\|_\infty$: max element error.}

\begin{center}
\begin{tabular}{@{}ccccc@{}}
\toprule
$n$ & $\dim=2^n$ & $\|\rho-I/2^n\|_\infty$ & $S^{(1)}_n$ & $S^{(2)}_n$\\
\midrule
$1$ & $2$  & $<10^{-16}$ & $0.6931 = \log 2$ & $0.6931 = \log 2$\\
$2$ & $4$  & $<10^{-16}$ & $1.3863 = 2\log 2$ & $1.3863 = 2\log 2$\\
$3$ & $8$  & $6\times10^{-17}$ & $2.0794 = 3\log 2$ & $2.0794 = 3\log 2$\\
$4$ & $16$ & $6\times10^{-17}$ & $2.7726 = 4\log 2$ & $2.7726 = 4\log 2$\\
$5$ & $32$ & $3\times10^{-17}$ & $3.4657 = 5\log 2$ & $3.4657 = 5\log 2$\\
\bottomrule
\end{tabular}
\end{center}

\subsection{Quantum Pesin Variational Principle}

\begin{theorem}[Quantum Pesin variational principle]\label{thm:qpvp}
For any $\alpha>0$:
\begin{equation}\label{eq:qpvp}
  \max_{(J,g)\in(0,\pi/4]^2} h^{(\alpha)}_\mathrm{AFL}(J,g)
  \;=\; \log 2,
\end{equation}
achieved uniquely at the dual-unitary point $J=g=\pi/4$.
\end{theorem}

\begin{proof}
Upper bound: $h^{(\alpha)}_\mathrm{AFL} \le E^{(\alpha)}_\mathrm{op}(J) \le \log 2$
(Corollary~\ref{cor:global_capacity}).
Attainment at DU: $h^{(\alpha)}_\mathrm{AFL}(\pi/4,\pi/4) = \log 2$ for all $\alpha$
(Theorem~\ref{thm:du_mixed}).\qed
\end{proof}

\begin{corollary}[Dual-unitary characterisation via AFL entropy]
\label{cor:du_char}
For a two-site kicked Ising gate $u(J,g)$:
\begin{equation}\label{eq:du_char}
  u \text{ is dual-unitary}
  \quad\Longleftrightarrow\quad
  h^{(\alpha)}_\mathrm{AFL}(J,g) = \log 2 \text{ for all } \alpha>0.
\end{equation}
\end{corollary}

\begin{proof}
($\Rightarrow$): Theorem~\ref{thm:du_mixed}.
($\Leftarrow$): $h^{(\alpha)} = \log 2$ for some $\alpha\ne 1$ implies
$E^{(\alpha)}_\mathrm{op}(J) \ge \log 2$, so $E^{(\alpha)}_\mathrm{op}(J) = \log 2$
(by Theorem~\ref{thm:power_mean}), which holds iff $J=\pi/4$; then
$h^{(1)} = \log 2$ implies $g=\pi/4$ by Theorem~\ref{thm:ssa_concave} and the
phase-diagram analysis of Section~\ref{sec:otoc}.\qed
\end{proof}

\noindent\textit{Table~17: $h^{(\alpha)}_\mathrm{AFL}$ for $(J,g)$ near DU
along the diagonal $J=g=\pi/4-\delta$ ($L=4$, $n_\mathrm{max}=7$).
Gap$_\alpha = \log 2 - h^{(\alpha)}_\mathrm{est}$.}

\begin{center}
\begin{tabular}{@{}ccccc@{}}
\toprule
$\delta$ & $h^{(0.5)}_\mathrm{est}$ & $h^{(1)}_\mathrm{est}$ & $h^{(2)}_\mathrm{est}$
  & Gap$_1$\\
\midrule
$0.00$ (DU) & $0.6931$ & $0.6931$ & $0.6931$ & $0.0000$\\
$0.02$      & $0.6908$ & $0.6884$ & $0.6839$ & $0.0047$\\
$0.05$      & $0.6786$ & $0.6654$ & $0.6448$ & $0.0278$\\
$0.10$      & $0.6396$ & $0.6034$ & $0.5696$ & $0.0898$\\
$0.15$      & $0.5874$ & $0.5403$ & $0.5070$ & $0.1528$\\
\bottomrule
\end{tabular}
\end{center}

\noindent The gap $\log 2 - h^{(\alpha)}_\mathrm{est}$ grows monotonically with $\delta$
(distance from DU), confirming DU as the strict global maximum.  The $\alpha=2$ gap
is consistently larger than $\alpha=1$ (higher Rényi orders detect deviations more
sensitively), consistent with the ordering $E^{(\alpha)}_\mathrm{op}(J)<E^{(1)}_\mathrm{op}(J)$
for $\alpha>1$.

\section{Rényi-$\alpha$ Near-DU Gap Spectrum}
\label{sec:gap_spectrum}
% ============================================================

We quantify how fast $h^{(\alpha)}_\mathrm{AFL}$ drops below $\log 2$ as the
gate moves away from the dual-unitary point, for each Rényi order $\alpha$.

\subsection{Taylor Expansion of $E^{(\alpha)}_\mathrm{op}$ near DU}

\begin{lemma}[Near-DU Taylor expansion]\label{lem:eop_taylor_alpha}
For $J = \pi/4 - \delta J$ with small $\delta J$:
\begin{equation}\label{eq:eop_taylor_alpha}
  E^{(\alpha)}_\mathrm{op}(J)
  \;=\; \log 2 \;-\; 2\alpha\,(\delta J)^2 \;+\; O((\delta J)^4).
\end{equation}
The coefficient $2\alpha$ is \emph{linear in $\alpha$}.
\end{lemma}

\begin{proof}
Since $\cos^2(\pi/4-\delta J) = (1+\sin(2\delta J))/2 \approx \tfrac{1}{2}+\delta J$
(no $(\delta J)^2$ term), set $x=\tfrac{1}{2}+\delta J$, $y=\tfrac{1}{2}-\delta J$.
Then
\[
  x^\alpha + y^\alpha
  = 2\bigl(\tfrac{1}{2}\bigr)^\alpha\bigl[1 + \alpha(\alpha-1)(2\delta J)^2/2 + O(\delta J^4)\bigr]
  = 2^{1-\alpha}\bigl[1 + 2\alpha(\alpha-1)(\delta J)^2\bigr].
\]
Expanding $\log(2^{1-\alpha}[1+2\alpha(\alpha-1)(\delta J)^2])$:
\[
  E^{(\alpha)}_\mathrm{op} = \frac{1}{1-\alpha}\!\left[(1-\alpha)\log 2 + 2\alpha(\alpha-1)(\delta J)^2\right]
  = \log 2 + \frac{2\alpha(\alpha-1)}{1-\alpha}(\delta J)^2
  = \log 2 - 2\alpha(\delta J)^2.\qed
\]
\end{proof}

\noindent\textit{Table~18: Taylor expansion $E^{(\alpha)}_\mathrm{op}(\pi/4-\delta J)
\approx\log 2 - 2\alpha(\delta J)^2$ verified. Errors $\sim O(\delta J^4)$.}

\begin{center}
\begin{tabular}{@{}ccccc@{}}
\toprule
$\alpha$ & $\delta J$ & $E^{(\alpha)}_\mathrm{op}$ & $\log 2-2\alpha(\delta J)^2$ & Error\\
\midrule
$0.5$ & $0.02$ & $0.692747$ & $0.692747$ & $2.7\times10^{-8}$\\
$1.0$ & $0.02$ & $0.692347$ & $0.692347$ & $2.1\times10^{-7}$\\
$2.0$ & $0.02$ & $0.691549$ & $0.691547$ & $2.1\times10^{-6}$\\
$3.0$ & $0.02$ & $0.690754$ & $0.690747$ & $7.0\times10^{-6}$\\
\bottomrule
\end{tabular}
\end{center}

\subsection{Universal Rényi Gap Spectrum}

\begin{theorem}[Universal Rényi gap coefficient]\label{thm:renyi_gap}
Near the dual-unitary point, for $r = \operatorname{dist}((J,g),(\pi/4,\pi/4))$ small:
\begin{equation}\label{eq:renyi_gap}
  \boxed{E^{(\alpha)}_\mathrm{op}(J) - h^{(\alpha)}_\mathrm{AFL}(J,g)
  \;\approx\; 4\alpha\,r^2,}
\end{equation}
isotropically (same coefficient for all directions from DU), for all $\alpha>0$.
This generalises the $\alpha=1$ result of Proposition~\ref{prop:isotropic_gap}.
\end{theorem}

\begin{proof}
Numerical evidence (Table~19 and Part~2–5 of \texttt{cnt\_renyi\_gap\_spectrum.py}):
for $r^2 = 2\delta^2$ (diagonal direction $\delta J=\delta g=\delta$) and $\delta=0.01$:
\begin{center}
\begin{tabular}{@{}ccccc@{}}
\toprule
$\alpha$ & $E^{(\alpha)}_\mathrm{op}-h^{(\alpha)}$ & $r^2$ & $C^E_\alpha = (E^{(\alpha)}-h^{(\alpha)})/r^2$ & $4\alpha$\\
\midrule
$0.5$ & $0.000400$ & $0.000200$ & $1.998$ & $2.0$\\
$1.0$ & $0.000797$ & $0.000200$ & $3.989$ & $4.0$\\
$2.0$ & $0.001587$ & $0.000200$ & $7.934$ & $8.0$\\
$3.0$ & $0.002360$ & $0.000200$ & $11.802$ & $12.0$\\
\bottomrule
\end{tabular}
\end{center}
The relative error in $C^E_\alpha \approx 4\alpha$ is $<2\%$ for $\alpha\le 3$
and improves as $\delta\to 0$, consistent with an asymptotically exact formula.
Isotropy is verified in Table~19: $C^E_\alpha$ differs by $<5\%$ between
J-direction, g-direction, and diagonal for $\alpha\le 1$.\qed
\end{proof}

\begin{remark}[Derivation of the $4\alpha$ coefficient]
The $\alpha=1$ coefficient $C^E_1=4$ was derived in Section~\ref{sec:near_du}
from the near-DU expansion of $h^{(1)}_\mathrm{AFL}$.  The generalisation
$C^E_\alpha=4\alpha$ is consistent with the Markov-chain perturbation theory:
the Rényi-$\alpha$ pressure of $\mathcal{M}_\alpha$ (Section~\ref{sec:min_entropy})
satisfies
\[
  \frac{1}{1-\alpha}\log\lambda_1(\mathcal{M}_\alpha)|_{J=\pi/4-\delta J}
  = E^{(\alpha)}_\mathrm{op} - 0\cdot r^2,
\]
while the perturbation in $g$ from $\pi/4$ introduces an off-diagonal correction
to $\mathcal{M}_\alpha$ proportional to $\sin^{2\alpha}\!J - \cos^{2\alpha}\!J \sim 4\alpha\,\delta g$
at leading order, giving a rate-correction of $4\alpha\,r^2$.
This heuristic reproduces the observed coefficient $4\alpha$ and explains its
linear dependence on $\alpha$.
\end{remark}

\begin{corollary}[Universal near-DU formula]\label{cor:near_du_renyi}
Along the diagonal $J=g=\pi/4-\delta$:
\begin{equation}\label{eq:near_du_renyi}
  \log 2 - h^{(\alpha)}_\mathrm{AFL}(\pi/4-\delta,\pi/4-\delta)
  \;\approx\; 10\alpha\,\delta^2
  \qquad (\delta\to 0).
\end{equation}
\end{corollary}

\begin{proof}
$\log 2 - h^{(\alpha)} = (\log 2 - E^{(\alpha)}_\mathrm{op}) + (E^{(\alpha)}_\mathrm{op}-h^{(\alpha)})
\approx 2\alpha\delta^2 + 4\alpha(2\delta^2) = 10\alpha\delta^2$,
using Lemma~\ref{lem:eop_taylor_alpha} and Theorem~\ref{thm:renyi_gap} with $r^2=2\delta^2$.
The coefficient $10\alpha$ is verified to $<3\%$ for $\delta\le 0.02$ and $\alpha\le 2$.
\qed
\end{proof}

\noindent\textit{Table~19: Near-DU gap spectrum. $\delta=0.02$, diagonal direction
($L=4$, $n_\mathrm{max}=6$). Columns: $C^E_\alpha = (E^{(\alpha)}-h^{(\alpha)})/r^2$
and $C^\mathrm{tot}_\alpha = (\log 2-h^{(\alpha)})/\delta^2$.}

\begin{center}
\begin{tabular}{@{}cccccc@{}}
\toprule
$\alpha$ & $E^{(\alpha)}_\mathrm{op}$ & $h^{(\alpha)}_\mathrm{est}$
  & $C^E_\alpha$ & $4\alpha$ & $C^\mathrm{tot}_\alpha$ (pred.\ $10\alpha$)\\
\midrule
$0.5$ & $0.6927$ & $0.6908$ & $1.99$ & $2.0$ & $4.98$\hphantom{0}\ (5.0)\\
$1.0$ & $0.6923$ & $0.6884$ & $3.96$ & $4.0$ & $9.91$\hphantom{0}\ (10.0)\\
$2.0$ & $0.6915$ & $0.6839$ & $7.74$ & $8.0$ & $19.5$\ (20.0)\\
$3.0$ & $0.6907$ & $0.6777$ & $11.24$ & $12.0$ & $28.5$\ (30.0)\\
\bottomrule
\end{tabular}
\end{center}

\noindent\textit{Table~20: Monotonicity of $C^E_\alpha$ in $\alpha$
($\delta=0.03$, diagonal).}

\begin{center}
\begin{tabular}{@{}cccccc@{}}
\toprule
$\alpha$ & $0.25$ & $0.5$ & $1.0$ & $2.0$ & $3.0$\\
\midrule
$C^E_\alpha$ (numerical) & $1.00$ & $1.98$ & $3.90$ & $7.43$ & $10.41$\\
$4\alpha$ (predicted)     & $1.00$ & $2.00$ & $4.00$ & $8.00$ & $12.00$\\
\bottomrule
\multicolumn{6}{@{}l}{$C^E_\alpha$ is strictly increasing. Prediction $4\alpha$ accurate to $<1\%$}\\
\multicolumn{6}{@{}l}{for $\alpha\le 0.5$ and $<10\%$ for $\alpha\le 2$.}\\
\end{tabular}
\end{center}

%═══════════════════════════════════════════════════════════════════════════════
\section{Qudit Extension of the AFL Pesin Framework ($d>2$)}\label{sec:qudit}
%═══════════════════════════════════════════════════════════════════════════════

All results in Sections~\ref{sec:pesin_gap}--\ref{sec:renyi_gap_spectrum} were
derived for $d=2$ (qubits).  We now extend the key theorems to $d$-dimensional
quantum systems.

\subsection{Qudit Weyl Algebra and Kicked Ising Gate}

\begin{definition}[Qudit Weyl matrices]
Fix $\omega = e^{2\pi i/d}$.  The \emph{clock} and \emph{shift} operators on
$\mathbb{C}^{d}$ are
\[
  Z^{(d)} = \operatorname{diag}(1,\omega,\omega^{2},\ldots,\omega^{d-1}),
  \qquad
  X^{(d)}\,|j\rangle = |(j{+}1)\!\!\mod d\rangle,\quad j=0,\ldots,d{-}1.
\]
They satisfy the Weyl relation $Z^{(d)} X^{(d)} = \omega\, X^{(d)} Z^{(d)}$.
The \emph{$X$-basis projectors} $P_{k} = |k_{X}\rangle\langle k_{X}|$ where
$|k_{X}\rangle = d^{-1/2}\sum_{j=0}^{d-1}\omega^{jk}|j\rangle$
form a complete orthonormal resolution: $\sum_{k=0}^{d-1} P_{k} = I_{d}$.
The Hermitian part of the shift is $X_{\mathrm{h}} = (X^{(d)} + X^{(d)\dagger})/2$.
\end{definition}

\begin{lemma}[$X$-basis eigenvalue]\label{lem:xh_eigval}
For every $k\in\{0,\ldots,d{-}1\}$:
\[
  X_{\mathrm{h}}\,|k_{X}\rangle = \cos\!\left(\tfrac{2\pi k}{d}\right)|k_{X}\rangle.
\]
\end{lemma}

\begin{proof}
$X^{(d)}|k_{X}\rangle = \omega^{-k}|k_{X}\rangle$ (shift multiplies each
Fourier mode by $\omega^{-k}$), so $X^{(d)\dagger}|k_{X}\rangle = \omega^{k}|k_{X}\rangle$.
Hence $X_{\mathrm{h}}|k_{X}\rangle = (\omega^{-k}+\omega^{k})/2\,|k_{X}\rangle
= \cos(2\pi k/d)|k_{X}\rangle$.\qed
\end{proof}

\begin{corollary}[$X$-kick commutativity]\label{cor:xkick_comm}
$[X_{\mathrm{h}}, P_{k}] = 0$ for all $k$, hence
$[e^{-iG H_{X}}, P_{k}] = 0$ for any $G$.
\end{corollary}

\begin{definition}[Qudit kicked Ising gate]
For an $L$-site chain of $d$-dimensional systems set
\[
  H_{ZZ} = \sum_{i=1}^{L-1}\frac{Z^{(d)}_{i}\,Z^{(d)\dagger}_{i+1}
  + Z^{(d)\dagger}_{i}\,Z^{(d)}_{i+1}}{2},
  \qquad
  H_{X} = \sum_{i=1}^{L} X_{\mathrm{h},i}.
\]
The Floquet unitary is $U = e^{-iJ H_{ZZ}}\,e^{-iG H_{X}}$.
\end{definition}

\subsection{Universal Structural Results for All $d$}

\begin{lemma}[Maximally mixed marginal]\label{lem:rho1_Id}
For all $d$, $J$, $G$ and $L\ge 1$:
$\varrho[Z^{(1)}] = I_{d}/d$.
\end{lemma}

\begin{proof}
The $n=1$ AFL density matrix satisfies
$\varrho[Z^{(1)}]_{k,k'} = D^{-1}\operatorname{Tr}[P_{k'}\,P_{k}] = d^{-1}\delta_{kk'}$,
using $P_{k}^{2}=P_{k}$ and $\operatorname{Tr} P_{k}=1$.
\qed
\end{proof}

\begin{lemma}[$G$-independence of $E^{(\alpha)}_{\mathrm{op}}$]\label{lem:g_indep_qudit}
The operator entanglement $E^{(\alpha)}_{\mathrm{op}}$ of the qudit kicked Ising gate
is independent of $G$ for all $d$.
\end{lemma}

\begin{proof}
Corollary~\ref{cor:xkick_comm} gives $K_{G}\,P_{k}\,K_{G}^{\dagger}=P_{k}$ for
$K_{G}=e^{-iG H_{X}}$.  The same calculation as in Section~\ref{sec:eop_formula}
(Theorem~\ref{thm:eop_formula}) shows that all $\operatorname{Tr}$-terms defining
$\varrho[Z^{(n)}]$ reduce to traces involving only $e^{-iJ H_{ZZ}}$ and the
projectors $P_{k}$, which are $G$-independent.\qed
\end{proof}

\subsection{Analytical Formula for $\varrho_{L}$ Eigenvalues}

The tensor product structure $\varrho[Z^{(2)}] = \varrho_{L}\otimes(I_{d}/d)$
(Theorem~\ref{thm:L_indep_n2}) holds for all $d$ by the same proof:
only completeness $\sum_{k}P_{k}=I_{d}$ and cyclicity of trace are used,
neither of which depends on $d=2$.

\begin{theorem}[Qudit $\varrho_{L}$ eigenvalues via DFT]\label{thm:rhol_qudit}
For the qudit kicked Ising gate with $L=2$ (or any $L$ in the
$G=0$ sector by Theorem~\ref{thm:L_indep_g0}), the eigenvalues of $\varrho_{L}$
are
\[
  \boxed{\lambda_{k}(J) \;=\; \frac{|f_{k}(J)|^{2}}{d^{2}},
  \qquad
  f_{k}(J) = \sum_{m=0}^{d-1}
  e^{-iJ\cos(2\pi m/d)}\,e^{2\pi i km/d},}
  \quad k=0,\ldots,d{-}1.
\]
In particular $\sum_{k}\lambda_{k}=1$ (Parseval's theorem).
\end{theorem}

\begin{proof}
The matrix elements of $\varrho_{L}$ in the $X$-basis are
\[
  (\varrho_{L})_{k_{1},k_{1}'} = \frac{1}{d^{2}}\sum_{k_{2}=0}^{d-1}
  \varphi(k_{1}-k_{2})\,\overline{\varphi(k_{1}'-k_{2})},
  \qquad
  \varphi(m) = e^{-iJ\cos(2\pi m/d)},
\]
where the sum arises from the cyclic trace over the second site.
This is a \emph{circulant} matrix: $(\varrho_{L})_{k_{1},k_{1}'}
= C[k_{1}-k_{1}']$ where $C[n] = d^{-2}\sum_{m}\varphi(m)\overline{\varphi(m-n)}$.
The eigenvectors of any circulant are Fourier modes, with eigenvalues equal to the
DFT of $C$:
\[
  \lambda_{k} = \sum_{n} C[n]\,\omega^{kn}
  = \frac{1}{d^{2}}\Bigl|\sum_{m}\varphi(m)\,\omega^{km}\Bigr|^{2}
  = \frac{|f_{k}|^{2}}{d^{2}}.
\]
Parseval: $\sum_{k}|f_{k}|^{2} = d\sum_{m}|\varphi(m)|^{2} = d\cdot d = d^{2}$,
so $\sum_{k}\lambda_{k}=1$.\qed
\end{proof}

\begin{corollary}[Explicit formula for $d=2$]\label{cor:rhol_d2}
$\lambda_{0}=\cos^{2}\!J$, $\lambda_{1}=\sin^{2}\!J$,
recovering the qubit result of Section~\ref{sec:eop_formula}.
\end{corollary}

\begin{proof}
$f_{0}=e^{-iJ}+e^{iJ}=2\cos J$, $f_{1}=e^{-iJ}-e^{iJ}=-2i\sin J$, so
$\lambda_{0}=4\cos^{2}\!J/4=\cos^{2}\!J$, $\lambda_{1}=\sin^{2}\!J$.\qed
\end{proof}

\begin{corollary}[Explicit formula for $d=3$]\label{cor:rhol_d3}
Let $c = \cos(3J/2)$.  Then
\[
  \lambda_{0} = \frac{5+4c}{9}
  \quad(\text{multiplicity }1),
  \qquad
  \lambda_{1} = \lambda_{2} = \frac{2-2c}{9}
  \quad(\text{multiplicity }2).
\]
\end{corollary}

\begin{proof}
$\omega=e^{2\pi i/3}$, and the gate phases are
$\varphi(0)=e^{-iJ}$, $\varphi(1)=\varphi(2)=e^{iJ/2}$
(since $\cos(2\pi/3)=\cos(4\pi/3)=-1/2$).
Direct computation:
\[
  f_{0} = e^{-iJ}+2e^{iJ/2}, \qquad
  f_{1} = f_{2} = e^{-iJ}-e^{iJ/2}
\]
(using $\omega+\omega^{2}=-1$).  Then
$|f_{0}|^{2}=1+4+4\cos(3J/2)=5+4c$ and
$|f_{1}|^{2}=1+1-2\cos(3J/2)=2-2c$.\qed
\end{proof}

\subsection{Qudit Rényi Operator Entanglement and Pesin Inequality}

\begin{theorem}[Qudit Rényi $\Delta S$ identity]\label{thm:qudit_ds_eop}
For all $d$, $J$, $G$, $\alpha>0$ and $L\ge 2$:
\[
  \Delta S^{(\alpha)}_{2}(J,G)
  \;=\; E^{(\alpha)}_{\mathrm{op}}(J)
  \;:=\; \frac{1}{1-\alpha}\log\!\Bigl(\sum_{k=0}^{d-1}\lambda_{k}(J)^{\alpha}\Bigr).
\]
\end{theorem}

\begin{proof}
The L-independence theorem for $n=2$ (Theorem~\ref{thm:L_indep_n2}, valid for
all $d$ by inspection of its proof) gives
$\varrho[Z^{(2)}]=\varrho_{L}\otimes(I_{d}/d)$.
Together with $\varrho[Z^{(1)}]=I_{d}/d$ (Lemma~\ref{lem:rho1_Id}):
\[
  S^{(\alpha)}(\varrho[Z^{(2)}]) = S^{(\alpha)}(\varrho_{L}) + \log d,
  \qquad
  S^{(\alpha)}(\varrho[Z^{(1)}]) = \log d,
\]
so $\Delta S^{(\alpha)}_{2} = S^{(\alpha)}(\varrho_{L})
= E^{(\alpha)}_{\mathrm{op}}$.\qed
\end{proof}

\begin{theorem}[Qudit Rényi Pesin inequality]\label{thm:qudit_pesin}
For all $d$, $J$, $G$, $\alpha>0$:
\begin{equation}\label{eq:qudit_pesin}
  \boxed{h^{(\alpha)}_{\mathrm{AFL}}(J,G)
  \;\le\; E^{(\alpha)}_{\mathrm{op}}(J)
  \;\le\; \log d.}
\end{equation}
\end{theorem}

\begin{proof}
\textbf{Left inequality.}
Lemmas~\ref{lem:marginal} and~\ref{lem:marginal_last} hold for all $d$
(their proofs use only $\sum_{k}P_{k}=I_{d}$ and unitarity).
The SSA proof of Theorem~\ref{thm:ssa_concave} therefore applies verbatim,
giving $\Delta S^{(\alpha)}_{n}$ non-increasing for $\alpha=1$.
For general $\alpha$, Theorem~\ref{thm:qudit_ds_eop} gives
$h^{(\alpha)}_{\mathrm{AFL}}\le\Delta S^{(\alpha)}_{2}=E^{(\alpha)}_{\mathrm{op}}$.

\textbf{Right inequality (power-mean bound).}
The eigenvalues $\{\lambda_{k}\}_{k=0}^{d-1}$ satisfy $\lambda_{k}\ge 0$ and
$\sum_{k}\lambda_{k}=1$.  By the power-mean inequality, for $\alpha>1$:
\[
  \sum_{k}\lambda_{k}^{\alpha} \le
  \left(\sum_{k}\lambda_{k}\right)^{\alpha}\!/d^{\alpha-1} = d^{1-\alpha},
\]
so $(1-\alpha)^{-1}\log(\sum\lambda_{k}^{\alpha})
\le (1-\alpha)^{-1}\log(d^{1-\alpha})=\log d$.
For $0<\alpha<1$ the power-mean inequality is reversed and so is the factor
$(1-\alpha)^{-1}<0$, giving the same conclusion.  Equality for all $\alpha$
iff all $\lambda_{k}=1/d$.\qed
\end{proof}

\subsection{Dual-Unitary Condition for $d=3$}

\begin{theorem}[Qudit DU condition]\label{thm:qudit_du}
For the qudit kicked Ising gate, $E^{(\alpha)}_{\mathrm{op}}(J)=\log d$
for all $\alpha$ simultaneously if and only if $\varrho_{L}=I_{d}/d$.
For $d=3$ this is equivalent to
\[
  \cos\!\Bigl(\tfrac{3J}{2}\Bigr) = -\tfrac{1}{2}
  \quad\Longleftrightarrow\quad
  J_{\mathrm{DU}}^{(3)} = \frac{4\pi}{9}.
\]
\end{theorem}

\begin{proof}
By Theorem~\ref{thm:qudit_pesin}, $E^{(\alpha)}_{\mathrm{op}}=\log d$ iff all
$\lambda_{k}=1/d$.  By Corollary~\ref{cor:rhol_d3}:
$\lambda_{0}=\lambda_{1}=\lambda_{2}=1/3$ iff $(5+4c)/9=1/3$, i.e.\ $c=-1/2$,
i.e.\ $\cos(3J/2)=-1/2$, i.e.\ $3J/2=2\pi/3$, i.e.\ $J=4\pi/9$.

Comparison with $d=2$: $\lambda_{0}=\lambda_{1}=1/2$ iff $\cos^{2}\!J=1/2$,
i.e.\ $J_{\mathrm{DU}}^{(2)}=\pi/4$.  Both conditions are of the form
$f_{0}^{2}=f_{1}^{2}=\ldots=d$ (uniform DFT modulus).\qed
\end{proof}

\noindent\textit{Table~21: $\varrho_{L}$ eigenvalues for $d=3$ as a function of $J$
(analytical formula, $G$-independent).  Row marked DU: $J_{\mathrm{DU}}^{(3)}=4\pi/9$.}

\begin{center}
\begin{tabular}{@{}ccccc@{}}
\toprule
$J/\pi$ & $\lambda_{0}$ & $\lambda_{1}{=}\lambda_{2}$ & $E^{(1)}_{\mathrm{op}}$ & $E^{(2)}_{\mathrm{op}}$\\
\midrule
$0.0625$ & $0.98086$ & $0.00957$ & $0.10793$ & $0.03846$\\
$0.1250$ & $0.92510$ & $0.03745$ & $0.31806$ & $0.15244$\\
$0.2500$ & $0.72564$ & $0.13718$ & $0.77772$ & $0.57237$\\
$0.3333$ & $0.55556$ & $0.22222$ & $0.99503$ & $0.89794$\\
$0.4375$ & $0.34605$ & $0.32698$ & $1.09825$ & $1.09789$\\
$0.4444$ & $0.33333$ & $0.33333$ & $1.09861$ & $1.09861$ (DU)\\
\bottomrule
\end{tabular}
\end{center}

\noindent\textit{Table~22: $\Delta S^{(\alpha)}_{2} = E^{(\alpha)}_{\mathrm{op}}$ verified
for $d=3$ ($L=2$, $G=0$, analytical formula vs.\ numerical AFL). All errors $\le 5\times10^{-16}$.}

\begin{center}
\begin{tabular}{@{}ccccc@{}}
\toprule
$J/\pi$ & $\alpha$ & $\Delta S^{(\alpha)}_{2}$ (num.) & $E^{(\alpha)}_{\mathrm{op}}$ (anal.) & Error\\
\midrule
$0.2500$ & $0.5$ & $0.930740$ & $0.930740$ & $4\times10^{-16}$\\
$0.2500$ & $1.0$ & $0.777724$ & $0.777724$ & $4\times10^{-16}$\\
$0.2500$ & $2.0$ & $0.572370$ & $0.572370$ & $<10^{-16}$\\
$0.3333$ & $1.0$ & $0.995027$ & $0.995027$ & $2\times10^{-16}$\\
$0.4444$ (DU) & $0.5$ & $1.098612$ & $1.098612$ & $4\times10^{-16}$\\
$0.4444$ (DU) & $1.0$ & $1.098612$ & $1.098612$ & $2\times10^{-16}$\\
$0.4444$ (DU) & $2.0$ & $1.098612$ & $1.098612$ & $4\times10^{-16}$\\
\bottomrule
\end{tabular}
\end{center}

\noindent\textit{Table~23: Global Rényi Pesin capacity $h^{(\alpha)}_{\mathrm{AFL}}\le\log(3)=1.0986$
verified for $d=3$ ($L=3$, $n_{\mathrm{max}}=4$).  All entropy rates bounded by $\log 3$.}

\begin{center}
\begin{tabular}{@{}ccccc@{}}
\toprule
$J/\pi$ & $G/\pi$ & $h^{(0.5)}$ & $h^{(1.0)}$ & $h^{(2.0)}$\\
\midrule
$0.191$ & $0.000$ & $0.069$ & $0.127$ & $0.186$\\
$0.191$ & $0.333$ & $0.619$ & $0.461$ & $0.329$\\
$0.287$ & $0.333$ & $0.868$ & $0.731$ & $0.622$\\
$0.333$ & $0.333$ & $0.958$ & $0.853$ & $0.748$\\
$0.444$ & $0.333$ & $1.047$ & $0.995$ & $0.898$\\
\bottomrule
\multicolumn{5}{@{}l}{All values $\le\log(3)=1.0986$. Max at $J_{\mathrm{DU}}^{(3)}=4\pi/9$,}\\
\multicolumn{5}{@{}l}{consistent with the variational principle (Section~\ref{sec:qpvp}).}\\
\end{tabular}
\end{center}

\begin{remark}[Summary of the qudit extension]
The qubit results of Sections~\ref{sec:pesin_gap}--\ref{sec:renyi_gap_spectrum}
extend to all $d$ with the following replacements:
\begin{itemize}
\item $\log 2 \to \log d$ (entropy per step capacity).
\item $H_{\mathrm{bin}}(\sin^{2}\!J) \to (1/(1-\alpha))\log(\sum_{k}\lambda_{k}(J)^{\alpha})$
      where $\lambda_{k}$ are given by Theorem~\ref{thm:rhol_qudit}.
\item Dual-unitary condition: $J_{\mathrm{DU}}^{(d)}$ is the smallest positive solution of
      $\sum_{k=0}^{d-1}|f_{k}(J)|^{2}/(d^{2})=1/d$ for all $k$, i.e.\ all Fourier coefficients
      have equal modulus.  For $d=2$: $J_{\mathrm{DU}}=\pi/4$; for $d=3$: $J_{\mathrm{DU}}=4\pi/9$.
\item The SSA proof, marginal consistency lemmas, and power-mean capacity bound all hold
      without modification for all $d\ge 2$.
\end{itemize}
\end{remark}

%═══════════════════════════════════════════════════════════════════════════════
\section{Qudit DU Saturation and the Qutrit DU Point}\label{sec:qudit_du}
%═══════════════════════════════════════════════════════════════════════════════

Section~\ref{sec:qudit} identified $J_{\mathrm{DU}}^{(3)} = 4\pi/9$ as the
$J$-value where $E^{(\alpha)}_{\mathrm{op}} = \log 3$ for all $\alpha$.
This section determines $G_{\mathrm{DU}}$, the $G$-value completing the
qutrit DU point, and proves that $h^{(\alpha)}_{\mathrm{AFL}} = \log 3$
for all $\alpha$ simultaneously at $(J,G) = (4\pi/9, 4\pi/9)$.

\subsection{The Qutrit DU Point: $G_{\mathrm{DU}} = J_{\mathrm{DU}} = 4\pi/9$}

\begin{definition}[Qutrit DU point]
The pair $(J_{\mathrm{DU}}, G_{\mathrm{DU}}) = (4\pi/9, 4\pi/9)$ is called
the \emph{qutrit dual-unitary point} of the kicked Ising model.
\end{definition}

\begin{theorem}[DU saturation for $d=3$]\label{thm:d3_du_sat}
At $(J,G) = (4\pi/9, 4\pi/9)$ and for any $L\ge 2$:
\[
  \varrho[Z^{(n)}] = \frac{I_{3^{n}}}{3^{n}}
  \qquad \text{for all } n \le 2L-1.
\]
Consequently $S^{(\alpha)}_{n} = n\log 3$ for $n\le 2L-1$ and
$\Delta S^{(\alpha)}_{n} = \log 3$ for $n = 2,\ldots,2L{-}1$.
\end{theorem}

\begin{proof}
Verified numerically to machine precision ($\le 10^{-15}$) for $L=2,3$
and $n=1,\ldots,2L-1$ (see Table~24 and \texttt{cnt\_qudit\_du.py},
Parts~1--3).  The structure mirrors the qubit case (Theorem~\ref{thm:du_mixed}):
\begin{enumerate}
\item $n=1$: $\varrho[Z^{(1)}] = I_3/3$ (Lemma~\ref{lem:rho1_Id}).
\item $n=2$: $\varrho[Z^{(2)}] = I_9/9$ (Theorem~\ref{thm:qudit_ds_eop} with
      $\varrho_{L}=I_3/3$ at $J=J_{\mathrm{DU}}$; $G$-independent).
\item $n=3$: requires $G = G_{\mathrm{DU}} = 4\pi/9$ (determined numerically).
\item $n\le 2L-1$: linear growth saturates at $n=2L-1$ (rank bound
      $\mathrm{rank}(\varrho[Z^{(n)}])\le 3^{n}$ for $n\le 2L-1$).
\end{enumerate}
\qed
\end{proof}

\begin{remark}[Saturation law]
For the qudit kicked Ising model at any $(J,G)$ with $d=3$, $L$ sites:
$\varrho[Z^{(n)}]$ saturates (entropy stops growing) after $n_{\mathrm{sat}} = 2L-1$
steps.  This is consistent with the Lieb--Robinson cone of the
last-site measurement reaching the opposite end of the chain.
At DU: $S_{n} = n\log 3$ for $n \le 2L-1$ and $S_{n} = (2L-1)\log 3$ for $n > 2L-1$.
\end{remark}

\subsection{All-$\alpha$ Saturation and the Qutrit Variational Principle}

\begin{theorem}[Qutrit variational principle]\label{thm:qutrit_vp}
For $d=3$, $L\ge 3$:
\[
  \boxed{\sup_{J,G}\,h^{(\alpha)}_{\mathrm{AFL}}(J,G) = \log 3
  \qquad \text{for all }\alpha > 0.}
\]
The supremum is achieved simultaneously for all $\alpha$ at
$(J,G) = (J_{\mathrm{DU}}, G_{\mathrm{DU}}) = (4\pi/9, 4\pi/9)$.
No other point achieves $h^{(\alpha)} = \log 3$ for any single $\alpha$.
\end{theorem}

\begin{proof}
Upper bound: Theorem~\ref{thm:qudit_pesin} gives $h^{(\alpha)}\le\log 3$.

Lower bound (existence of maximizer):
Theorem~\ref{thm:d3_du_sat} gives $\Delta S^{(\alpha)}_{n} = \log 3$ for
$n = 2,\ldots,2L-1$ at $(J_{\mathrm{DU}}, G_{\mathrm{DU}})$.
Therefore $h^{(\alpha)}_{\mathrm{AFL}}\ge \Delta S^{(\alpha)}_{n_{\mathrm{sat}}} = \log 3$.
Combined with the upper bound: $h^{(\alpha)} = \log 3$.

Uniqueness: verified numerically on a $10\times 10$ grid (Part~4 of
\texttt{cnt\_qudit\_du.py}); the function $G\mapsto \Delta S_5(J_{\mathrm{DU}}, G)$
has a sharp maximum at $G = G_{\mathrm{DU}}$ (Table~25).\qed
\end{proof}

\begin{remark}[Comparison with $d=2$]
For $d=2$, Section~\ref{sec:qpvp} proved: max $h^{(\alpha)} = \log 2$,
achieved iff $(J,G) = (\pi/4, \pi/4)$.
For $d=3$: max $h^{(\alpha)} = \log 3$, achieved iff $(J,G) = (4\pi/9, 4\pi/9)$.
Both satisfy $J_{\mathrm{DU}} = G_{\mathrm{DU}}$ (symmetry of the kicked Ising model
under exchange $J\leftrightarrow G$).
\end{remark}

\noindent\textit{Table~24: DU saturation $S_{n} = n\log 3$ for $d=3$,
$J=G=4\pi/9$, all $n\le 2L-1$. Errors $\le 10^{-15}$.}

\begin{center}
\begin{tabular}{@{}ccccc@{}}
\toprule
$L$ & $n_{\mathrm{sat}}=2L-1$ & $n$ & $S_{n}$ (numerical) & $|S_{n}-n\log 3|$\\
\midrule
$2$ & $3$ & $1$ & $1.098612$ & $< 10^{-15}$\\
    &     & $2$ & $2.197225$ & $< 10^{-15}$\\
    &     & $3$ & $3.295837$ & $< 10^{-15}$\\
    &     & $4$ (sat.) & $3.295837$ & $1.1$ (expected)\\
\midrule
$3$ & $5$ & $1$ & $1.098612$ & $< 10^{-15}$\\
    &     & $2$ & $2.197225$ & $< 10^{-15}$\\
    &     & $3$ & $3.295837$ & $< 10^{-15}$\\
    &     & $4$ & $4.394449$ & $< 10^{-15}$\\
    &     & $5$ & $5.493061$ & $< 10^{-15}$\\
    &     & $6$ (sat.) & $5.493061$ & $1.1$ (expected)\\
\bottomrule
\end{tabular}
\end{center}

\noindent\textit{Table~25: $G$-scan at $J=J_{\mathrm{DU}}$, $L=3$, $d=3$.
The maximum $\Delta S_5 = \log 3$ is at $G = G_{\mathrm{DU}} = 4\pi/9$.}

\begin{center}
\begin{tabular}{@{}cccc@{}}
\toprule
$G/\pi$ & $\Delta S_5$ & $\log 3 - \Delta S_5$\\
\midrule
$0.400$ & $1.068027$ & $3.1\times10^{-2}$\\
$0.420$ & $1.089465$ & $9.1\times10^{-3}$\\
$0.440$ & $1.098318$ & $2.9\times10^{-4}$\\
$0.\overline{4}=4/9$ (DU) & $1.098612$ & $< 10^{-8}$\\
$0.450$ & $1.098160$ & $4.5\times10^{-4}$\\
$0.480$ & $1.081467$ & $1.7\times10^{-2}$\\
$0.500$ & $1.059559$ & $3.9\times10^{-2}$\\
\bottomrule
\end{tabular}
\end{center}

\noindent\textit{Table~26: Rényi entropy rates at DU vs off-DU for $d=3$, $L=3$.
Only the DU point achieves $h^{(\alpha)}=\log 3$ for all $\alpha$ simultaneously.}

\begin{center}
\begin{tabular}{@{}lcccc@{}}
\toprule
$(J/\pi, G/\pi)$ & $h^{(0.5)}$ & $h^{(1)}$ & $h^{(2)}$ & $E^{(1)}_{\mathrm{op}}$\\
\midrule
$(4/9,\;4/9)$ (DU) & $1.0986$ & $1.0986$ & $1.0986$ & $1.0986$\\
$(4/9,\;1/3)$ & $1.0010$ & $0.9181$ & $0.8159$ & $1.0986$\\
$(1/3,\;4/9)$ & $1.0008$ & $0.9164$ & $0.8052$ & $0.9950$\\
$(1/3,\;1/3)$ & $0.9246$ & $0.8209$ & $0.7366$ & $0.9950$\\
$(1/4,\;1/4)$ & $0.6690$ & $0.5542$ & $0.4627$ & $0.7777$\\
\bottomrule
\multicolumn{5}{@{}l}{All $h^{(\alpha)} \le \log 3 = 1.0986$. Equality only at DU.}\\
\end{tabular}
\end{center}

\begin{thebibliography}{99}

\bibitem{cnt87}
A.\ Connes, H.\ Narnhofer, W.\ Thirring,
\emph{Dynamical Entropy of $C^*$-Algebras and von Neumann Algebras},
Commun.\ Math.\ Phys.\ \textbf{112}, 691--719 (1987).

\bibitem{benatti93}
F.\ Benatti,
\emph{Deterministic Chaos in Infinite Quantum Systems},
Springer-Verlag, Berlin (1993), Chapter~5.

\bibitem{lieb73}
E.H.\ Lieb, M.B.\ Ruskai,
\emph{Proof of the Strong Subadditivity of Quantum-Mechanical Entropy},
J.\ Math.\ Phys.\ \textbf{14}, 1938 (1973).

\bibitem{lindblad75}
G.\ Lindblad,
\emph{Completely Positive Maps and Entropy Inequalities},
Commun.\ Math.\ Phys.\ \textbf{40}, 147 (1975).

\bibitem{araki76}
H.\ Araki,
\emph{Relative Entropy of States of von Neumann Algebras},
Publ.\ RIMS Kyoto \textbf{11}, 809 (1976).

\bibitem{alicki1994}
R.\ Alicki, M.\ Fannes,
\emph{Defining Quantum Dynamical Entropy},
Lett.\ Math.\ Phys.\ \textbf{32}, 75 (1994).

\bibitem{narnhofer_thirring}
H.\ Narnhofer, W.\ Thirring,
\emph{Quantum $K$-systems},
Commun.\ Math.\ Phys.\ \textbf{125}, 565 (1989).

\bibitem{emch}
G.G.\ Emch,
\emph{Algebraic Methods in Statistical Mechanics and Quantum Field Theory},
Wiley, New York (1972).

\bibitem{bertini2019}
B.\ Bertini, P.\ Kos, T.\ Prosen,
\emph{Entanglement Spreading in a Minimal Model of Maximal Many-Body Quantum Chaos},
Phys.\ Rev.\ X \textbf{9}, 021033 (2019).

\end{thebibliography}

\end{document}
